\subsection{The algebra actually written on 155/002}

Throughout this subsection, $k$ is a field and $L/k$ is finite. All tensor
categories and tensor functors are symmetric and unital. We make this convention
explicit: without symmetry the commutativity of $U$ in the manuscript is not
the appropriate conclusion. A pro-vector space means the inverse system dual
to an ordinary vector space; completed tensor products have this meaning.
All ideals in such a space are closed. Finite-dimensional examples need no
completion.

\begin{saresult}{B.001}{Proved}{The idealizer is the Takeuchi product}
Let $A$ be a unital $k$-algebra containing $L$ by $s:L\longrightarrow A$.
Put $D=A\otimes_k A$ and let $J$ be the right ideal generated by
$s(\lambda)\otimes1-1\otimes s(\lambda)$, for $\lambda\in L$.
The quotient $D/J$ is the tensor product of the two left $L$-modules $A$.
The image of the left idealizer $\widetilde J=\{d:dJ\subset J\}$ is
\[
 A\mathbin{\times_L}A=
 \left\{\sum_i a_i\otimes_L b_i:
 \sum_i a_i s(\lambda)\otimes_L b_i
 =\sum_i a_i\otimes_L b_i s(\lambda)
 \text{ for every }\lambda\in L\right\}.
\]
Multiplication is induced by factorwise multiplication in $D$. Thus
$\widetilde J/J\simeq A\times_L A$ as algebras. The same statement holds
for the completed products of 155/002.
\end{saresult}
\begin{proof}
The defining relations for $D/J$ are
$s(\lambda)a\otimes b=a\otimes s(\lambda)b$, exactly the balancing
relations of the indicated tensor product. An element $d$ idealizes $J$
if and only if $d(s(\lambda)\otimes1-1\otimes s(\lambda))$ lies in $J$
for every $\lambda$, because these differences generate $J$ on the right.
This is the displayed condition. The idealizer is a subalgebra, and $J$
is a two-sided ideal in it: right multiplication preserves $J$ by its
definition, and left multiplication by an idealizer element preserves it
by the idealizer condition. This proves the algebra assertion, without
incorrectly assigning an algebra structure to all of $A\otimes_L A$.
The completed assertion follows by the same closed-relations argument,
or by taking inverse limits of the finite linear quotients.
\end{proof}

\begin{saresult}{B.002}{Refuted}{The displayed conditions on 155/002 omit a unit constraint}
Even for $L=k=\mathbb F_2$ and finite-dimensional $U$, the data 1), 2),
(i), and (ii) on 155/002 do not imply the asserted tensor-category
reconstruction. The missing constraint is the character identity for
the coalgebra counit. In the neutral case it says that this counit is
an algebra homomorphism. Consequently R7 and L6 require this correction.
\end{saresult}
\begin{proof}
Let $S=\{a,b,z\}$, with $a^2=a$, $b^2=b$, $ab=ba=z$, and $zs=sz=z$
for every $s\in S$. This is a commutative associative semigroup with no
identity. Let $U=kS$ be its semigroup algebra. It is nevertheless a
unital algebra: its identity is $e_a+e_b-e_z$. The map
$\varepsilon_U:U\longrightarrow k$, $e_s\longmapsto1$, is a unital
algebra homomorphism. Take $A=U^\vee=k^S$, with pointwise multiplication
and identity the constant function $1$, which is precisely
$\varepsilon_U$. Condition (i) is the ordinary inclusion of scalar
functions. The transpose of the multiplication of $U$ is
\[
 \Delta:A\longrightarrow A\otimes_k A=k^{S\times S},
 \qquad \Delta(f)(s,t)=f(st).
\]
It is a unital algebra homomorphism. Since $J=0$, this verifies (ii)
including its strongest usual unital interpretation. Coassociativity
and cocommutativity follow from the semigroup law. The coalgebra counit
dual to $1_U$ is
\[
 e:A\longrightarrow k,\qquad e(f)=f(a)+f(b)-f(z).
\]
For the indicator functions $p_a,p_b$, we have
$p_ap_b=0$ and $e(p_a)=e(p_b)=1$; hence $e$ is not multiplicative.
All these calculations remain valid in characteristic two.

There is also an obstruction independent of terminology. Finite
$A$-modules have three one-dimensional simples $V_s$, and the proposed
tensor product satisfies $V_s\otimes V_t=V_{st}$. A unit whose image
under the proposed fibre functor is $k$ must be some $V_s$. It would
require $st=t$ for all $t$, but $S$ has no such $s$. Thus the proposed
category has no permitted tensor unit. The attached program
\texttt{check\_B\_finite.py} exhausts all eight elements of $U$ over
$\mathbb F_2$, all $512$ associativity triples, all $64$ augmentation
pairs, and all $64$ coproduct-multiplicativity pairs. It also checks both
counit identities and all three possible one-dimensional units.
\end{proof}

\begin{saresult}{B.003}{Proved under stated hypotheses}{Corrected pro-bialgebroid reconstruction}
In the manuscript data, assume that $A$ is a pseudocompact $k$-algebra
and that the multiplication is continuous. Write $e:A\longrightarrow L$
for the transpose of $1_U$. Add
\begin{equation}\label{eq:B-character}
 e(1)=1,\qquad
 e(ab)=e\bigl(a\,s(e(b))\bigr)\quad(a,b\in A).
\end{equation}
The first equality already follows if the stated augmentation of $U$
is unital. With this addition, the displayed data are exactly the data
of a continuous cocommutative left $L$-bialgebroid with equal source
and target $s$. They give a symmetric abelian tensor category of
continuous $A$-modules finite-dimensional over $L$, with exact faithful
forgetful functor to $\operatorname{Modf}(L)$. Conversely every
essentially small abelian $k$-linear symmetric tensor category with
such a fibre functor reconstructs these data. Both reconstructions are
inverse up to an equivalence respecting the fibre functors and the
induced isomorphism of the pro-algebras.
\end{saresult}
\begin{proof}
For the forward direction, the tensor product of two modules has
underlying space $M\otimes_L N$ and action
\[
 a(m\otimes n)=\sum a_{(1)}m\otimes a_{(2)}n.
\]
The sum is interpreted in a finite quotient when acting on $M,N$.
B.001 proves that balancing over $L$ is respected. Multiplicativity
of $\Delta$ proves the module law; coassociativity gives associativity
of tensor products; cocommutativity gives the usual symmetry. The
tensor unit is $L$, with action
\[
 a\cdot\lambda=e(a s(\lambda)).
\]
Equation \eqref{eq:B-character} says exactly that this is an $A$-action:
$(ab)\cdot\lambda=a\cdot(b\cdot\lambda)$. The ordinary coalgebra
counit identities, applied to $\Delta(a s(\lambda))$, show that
$L\otimes_L M\longrightarrow M$ and $M\otimes_L L\longrightarrow M$
are $A$-linear. The action of $s(\lambda)$ is ordinary scalar
multiplication, since $e$ is left $L$-linear. All coherence maps are
the usual linear ones. Kernels and cokernels retain continuous
$A$-actions and their finite dimension; the forgetful functor is exact
and faithful. Tensoring is exact because $L$ is a field.

For reconstruction, put $W=\operatorname{Res}_{L/k}F$ and give
$A=\operatorname{Nat}_k(W,W)$ its finite topology. The category is
locally finite over $k$: its Hom spaces inject into finite-dimensional
linear Hom spaces, and every strictly increasing chain of subobjects
strictly increases $\dim_LF(X)$. The coalgebra reconstruction theorem
in Deligne \cite[6.1--6.2]{Deligne1990}, applied to $W$, therefore
identifies the category with the finite-dimensional continuous
$A$-modules. This is an application of the cited theorem, with all its
finiteness hypotheses checked, rather than an assertion that an
arbitrary faithful functor suffices. Scalar multiplication on $F(X)$
gives $s:L\longrightarrow A$.

For completeness, the algebraic passage from the tensor product to
$\Delta$ can be checked directly after this equivalence. A natural
endomorphism of the bifunctor $(M,N)\longmapsto M\otimes_L N$ is
determined on pairs of finite regular quotients $A/I,A/J$ by its value
at $1\otimes1$. Naturality under the maps $A/I\longrightarrow M$,
$a\longmapsto am$, and their analogues for $N$, determines all its
values. Naturality under right multiplication by $s(\lambda)$ is
exactly the Takeuchi condition in B.001. Compatible values on the
regular quotients are therefore precisely the elements of the
completed Takeuchi product. The action of each $a\in A$ on the
tensor-product object yields $\Delta(a)$. Associativity, symmetry,
and the action law give the coproduct axioms. The action on the tensor
unit gives $e(a)=a_1(1)$ and hence \eqref{eq:B-character}. Dualizing
the coproduct gives the commutative unital $L$-algebra $U$; dualizing
$1_A$ gives its augmentation. Finally the same regular-quotient
argument identifies the natural endomorphism algebra of the
forgetful functor on the reconstructed category with $A$ itself.
This proves both inverse comparisons. The categorical equivalence is
also the finite-field-base case of Hai
\cite[Theorem 2.1.2 and Corollary 2.2.5]{Hai2008}, with the manuscript's
pro-dual convention in place of his coring convention.
\end{proof}

\begin{saresult}{B.004}{Known}{The precise answer to 155/004, including inseparability}
Assume that the category in B.003 is rigid and that
$k\simeq\operatorname{End}(1)$. Its scalar extension to $L$ has an exact
faithful fibre functor induced by $F$. In the module description this
is the restricted diagonal functor on finite
$(A\widehat\otimes_k L)$-modules written on 155/004. It is exact and
faithful for every finite field extension $L/k$, including inseparable
ones. By contrast, the diagonal functor on all
$(L\otimes_k L)$-modules need not be exact, and is not generally
faithful even for a nontrivial separable extension.
\end{saresult}
\begin{proof}
Deligne \cite[Theorem 1.12, proved in 6.8 and 6.14--6.15]{Deligne1990}
identifies the rigid category with representations of the affine
groupoid $G\rightrightarrows\operatorname{Spec}L$, whose map to
$\operatorname{Spec}(L\otimes_k L)$ is faithfully flat. After extension
of the ground field to $L$, the object scheme has its diagonal
$L$-point. Restriction of an equivariant vector bundle to this point
is exact because the underlying modules in the representation
category are locally free. It is faithful because the groupoid is
transitive: after the faithfully flat base change parametrizing
arrows from that point to all objects, every equivariant map is
transported from its value at the point. A map vanishing there thus
vanishes after a faithfully flat base change and hence vanishes.
The groupoid and scalar-extension descriptions agree by Deligne's
scalar-extension construction in \S5 and its tensor-product
compatibility in 5.18--5.21. On modules the restriction is precisely
$M'\longmapsto M'\otimes_{L\otimes_k L}L$.

The unrestricted assertions are elementary. If
$L=k(u)$, $u^p=t\in k$ is purely inseparable of degree $p$, then
$L\otimes_k L\simeq L[\delta]/(\delta^p)$ and the diagonal quotient
is $\delta\mapsto0$. Tensoring the injection
$(\delta^{p-1})\hookrightarrow L[\delta]/(\delta^p)$ with $L$ gives
a nonzero source and a zero map. Thus it is not exact. For a nontrivial
finite separable extension, the diagonal component splits off as an
idempotent summand of $L\otimes_k L$; the other nonzero summand is
killed, so faithfulness fails. Neither test is a counterexample to
the restricted functor: arbitrary modules over the scalar subring
need not carry the required $A\otimes_k L$-action.
\end{proof}

\begin{saresult}{B.005}{Proved}{What completes the duality calculation on 155/006--010}
Choose a nonzero $k$-linear functional $\ell:L\longrightarrow k$ and
write $U=\operatorname{Hom}^{\mathrm{cont}}_L(A,L)$, where $A$ is a left
$L$-module by $s$. Give $U$ the left $A$-action
$(a u)(b)=u(ba)$. Then
\[
 \operatorname{Hom}_A(M,U)\simeq\operatorname{Hom}_L(M,L),
 \qquad \phi\longmapsto\bigl(m\longmapsto\phi(m)(1)\bigr),
 \qquad \operatorname{End}_A(U)\simeq A^{\mathrm{op}}.
\]
The last isomorphism uses $\ell$ when $L\ne k$. A rigid tensor
structure supplies a continuous anti-isomorphism
$a_\ell:A\longrightarrow A^{\mathrm{op}}$, fixing $s(L)$, and hence a
second, commuting $A$-action on $U$. Evaluation and coevaluation are
\[
 \operatorname{ev}_M(m\otimes\phi)=\phi(m)(1),\qquad
 \operatorname{coev}_M(1)=\sum_i m_i\otimes\phi_i,
\]
where $(m_i), (\phi_i)$ are dual $L$-bases under the first isomorphism.
Conversely a continuous anti-isomorphism fixing $s(L)$ gives rigidity
exactly when these two maps are $A$-linear for every finite continuous
$A$-module $M$. Their two triangle identities always hold as linear
maps, and therefore hold in the module category when the maps are
$A$-linear. For pro-$A$, the manuscript's instruction to test only
$M=A$ must be replaced by tests on finite modules; $A$ need not itself
be an object of the category.
\end{saresult}
\begin{proof}
Given $f\in\operatorname{Hom}_L(M,L)$, set
$\phi_f(m)(b)=f(bm)$. This is left $L$-linear in $b$, continuous, and
$A$-linear in $m$. It inverts the displayed evaluation-at-$1$ map.
The bilinear form $(x,y)\longmapsto\ell(xy)$ on the field $L$ is
nondegenerate: for $x\ne0$, choose $z$ with $\ell(z)\ne0$ and put
$y=x^{-1}z$. Consequently
\[
 U\longrightarrow\operatorname{Hom}^{\mathrm{cont}}_k(A,k),
 \qquad u\longmapsto\ell\circ u
\]
is an isomorphism. Its inverse sends $g$ to the unique $u$ satisfying
$\ell(\lambda u(b))=g(s(\lambda)b)$ for all $\lambda\in L$.
The map intertwines the coinduced $A$-actions. Continuous duality
therefore gives
\[
 \operatorname{End}_A(U)
 \simeq\operatorname{End}_{A^{\mathrm{op}}}^{\mathrm{cont}}(A)^{\mathrm{op}}
 \simeq A^{\mathrm{op}},
\]
because a continuous endomorphism of the right regular module is
left multiplication by its value at $1$. Under this identification,
scalar multiplication on the values of $u$ corresponds to $s(L)$.

Suppose now that $D$ is the rigid duality. The fibre functor identifies
$F(DM)$ with the $L$-dual of $F(M)$. Pair it with $F(M)$ by
$B_M(\phi,m)=\ell(\phi(m))$. For $a\in A$, define
$a_\ell(a)_M$ as the adjoint of $a_{DM}$ under this nondegenerate
$k$-pairing. Naturality of the pairing makes this a natural
endomorphism of $F$ as a $k$-valued functor. Transposition reverses
products, and the bidual isomorphism gives the inverse map.
Multiplication by $\lambda\in L$ is self-adjoint, so $a_\ell$ fixes
$s(L)$. This is the required anti-isomorphism. The above adjunction
then identifies the dual module with $\operatorname{Hom}_A(M,U)$.
Its evaluation and coevaluation are the images of the corresponding
morphisms of the given rigid category, and hence are $A$-linear.

Conversely the stated anti-isomorphism defines a contravariant
equivalence on finite modules through $\operatorname{Hom}_A(-,U)$;
on underlying spaces it is exact ordinary $L$-duality. If evaluation
and coevaluation are $A$-linear, they give duality in the tensor
category. Explicitly the first triangle sends
$m$ to $\sum_i m_i\phi_i(m)=m$, and the second sends
$\phi$ to $\sum_i\phi(m_i)\phi_i=\phi$. These computations also prove
basis independence. A fully explicit criterion for $A$-linearity is
\[
 \sum\operatorname{ev}_M(a_{(1)}m\otimes a_{(2)}\phi)
 =a\cdot\operatorname{ev}_M(m\otimes\phi),
\]
\[
 \sum_i\sum a_{(1)}m_i\otimes a_{(2)}\phi_i
 =e(a)\sum_i m_i\otimes\phi_i.
\]
Here $a\cdot\lambda=e(a s(\lambda))$, so both equations are typed
over $L$. These are the evaluation and coevaluation constraints
missing from a bare anti-homomorphism. A finite separable extension
has the canonical permissible choice $\ell=\operatorname{Tr}_{L/k}$;
for a purely inseparable extension the trace is zero and cannot be
used for this purpose. We assert existence of a choice, not a
canonical inseparable trace.
\end{proof}

\subsection{The divided-power analyser and the missing biring data}

\begin{saresult}{B.006}{Proved}{The nonunital divided-power analyser}
Let
\[
 \Omega_{\mathrm{DP}}=\bigoplus_{n\ge1}\mathbb Z\gamma_n,
 \qquad \gamma_n=T^n/n!\in\mathbb Q[T].
\]
With ordinary polynomial addition, multiplication, and composition,
this is a nonunital analyser satisfying 153/006, axioms 1) and 2),
with composition unit $T=\gamma_1$ and constant subring zero. Its
natural domain of action is any divided-power ideal $I$ in a
commutative unital ring, by
$\left(\sum a_n\gamma_n\right)(x)=\sum a_n\gamma_n^I(x)$.
It does not satisfy the optional axiom 3a) literally printed on
153/006. The correct nonunital coaddition has its target in the
augmentation ideal of
$\mathbb Z[\gamma_1,\gamma_2,\ldots]\otimes
 \mathbb Z[\gamma_1,\gamma_2,\ldots]$, not in
$\Omega_{\mathrm{DP}}\otimes\Omega_{\mathrm{DP}}$ alone.
\end{saresult}
\begin{proof}
The product is
$\gamma_m\gamma_n=\binom{m+n}{m}\gamma_{m+n}$.
For $G=\sum_{i=1}^d a_i\gamma_i$ and $m\ge1$, the coefficient of
$\gamma_N$ in $\gamma_m(G)$ is
\[
 \sum_{\substack{r_1+\cdots+r_d=m\\\sum i r_i=N}}
 \frac{N!}{\prod_i r_i!(i!)^{r_i}}\prod_i a_i^{r_i}.
\]
The factorial quotient counts partitions of a labelled $N$-element
set into $r_i$ unordered blocks of size $i$ for each $i$, and is an
integer. Thus composition closes in $\Omega_{\mathrm{DP}}$.
Its associativity and unit follow from the injective embedding into
$\mathbb Q[T]$, and substitution in the outer variable is a
homomorphism of nonunital rings. A constant under composition must
equal its value on $0$, which is $0$ here. The divided-power axioms
give exactly the displayed product and substitution identities on
any divided-power ideal, so the stated action is defined even when
the ambient ring has torsion; division in that ring is neither
used nor asserted.

In any nonzero analyser with zero constant subring, optional 3a)
is impossible: put $G''=0$ there. Each $F_i''\circ0$ is a constant,
hence zero, so the right side vanishes, whereas for $F=T$ the left
side is $G'$. In the divided-power case the valid formula is
\[
 \gamma_n(x+y)=\gamma_n(x)+\gamma_n(y)
   +\sum_{i=1}^{n-1}\gamma_i(x)\gamma_{n-i}(y).
\]
The first two terms explain the required unitalization of the target.
For multiplication the formula
$\gamma_n(xy)=n!\gamma_n(x)\gamma_n(y)$ already lies in the
nonunital tensor square. The attached program verifies the pure
composition coefficients for $1\le m,n\le12$ and $461$ Bell
coefficients with at most six blocks of five possible sizes; the
set-partition argument, not this finite range, proves integrality
for all indices.
\end{proof}

\begin{saresult}{B.007}{Proved}{Exact extra data turning a unital analyser into a plethory}
Let $P$ be a unital analyser and let $k$ be a specified unital subring
of its constant ring. A compatible $k$-plethory structure is exactly
the following additional data and conditions: a $k$--$k$ biring
structure on the $k$-algebra $P$, with coaddition and comultiplication
realizing the finite decompositions in 3a) and 3b), with
$\beta(c)(F)=F\circ c$; and the requirement that the induced map
$P\odot_k P\longrightarrow P$, $F\odot G\longmapsto F\circ G$,
and $k[T]\longrightarrow P$, $T\longmapsto T$, are biring morphisms.
Here a biring includes additive and multiplicative counits, additive
coinverse, coassociativity, cocommutativity, codistributivity, and the
scalar map $\beta:k\longrightarrow\operatorname{Hom}_{k\text{-alg}}(P,k)$.
None of these coalgebra identities is encoded by unspecified
existence of decompositions alone.

There is a useful sufficient criterion with no choice of additional
operations: if the constant ring is a field $k$, $P$ acts faithfully
on $k$, and 1)--3) hold with the unitality convention of 153/006,
then there is a unique compatible $k$-plethory structure.
\end{saresult}
\begin{proof}
The first assertion is the explicit form of a monoid object in the
category of $k$--$k$ birings; it is the definition of Borger--Wieland
\cite[1.1, 1.3, 1.8]{BorgerWieland2005}. To check that it applies
to the given analyser, the relations defining $P\odot_k P$ are
linearity and multiplicativity in the first variable, the two
decomposition formulas in the second variable, and
$F\odot c=\beta(c)(F)$. These hold exactly by the specified
compatibility, so composition descends. The analyser monoid
identities give associativity and the two unit identities; the
remaining requirement is precisely that they are maps in the
biring category. Conversely every plethory supplies all these
data and its elementwise operations satisfy the analyser axioms.

For the sufficient criterion, identify $P$ with a $k$-subalgebra of
$k^k$ by faithful evaluation on constants. For every positive integer
$r$, the evaluation map
\[
 P^{\otimes_k r}\longrightarrow k^{k^r},\qquad
 F_1\otimes\cdots\otimes F_r\longmapsto
 \bigl((x_1,\ldots,x_r)\mapsto\prod_i F_i(x_i)\bigr)
\]
is injective. To see this for two factors, take a finite linearly
independent collection of functions in the first factor. The
evaluation vectors span its full linear dual, since a nonzero
linear combination cannot vanish at every point. Evaluating a
tensor at sufficiently many first coordinates therefore recovers
each coefficient function in the second factor. Induction proves
the assertion for every $r$.

The decompositions in 3a), 3b), evaluated on constants, now define
unique $\Delta_+F,\Delta_\times F\in P\otimes_k P$. Every biring
identity can be tested by the injective evaluation map with two,
three, or four variables, where it reduces respectively to an
ordinary ring identity in $k$. The counits are $F\mapsto F(0)$
and $F\mapsto F(1)$, the coinverse is $F\mapsto F\circ(-T)$,
and $\beta(c)$ is evaluation at $c$. The maps $\Delta_+$ and
$\Delta_\times$ are $k$-algebra maps by the same test. The
decomposition relations for composition already hold on all
elements of $P$ by hypothesis; hence composition descends to
$P\odot_k P$. Its compatibility with coaddition follows, on a
generator $F\odot G$, by evaluating both sides at $x,y\in k$:
both give $F(G(x+y))$. The multiplicative test gives $F(G(xy))$;
the counit, scalar, and coinverse tests are the corresponding
one-variable identities. Injectivity proves the biring identities.
The unit $k[T]\longrightarrow P$ has the required co-operations
because $T$ evaluates as the identity function. This completes the
construction and proves uniqueness.
\end{proof}

\begin{saresult}{B.008}{Refuted}{Freeness is insufficient for the ring extension on 153/004}
After the manuscript replaces a fraction-field inclusion by an
arbitrary inclusion $R\subset S$ of commutative rings, the marginal
claim ``OK if $\Omega$ is free'' does not ensure well-determined
coaddition or comultiplication tensors. It fails even when $R,S$
are fields. The correction ``$R$ is a domain'' in the earlier study
does not suffice either.
\end{saresult}
\begin{proof}
Take $R=\mathbb F_2$, $S=\mathbb F_4=\mathbb F_2[\alpha]$ with
$\alpha^2+\alpha+1=0$, and
\[
 \Omega=\{F\in S[T]:F(R)\subset R\},\qquad q=T^2-T.
\]
It is a free $R$-module, since $R$ is a field. Both $q$ and
$\alpha q$ belong to $\Omega$ and are linearly independent over
$R$. Thus
\[
  z=(\alpha q)\otimes q-q\otimes(\alpha q)
\]
is nonzero in $\Omega\otimes_R\Omega$. For example, extend
$q,\alpha q$ to an $R$-basis and apply the tensor product of their
two coordinate functionals. Nevertheless its image in $S[X,Y]$
is zero. A decomposition of the zero polynomial already has the
two distinct preimages $0,z$. Hence tensors in the claimed
expansions are not well-determined. The attached computation
exhausts the $16$ elements of the tensor square of
$\operatorname{span}_R(q,\alpha q)$: its polynomial image has four
elements and its kernel has four elements, including the displayed
witness. There is no finite-degree inference here: the witness in
this subspace remains nonzero in the full tensor product because
vector-space inclusions split.
\end{proof}

\begin{saresult}{B.009}{Proved}{The precise tensor criterion and the valid free-module theorem}
For an inclusion $R\subset S$, let
$\Omega=\{F\in S[T]:F(R)\subset R\}$ and let
\[
 \theta_2:\Omega\otimes_R\Omega\longrightarrow S[X,Y],
 \qquad G\otimes H\longmapsto G(X)H(Y).
\]
The polynomial identities on 153/004 determine one and only one
tensor for every $F$ if and only if $\theta_2$ is injective and both
$F(X+Y)$ and $F(XY)$ belong to its image for every $F\in\Omega$.
For a domain $R$ with fraction field $S$, freeness of $\Omega$ over
$R$ implies these conditions and, more strongly, gives isomorphisms
\[
 \theta_n:\Omega^{\otimes_R n}\xrightarrow{\sim}
 \{P\in S[T_1,\ldots,T_n]:P(R^n)\subset R\}
 \quad(n\ge1).
\]
In this case the coaddition, comultiplication, and ordinary
composition give a unique compatible $R$-plethory structure.
\end{saresult}
\begin{proof}
Existence of the requested tensors is exactly the image condition.
Uniqueness for $F=0$ already forces $\ker\theta_2=0$, and injectivity
then gives uniqueness for every $F$. This proves the exact criterion
for an arbitrary inclusion, including rings with zero divisors.

Now let $S=\operatorname{Frac}(R)$. If $R$ is finite it is a field,
$S=R$, and $\Omega=R[T]$, so the conclusion is immediate. Otherwise
$R$ is infinite. Let $\mathcal B$ be an $R$-basis of $\Omega$.
Since $R[T]\subset\Omega\subset S[T]$, localization gives
$S\otimes_R\Omega=S[T]$, and $\mathcal B$ is an $S$-basis of
$S[T]$. A polynomial $P(X,Y)$ integer-valued on $R^2$ has a unique
finite expansion $\sum_i b_i(X)h_i(Y)$ with $b_i\in\mathcal B$ and
$h_i\in S[Y]$. For every $y\in R$, the polynomial $P(X,y)$ belongs
to $\Omega$. Uniqueness of the $S$-basis expansion then implies
$h_i(y)\in R$. Therefore $h_i\in\Omega$, which proves surjectivity
of $\theta_2$. Injectivity follows because the tensor products of
the basis vectors remain linearly independent after localization.
The same argument in successive variables gives every $\theta_n$.

Transport the polynomial substitutions $T\mapsto X+Y$ and
$T\mapsto XY$ across $\theta_2$. The biring identities hold in
$S[X_1,\ldots,X_n]$, so injectivity of $\theta_n$ proves them in
the tensor powers. Evaluation at constants, negation, and the
inclusion $R[T]\subset\Omega$ provide the remaining biring maps.
Composition closes because an integer-valued polynomial sends
$R$ into $R$, and its compatibility identities are ordinary
substitution identities, again transported across the injective
$\theta_n$. Hence it defines the plethory structure. This proves
the manuscript's free-module argument under the fraction-field
hypothesis it originally had, but not under the later, weaker
hypothesis. Elliott's Theorem 4 and Proposition 13
\cite{Elliott2010} place this criterion in the more general
theory of weak polynomial compositeness. Its questions record the
state of that 2010 paper, not a claim that they remain open today;
B.016 supplies a domain where even multiplicative existence fails.
\end{proof}

\begin{saresult}{B.010}{Proved}{Faithfulness of evaluation on constants}
For any analyser $\Omega$, evaluation on its constant subring
$\Omega_0$ is injective exactly when the constants separate
elements of $\Omega$ under composition. For the polynomial analyser
of an inclusion $R\subset S$, this has the following explicit
coefficient criterion:
\[
 \bigcap_{r\in R}\ker\bigl(S^{d+1}\longrightarrow S,
 (a_0,\ldots,a_d)\longmapsto\sum_{i=0}^d a_i r^i\bigr)=0
 \quad\hbox{for every }d\ge0.
\]
In particular, if $S$ is a domain it is injective exactly when
$R$ is infinite. More generally it is injective if for every $d$
there are $r_0,\ldots,r_d\in R$ such that
$\prod_{i<j}(r_j-r_i)$ is a non-zero-divisor in $S$.
The unqualified desirability statement on 153/008 is not a theorem:
$\mathbb F_q[T]$ and the divided-power analyser B.006 both give
nonfaithful examples.
\end{saresult}
\begin{proof}
The first assertion is the kernel definition of the indicated
representation, stated to distinguish an analyser from its
possibly nonfaithful action on constants. In the polynomial case,
the constants are exactly $R$: a polynomial constant under
composition equals its composition with $0$, and therefore is
constant as a polynomial. Every polynomial vanishing on $R$ lies
in $\Omega$. Thus the kernel is the null-polynomial ideal
$\{F\in S[T]:F(r)=0\text{ for all }r\in R\}$, and collecting
coefficients by degree gives the displayed exact criterion.
A nonzero polynomial over a domain has at most its degree many
roots, proving the infinite-domain case. If $R$ is finite,
$\prod_{r\in R}(T-r)$ is a nonzero monic null polynomial, so
evaluation is not injective even if $S$ has zero divisors.
Under the stated non-zero-divisor condition, the Vandermonde
matrix on $r_0,\ldots,r_d$ has determinant
$\prod_{i<j}(r_j-r_i)$. Multiplication by its adjugate shows that
a vector in its kernel is annihilated by that determinant and
is therefore zero. This proves the sufficient condition in
general. The specific null polynomial over $\mathbb F_q$ is
$T^q-T$. B.006 has zero constant ring but contains the nonzero
element $T$, so its action on constants is also nonfaithful.
\end{proof}

\begin{saresult}{B.011}{Refuted}{A finite unital analyser satisfying all three axioms need not be a plethory}
The $64$-element algebra
\[
 P=\mathbb F_2[T]/\bigl((T^2-T)^3\bigr)
\]
with ordinary polynomial substitution is a unital analyser
satisfying 153/006, 1)--3), including both optional decomposition
axioms. It has no compatible plethory structure. Thus the missing
biring data in B.007 cannot be inferred from those axioms.
\end{saresult}
\begin{proof}
Put $q=T^2-T$. For every $G\in\mathbb F_2[T]$, the polynomial
$G^2-G$ vanishes at $0,1$, so it is divisible by $q$.
Consequently $q(G)^3$ is divisible by $q^3$; changing an outer
polynomial modulo $q^3$ does not change its substituted value.
Changing the inner polynomial by $q^3H$ also does not change the
result modulo $q^3$, by the factor theorem applied to
$F(X)-F(Y)$. Thus composition descends. Its associativity, unit
$T$, and the homomorphism property in the outer variable descend
from polynomials. The ring unit is the constant $1$.

For a representative $F(T)=\sum_{n=0}^5 a_nT^n$, expand
$F(X+Y)$ by the binomial theorem and $F(XY)$ term by term.
These expansions are finite sums of products of one-variable
polynomials. Substituting arbitrary elements $G,H\in P$ gives
exactly axioms 3a), 3b). Their dependence on the choice of a
representative is immaterial to the existential axioms.

If a compatible plethory existed, its unit map
$\mathbb F_2[T]\longrightarrow P$ would be a biring map.
Therefore $\Delta_+(T)=T\otimes1+1\otimes T$ and
$\Delta_+(q)=q\otimes1+1\otimes q$. Since $\Delta_+$ is an
algebra map and $q^3=0$ in $P$, we would obtain
\[
 0=\Delta_+(q^3)=q^2\otimes q+q\otimes q^2.
\]
This tensor is nonzero: $q,q^2$ are linearly independent in the
six-dimensional vector space $P$, and their tensor products are
independent. This is a contradiction. The attached program
exhaustively checks all $64^3=262144$ triples for composition
associativity, both homomorphism identities, and both optional
decomposition identities; it also computes the nonzero tensor.
\end{proof}

\subsection{Comparison with the primary reconstruction sources}
\begin{saresult}{B.012}{Known}{What cote 155 does and does not anticipate}
After adding B.003's unit constraint, the algebra on 155/002 is a
cocommutative equal-source-and-target instance of Takeuchi's
$\times_L$-bialgebra, presented as a pro-dual. Without that constraint,
it is not a bialgebroid, by B.002. The manuscript's reconstruction is a
special case of the general reconstruction theorems below, with a
separate completion step; its rigid case has the non-neutral groupoid
interpretation of Deligne. Thus L6 is valid as a qualified structural
comparison, and is false if it asserts that the literal displayed data
already satisfy the modern definition. No priority claim or claim of
independence follows from these comparisons.
\end{saresult}
\begin{proof}
The primary sources have been opened and the following exact portions
checked.

Takeuchi \cite[\S1, especially 1.1--1.2; Definition 4.5, pp.474--475]
{Takeuchi1977} defines the balancing-and-invariance product and requires
both the coproduct and the map to $\operatorname{End}_k(L)$ to be
algebra maps. Substituting the two equal copies of the commutative base
$L$ gives B.001. His map to $\operatorname{End}_k(L)$ is
$a\mapsto(\lambda\mapsto e(a s(\lambda)))$ in B.003, so its
multiplicativity is exactly the missing character equation. This identifies
the structure, including what the manuscript omitted; the definition
itself is not an abelian-category reconstruction theorem.

Schauenburg \cite[Theorem 5.1 and Remark 5.3; author's preprint pp.15--17]
{Schauenburg1998} proves that an algebra over $R\otimes R^{\mathrm{op}}$
has a $\times_R$-bialgebra structure precisely when its module category
has a monoidal structure with strictly monoidal underlying functor to
$R$-bimodules. Remark 5.3 removes strictness by transport of structure.
The proof reconstructs the coproduct from the regular module and the
counit from the action on the tensor unit. This is exactly B.003's
algebra-to-tensor-structure comparison before completion. For the
manuscript's pro-algebra, the regular module need not be finite, so one
must replace it by all finite regular quotients, as our proof does.
Schauenburg's theorem by itself does not reconstruct an arbitrary abelian
category from a faithful exact functor.

Hai \cite[Theorem 2.1.2, Theorem 2.2.4, Corollary 2.2.5, and
Theorem 2.2.6]{Hai2008} supplies that reconstruction by a coend: local
finiteness, exactness, faithfulness and finite projectivity give the
abelian equivalence; monoidality gives the bialgebroid; left and right
rigidity give the corresponding Hopf structures. B.003 checked local
finiteness, and finite-dimensional $L$-spaces are projective, so these
hypotheses hold here. His algebra is the coend acting by coactions;
the manuscript's $A$ is the continuous pro-dual acting by actions.
They must not be identified without this dualization or with a
finite-dimensional dual unless finiteness has actually been proved.

Finally Schauenburg \cite[\S3 and Theorems 5.1, 5.5]{Schauenburg2003}
identifies $\times_R$-bialgebras over a Frobenius-separable base with
weak bialgebras. A finite separable field extension admits that
Frobenius-separable structure; a purely inseparable extension does not.
Thus this particular weak-bialgebra reformulation has a separability
restriction, while B.003 and Deligne's transitive-groupoid theorem in
B.004 do not. None of these theorems makes the unrestricted diagonal
functor on all $L\otimes_kL$-modules faithful or exact. This separates
the three logically different uses of separability and completes the
requested comparison.
\end{proof}

\subsection{Further classifications and an explicit domain obstruction}
\begin{saresult}{B.013}{Proved}{Complete answer for a finite base ring}
Let $R\ne0$ be a finite commutative ring and let $R\subset S$ be a
unital inclusion. For the module $\Omega$ and map $\theta_2$ of B.009,
the following are equivalent: $\theta_2$ is injective; the identities
on 153/004 determine unique coaddition and comultiplication tensors;
and $S=R$. In the last case $\Omega=R[T]$ with its usual plethory
structure. This classifies arbitrary extensions of a finite base,
including extensions that are infinite or have zero divisors.
\end{saresult}
\begin{proof}
Set $P(T)=\prod_{r\in R}(T-r)$ and $d=|R|$. This is a monic
polynomial of degree $d$ vanishing on $R$. Division by $P$, valid over
any commutative coefficient ring because $P$ is monic, gives
\[
 \Omega=E\oplus P S[T],\qquad
 E=\{H\in S[T]:\deg H<d,\ H(R)\subset R\}.
\]
Indeed $F=H+PQ$ is integer-valued on $R$ exactly when $H$ is.
The summand $SP$ of $PS[T]$ is itself split: project $Q$ to its
constant coefficient. Thus $S\otimes_R S$ embeds as a direct summand
of $\Omega\otimes_R\Omega$ by $s\otimes t\mapsto(sP)\otimes(tP)$.
On this summand $\theta_2$ is
\[
 s\otimes t\longmapsto stP(X)P(Y).
\]
Injectivity of $\theta_2$ forces multiplication
$\mu:S\otimes_R S\to S$ to be injective. It is always surjective,
so it is an isomorphism.

Write $J$ for the Jacobson radical of $R$. It is nilpotent: every
$x\in J$ has repeating powers, say $x^a=x^{a+b}$, and $1-x^b$ is a
unit, whence $x^a=0$; a finite generating set of this nil ideal gives
$J^N=0$ by the pigeonhole principle. The Chinese remainder theorem
identifies $R/J$ with a finite product of finite fields. After base
change, multiplication for $T=S/JS$ over $R/J$ is still an
isomorphism. On a field factor $k$, either $T_k=0$, or choose a
$k$-linear functional $\ell:T_k\to k$ with $\ell(1)=1$.
For every $a\in T_k$, injectivity of multiplication gives
$a\otimes1=1\otimes a$. Applying $\ell\otimes1$ yields
$a=\ell(a)1$. Thus every nonzero factor equals $k$ and
$R/J\to S/JS$ is surjective. Consequently $S=R+JS$.
Successive substitution gives $S=R+J^nS$ for all $n$, and $J^N=0$
gives $S=R$. Conversely, if $S=R$ then $\Omega=R[T]$ and every
$\theta_n$ is the standard polynomial-ring isomorphism. B.009 now
proves the equivalence with the requested uniqueness statement.
This assertion concerns uniqueness of the polynomial tensor
expansions; it does not infer nonexistence of every conceivable
abstract biring structure merely from a nonzero tensor kernel.
\end{proof}

\begin{saresult}{B.014}{Proved}{All field bases and the rational-coefficient reduction}
Let $R$ be an infinite domain, $K=\operatorname{Frac}(R)$, and let
$R\subset S$ with $S$ torsion-free as an $R$-module. Inside
$K\otimes_R S$ define the overring
\[
 S_{\mathrm{rat}}=S\cap(K\cdot1).
\]
Then
\[
 \Omega(R,S)=\{F\in S_{\mathrm{rat}}[T]:F(R)\subset R\}.
\]
Moreover $\ker\theta_2$ is exactly the $R$-torsion submodule of
$\Omega\otimes_R\Omega$. In particular, for an infinite field base
$R$, every unital extension $R\subset S$ gives $\Omega=R[T]$ and
unique standard co-operations. Together with B.013, this is a
complete classification for every field base: an infinite field
allows every extension; a finite field allows unique polynomial
tensor expansions exactly for the identity extension.
\end{saresult}
\begin{proof}
Torsion-freeness embeds $S$ in $K\otimes_R S$. If $F$ has degree
at most $d$, choose distinct $r_0,\ldots,r_d\in R$. Its coefficient
vector is obtained from $(F(r_0),\ldots,F(r_d))\in R^{d+1}$ by the
inverse Vandermonde matrix over $K$. Therefore each coefficient
belongs to $K\cdot1$ as well as to $S$, proving the first equality.
The converse inclusion is immediate. Since
$R[T]\subset\Omega\subset K[T]$, localization identifies
$K\otimes_R\Omega$ with $K[T]$; after localization $\theta_2$ is the
injective standard map $K[X]\otimes_KK[Y]\to(K\otimes_RS)[X,Y]$.
Its kernel is therefore torsion. Conversely a torsion tensor maps
to zero because the target $S[X,Y]$ is $R$-torsion-free. This proves
the asserted equality. If $R=K$ is an infinite field, then
$S_{\mathrm{rat}}=R$ and every polynomial in $R[T]$ preserves $R$,
which proves the field assertion.
\end{proof}

\begin{saresult}{B.015}{Proved}{A two-variable difference obstruction}
Let $(R,\mathfrak m)$ be a local domain with fraction field $K$.
Suppose $u_n\in R$ have the following property: for every finite
set $\mathcal H\subset\operatorname{Int}(R)$ there is an $n$ with
$h(u_n)-h(0)\in\mathfrak m$ for every $h\in\mathcal H$.
If $g\in\operatorname{Int}(R)$ satisfies $g(0)=0$ and
$g(u_n^2)\notin\mathfrak m^2$ for every $n$, then
\[
 g(XY)\notin\operatorname{im}\bigl(
 \operatorname{Int}(R)\otimes_R\operatorname{Int}(R)
 \longrightarrow K[X,Y]\bigr).
\]
Consequently no compatible biring, and hence no compatible plethory,
can have its polynomial inclusion $R[T]\to\operatorname{Int}(R)$
as a biring morphism.
\end{saresult}
\begin{proof}
If $g(XY)=\sum_{i=1}^N f_i(X)h_i(Y)$, choose $u=u_n$ for the
finite collection of all the $f_i,h_i$. Take the difference between
the values at $(u,u),(u,0),(0,u),(0,0)$ with signs $+,-,-,+$.
On the right it is
\[
 \sum_i\bigl(f_i(u)-f_i(0)\bigr)
          \bigl(h_i(u)-h_i(0)\bigr)\in\mathfrak m^2.
\]
On the left it is $g(u^2)$, a contradiction. For the last assertion,
localize a putative multiplicative co-operation to $K$.
Because $K\otimes_R\operatorname{Int}(R)=K[T]$ and the unit is a
biring map, its localized value on $T$ is $XY$, so its value on
$g$ must be $g(XY)$. Its unlocalized value would be precisely a
tensor preimage whose nonexistence was just proved. This argument
requires neither injectivity of $\theta_2$ nor a general
classification theorem for integer-valued polynomial rings.
\end{proof}

\begin{saresult}{B.016}{Refuted}{Explicit one-dimensional local domains without comultiplication}
For every finite field $k=\mathbb F_q$, put
\[
 T=k[t]_{\{h:h(0)h(1)\ne0\}},\qquad m=t(t-1),\qquad
 R=k+mT\subset K=k(t).
\]
Then $R$ is a one-dimensional Noetherian local domain, but
$\operatorname{Int}(R)$ has no compatible comultiplication tensor
for a specific polynomial of degree $q^2$. In particular the
fraction-field hypothesis alone does not repair the unrestricted
claim on 153/004; freeness or another genuine tensor condition is
needed. The case $q=2$ was found in the opened primary preprint
\cite[Theorem 3.4, pp.7--8]{JiangEtAl2026}; the argument here treats
all finite fields and proves the final biring obstruction directly.
\end{saresult}
\begin{proof}
The two maximal ideals of $T$ are $tT$ and $(t-1)T$, with residue
field $k$. Their intersection is $M=mT$. Thus $R$ consists exactly
of those elements of $T$ with equal residues at the two points.
An element with nonzero common residue has its inverse in $T$ with
equal residues, hence its inverse lies in $R$. Therefore $R$ is
local with maximal ideal $M$ and $R/M=k$. It is a domain because
it is a subring of $k(t)$, and its fraction field is $k(t)$ since
$m,mt\in R$ and $t=(mt)/m$.

For the stated finiteness properties, let $A=k[m]_{(m)}$. The ring
$A[t]$ is free over $A$ on $1,t$, using $t^2=t+m$. If
$h(t)\in k[t]$ is nonzero at $0,1$, its product with $h(1-t)$
belongs to $k[m]$ and has nonzero constant term. Thus $h(t)$ is
invertible in $A[t]$, and $T=A\oplus At$. It follows that
\[
 R=A\oplus A(mt)\simeq
 A[v]/(v^2-mv-m^3).
\]
This is a finite $A$-algebra and is Noetherian. Any nonzero prime
of $R$ meets $A$ nontrivially: multiplication by a nonzero element
of the prime on the free rank-two $A$-module $R$ has nonzero
determinant, and its characteristic equation puts that nonzero
determinant in the prime. Hence every nonzero prime contracts to
$(m)$ and is the unique maximal ideal $M$. The chain
$(0)\subsetneq M$ proves that the dimension is one.

Let $v_0,v_1$ be the valuations at $t=0,1$, and set
\[
 u_n=t(t-1)^{n+1},\qquad
 a(X)=\frac{X^q-X}{m},\qquad g(X)=a(X)^q-a(X).
\]
For $x\in R$, $x^q-x\in M$, so $a(x)\in T$. Both residues of
$a(x)$ belong to $\mathbb F_q$, whence $g(x)\in M\subset R$.
Thus $g\in\operatorname{Int}(R)$ and $g(0)=0$. Since
$v_0(u_n)=1$,
\[
 v_0\bigl(a(u_n^2)\bigr)
 =v_0\bigl((u_n^{2q}-u_n^2)/m\bigr)=1,
 \qquad v_0\bigl(g(u_n^2)\bigr)=1.
\]
Every element of $M^2=m^2T$ has $v_0\ge2$, so the second hypothesis
of B.015 holds. For a finite collection of polynomials
$h(X)=\sum_r c_{h,r}X^r$, choose $n$ so large that
$v_1(c_{h,r})+r(n+1)>0$ for every nonconstant coefficient in the
collection. Then $h(u_n)-h(0)\in(t-1)T$; both values lie in $R$,
so their difference belongs to $(t-1)T\cap R=M$. This is the first
hypothesis of B.015, which proves the obstruction.

For completeness, $g$ has leading term $m^{-q}X^{q^2}$ and hence
has degree exactly $q^2$. The attached symbolic check verifies the
valuation formulas for $q=2,3,4,5,7,8,9$ and $1\le n\le12$;
the valuation argument above proves them for every prime power
and every positive integer $n$.
\end{proof}

\begin{saresult}{B.017}{Proved}{Which base rings work for every extension}
For a nonzero commutative ring $R$, the following are equivalent:
\begin{enumerate}
\item every residue field $R/\mathfrak m$ at a maximal ideal is infinite;
\item for every unital inclusion $R\subset S$, every polynomial in
$S[T]$ preserving $R$ already belongs to $R[T]$;
\item for every such inclusion, the identities of 153/004 determine
unique coaddition and comultiplication tensors.
\end{enumerate}
Thus no structural classification gate remains for bases all of whose
residue fields are infinite. This includes every ring containing an
infinite field. If $R$ has even one finite residue field, the square-zero
extension constructed below already violates uniqueness.
\end{saresult}
\begin{proof}
Assume condition 1 and let $F=\sum_{i=0}^d a_iT^i\in S[T]$
satisfy $F(R)\subset R$. Fix a maximal ideal $\mathfrak m$ of $R$.
Choose $d+1$ elements of $R$ with distinct images in the infinite
field $R/\mathfrak m$. Their Vandermonde determinant is outside
$\mathfrak m$, so the inverse matrix exists over $R_{\mathfrak m}$.
Apply it to the evaluations of $F$. Each $a_i$ has image in
$R_{\mathfrak m}\subset S_{\mathfrak m}$, so the class of $a_i$ in
the $R$-module $S/R$ vanishes after localization at every maximal
ideal. Such a class is zero: if a nonzero element $x$ of a module
had this property, its proper annihilator would be contained in
some maximal ideal $\mathfrak m$, whereas $x/1=0$ would supply
an annihilator outside $\mathfrak m$. Hence every $a_i\in R$,
proving condition 2. Condition 2 gives $\Omega=R[T]$ and the
usual tensor isomorphisms, hence condition 3.

For the converse let $k=R/\mathfrak m=\mathbb F_q$ and form
\[
 S=R\oplus k\epsilon,\qquad
 (r,a\epsilon)(s,b\epsilon)
  =(rs,(\bar r b+\bar s a)\epsilon),\qquad \epsilon^2=0.
\]
The map $r\mapsto(r,0)$ is a unital inclusion. Writing a polynomial
uniquely as $F_0+\epsilon H$ gives
\[
 \Omega=R[T]\oplus\epsilon(T^q-T)k[T],
\]
because the null polynomials on $\mathbb F_q$ are exactly the
multiples of $T^q-T$; division by this monic polynomial proves
that assertion. In particular condition 2 fails. Put
$e=\epsilon(T^q-T)$. The module $ke$ is a direct $R$-summand of
$\Omega$: in the second summand divide by $T^q-T$ and extract
the constant coefficient. Hence $e\otimes e$ is nonzero in
$\Omega\otimes_R\Omega$, since
$ke\otimes_Rke\simeq k\otimes_Rk\simeq k$.
Its polynomial image is zero because $\epsilon^2=0$.
Condition 3 therefore fails as well. This proves both converses
and the complete equivalence. The finite $R=k=\mathbb F_2$
instance is checked explicitly in \texttt{check\_B\_domains.py}.
\end{proof}

\begin{saresult}{B.018}{Proved}{The free-module repair for every torsion-free extension of an infinite domain}
Let $R$ be an infinite domain and $R\subset S$ a unital inclusion
with $S$ torsion-free over $R$. If
$\Omega=\{F\in S[T]:F(R)\subset R\}$ is free over $R$, then
for every $n\ge1$ the natural map is an isomorphism
\[
 \Omega^{\otimes_R n}\xrightarrow{\sim}
 \{P\in S[T_1,\ldots,T_n]:P(R^n)\subset R\}.
\]
Consequently the two requested tensors and the compatible plethory
structure exist uniquely. This extends B.009 from a fraction field
to every torsion-free overring of an infinite domain; B.008 shows
why the word infinite cannot simply be omitted.
\end{saresult}
\begin{proof}
By B.014 replace the coefficient ring by
$L=S_{\mathrm{rat}}\subset K=\operatorname{Frac}(R)$.
The same reduction holds in any finite number of variables:
use a rectangular grid of distinct elements of $R$ and invert
the Vandermonde matrix in each variable. Every coefficient of an
integer-valued multivariable polynomial then lies in $K\cdot1$
and hence in $L$.

Let $\mathcal B$ be an $R$-basis of $\Omega$. Localization shows
that it is a $K$-basis of $K[T]$. It is in fact an $L$-basis of
$L[T]$: independence follows after embedding in $K[T]$, and
spanning follows because every monomial $T^n$ already lies in
$\Omega$ and thus has a finite expansion with coefficients in $R$.
Now expand $P(X,Y)\in L[X,Y]$, integer-valued on $R^2$, uniquely
as $\sum_i b_i(X)h_i(Y)$ using this $L$-basis. The coefficients
$h_i$ lie in $L[Y]$. At every $y\in R$, the polynomial $P(X,y)$
belongs to $\Omega$, so uniqueness in the $K$-basis forces
$h_i(y)\in R$. Thus $h_i\in\Omega$. This gives surjectivity.
The tensors of the $R$-basis remain linearly independent over $K$,
which proves injectivity. Iteration gives the assertion for every
$n$. The same transport of polynomial identities as in B.009
then gives the biring and plethory axioms, proving the conclusion.
\end{proof}

\begin{saresult}{B.019}{Proved}{A finite dual-basis criterion, including every projective case}
In the infinite-domain, torsion-free-extension setting of B.014,
assume the following condition on the $R$-module $\Omega$: for every
finite subset $E\subset\Omega$ there are finitely many
$b_i\in\Omega$ and $R$-linear maps $\lambda_i:\Omega\to R$ such that
\[
 x=\sum_i b_i\lambda_i(x)\qquad(x\in E).
\]
Then every multivariable tensor map in B.018 is an isomorphism,
and $\Omega$ has the unique compatible plethory structure.
In particular this holds if $\Omega$ is projective over $R$;
freeness is unnecessary.
\end{saresult}
\begin{proof}
Reduce to the coefficient overring $L\subset K$ by B.014 and
its grid argument. Every $\lambda_i$ extends uniquely to a
$K$-linear functional on $K\otimes_R\Omega=K[T]$.
It maps $L[T]$ into $L$: its value on each monomial is in $R$,
and linearity gives the assertion for arbitrary coefficients in $L$.

Let $P(X,Y)\in L[X,Y]$ be integer-valued on $R^2$. Choose a nonzero
$d\in R$ clearing the denominators of its finitely many coefficient
polynomials in $X$, so those coefficients of $dP$ belong to $R[X]$
inside $\Omega$. Choose the finite dual-basis identity on this finite
set. After extension to $K$, it holds on their $K$-linear span,
so
\[
 P(X,Y)=\sum_i b_i(X)h_i(Y),\qquad
 h_i(Y)=\lambda_{i,K}\bigl(P(X,Y)\bigr).
\]
The preceding observation gives $h_i\in L[Y]$. At $y\in R$,
$P(X,y)\in\Omega$, so $h_i(y)=\lambda_i(P(X,y))\in R$.
Thus $h_i\in\Omega$, proving surjectivity.

For injectivity, let $z=\sum_j x_j\otimes y_j$ have zero polynomial
image and choose a finite dual-basis identity on the $x_j$. Then
\[
 z=\sum_i b_i\otimes\left(\sum_j\lambda_i(x_j)y_j\right).
\]
Applying $\lambda_{i,K}$ in the first variable to the zero polynomial
image gives $\sum_j\lambda_i(x_j)y_j=0$ in $K[Y]$, hence in $\Omega$.
Thus $z=0$. The same argument one variable at a time proves both
assertions for all tensor powers. Polynomial identities then give
the plethory as in B.009.

Finally a projective module is a summand of a free module. Embed
$\Omega$ in a free module and project back. The finitely many free
coordinates occurring in a finite subset $E$ give coordinate
functionals $\lambda_i$ on $\Omega$ and projected free basis vectors
$b_i$ with exactly the required identity. This verifies the
hypothesis rather than assuming that an arbitrary torsion-free
module has a dual basis.
\end{proof}

\begin{saresult}{B.020}{Proved}{Every infinite Noetherian normal base with torsion-free coefficients}
Let $R$ be an infinite Noetherian normal domain and $R\subset S$
an arbitrary $R$-torsion-free unital extension. Then $\Omega$ is a
projective $R$-module, all the multivariable tensor maps are
isomorphisms, and the unique co-operations and compatible plethory
exist. More explicitly, put $K=\operatorname{Frac}(R)$ and
$L=S_{\mathrm{rat}}\subset K$. For $n\ge0$ put
\[
 w_q(n)=\sum_{a\ge1}\left\lfloor\frac{n}{q^a}\right\rfloor,
 \qquad
 I_n=\prod_{\substack{\operatorname{ht}\mathfrak p=1\\
                     |R/\mathfrak p|=q<\infty}}
                 \mathfrak p^{-w_q(n)}.
\]
Only finitely many factors are nontrivial. There are monic
$P_n\in R[T]$ of degree $n$ such that
\[
 \operatorname{Int}(R)=\bigoplus_{n\ge0}I_nP_n,
 \qquad
 \Omega=\bigoplus_{n\ge0}(I_n\cap L)P_n.
\]
Each fractional ideal $I_n\cap L$ is invertible. These formulas
include nonprincipal ideals and therefore do not assume a regular
free polynomial basis.
\end{saresult}
\begin{proof}
We first record precisely the commutative-algebra inputs. A
Noetherian normal domain is the intersection of its height-one
localizations in its fraction field, and each such localization
is a discrete valuation ring. These are the opened Stacks Project
Lemmas 10.157.6 and 10.119.7, tags 031T and 00PD
\cite{StacksNormalDomains}. A finitely presented module whose
localizations at all maximal ideals are free is finite projective,
by its Lemma 10.78.2, tag 00NX. We apply these results only with
exactly the stated Noetherian and normal hypotheses.

There are finitely many maximal ideals with residue field of
cardinality at most $n$. To prove this, for each prime power $q\le n$
let $J_q$ be the ideal generated by $x^q-x$, $x\in R$.
Every prime quotient of $R/J_q$ is a domain in which every element
is a root of $T^q-T$, and therefore is a finite field with at most
$q$ elements. All its primes are consequently maximal. A
Noetherian ring has finitely many minimal primes, so a
zero-dimensional Noetherian ring has finitely many primes.
Every maximal ideal with residue field $\mathbb F_q$ contains
$J_q$, proving finiteness. A height-one prime with finite residue
ring is maximal. Such a prime $\mathfrak p$ is invertible: at
$\mathfrak p$ it is principal in a discrete valuation ring, and at
every other maximal ideal its localization is the whole ring.
It is finitely presented because $R$ is Noetherian, so the cited
finite-projective criterion applies. In particular every $I_n$
in the statement is an invertible fractional ideal.

Here is the needed discrete-valuation calculation with proof.
Let $V$ be a discrete valuation ring with uniformizer $\pi$ and
residue field of cardinality $q$. Choose representatives
$c_0=0,c_1,\ldots,c_{q-1}$ of its residue classes. If
$n=\sum_j n_jq^j$, define $a_n=\sum_jc_{n_j}\pi^j$ and
$Q_n(T)=\prod_{i<n}(T-a_i)$. The first unequal digits show that
$a_i\equiv a_j\pmod{\pi^r}$ exactly when
$i\equiv j\pmod{q^r}$. Among $0,\ldots,n-1$, every residue class
modulo $q^r$ occurs at least $\lfloor n/q^r\rfloor$ times, and the
class of $n$ occurs exactly that often. Counting the factors
congruent to an arbitrary $x\in V$ modulo $\pi^r$ gives
\[
 v(Q_n(x))\ge w_q(n),\qquad v(Q_n(a_n))=w_q(n).
\]
Thus $Q_n/\pi^{w_q(n)}$ is integer-valued and has unit value at
$a_n$. Successively evaluating a polynomial at $a_0,a_1,\ldots$
shows that these polynomials are a $V$-basis of
$\operatorname{Int}(V)$: subtract its degree-zero coefficient,
then its coefficient of $Q_1/\pi^{w_q(1)}$, and so forth; at step
$n$ the diagonal value is a unit. The possible leading coefficients
in degree at most $n$ are exactly $\pi^{-w_q(n)}V$.

Apply this calculation at a height-one maximal ideal
$\mathfrak p$ of $R$ with finite residue field. A uniformizer in
$R_{\mathfrak p}$ can be chosen from $R$, since one generator of
the finitely generated ideal $\mathfrak p$ has valuation one.
The residue representatives can also be chosen from $R$.
Thus every node $a_i$ above belongs to $R$. The same successive
evaluation argument therefore bounds the leading coefficient of
any $F$ with $F(R)\subset R_{\mathfrak p}$ by
$-w_q(n)$; it does not assume a localization theorem for
integer-valued polynomials. At any other height-one prime,
$R/\mathfrak p$ is infinite, and Vandermonde evaluation at $n+1$
distinct residue classes makes every coefficient integral in
$R_{\mathfrak p}$. These observations show that the leading
coefficient of every degree-$n$ element of $\operatorname{Int}(R)$
belongs to $I_n$. Indeed these are exactly its height-one valuation
bounds: if $c$ satisfies them, then $cb\in R$ for every
$b\in I_n^{-1}$ by the intersection theorem, and
$(I_n^{-1})^{-1}=I_n$ since $I_n$ is invertible.

For the reverse inclusion, let $E_n$ be the finite set of primes
contributing nontrivially to $I_n$. For each $\mathfrak p\in E_n$
construct the monic $Q_{n,\mathfrak p}\in R[T]$ just described.
By the Chinese remainder theorem choose a monic $P_n\in R[T]$
of degree $n$ whose lower coefficients are congruent to those of
$Q_{n,\mathfrak p}$ modulo $\mathfrak p^{w_q(n)}$, simultaneously
for all $\mathfrak p\in E_n$. The relevant ideal powers are
pairwise comaximal. Its values have valuation at least $w_q(n)$
at each of those primes. Consequently $cP_n$ is integer-valued
for every $c\in I_n$: it is integral at every height-one prime,
and their intersection is $R$. This proves that $I_n$ is exactly
the leading-coefficient ideal. Subtracting $cP_n$ and descending
in degree proves
$\operatorname{Int}(R)=\bigoplus I_nP_n$; uniqueness follows from
the distinct monic degrees.

Because every $P_n$ is monic with coefficients in $R$, that same
descending-degree procedure shows that its expansion has all
coefficients in $L$ precisely when the original polynomial does.
This proves $\Omega=\bigoplus (I_n\cap L)P_n$.
Set $J_n=I_n\cap L$. It is finitely presented, being a submodule of
the finite fractional ideal $I_n$ over the Noetherian ring $R$.
Localization commutes with this intersection. At a maximal ideal
outside $E_n$, it gives $(J_n)_{\mathfrak m}=R_{\mathfrak m}$.
At $\mathfrak p\in E_n$, the ring $L_{\mathfrak p}$ is an overring
of the discrete valuation ring $R_{\mathfrak p}$ in $K$, and is
either that ring or $K$. To verify this dichotomy, an element of
negative valuation, multiplied by a unit and a suitable power of
$\pi$, supplies $\pi^{-1}$, and inverting $\pi$ gives $K$.
Accordingly $(J_n)_{\mathfrak p}$ is either $R_{\mathfrak p}$ or
$\pi^{-w_q(n)}R_{\mathfrak p}$. Every localization is therefore
free of rank one, so $J_n$ is invertible by the cited criterion.
The displayed direct sum is projective. B.019 now proves all the
tensor isomorphisms and the plethory assertion.
\end{proof}

\begin{saresult}{B.021}{Proved}{An explicit infinite-rank patching lemma}
Let $V\to D$ be a surjective ring homomorphism, let $k\to D$ be a
second homomorphism, and put $R=V\times_D k$. Fix a set $I$ and an
isomorphism $\alpha:D^{(I)}\to D^{(I)}$. Define
\[
 M_\alpha=\{(b,c)\in V^{(I)}\oplus k^{(I)}:
                  \bar b=\alpha(\bar c)\}.
\]
Then $M_\alpha$ is projective over $R$, even for infinite $I$.
More precisely $M_\alpha\oplus M_{\alpha^{-1}}$ is free over $R$.
\end{saresult}
\begin{proof}
All matrices below act on direct sums, so each column has finite
support. Matrix products are well defined: applying a matrix to a
finite-support vector again gives a finite-support vector.
In the endomorphism ring of $D^{(I)}$ one has
\[
 \begin{pmatrix}\alpha&0\\0&\alpha^{-1}\end{pmatrix}
 =
 \begin{pmatrix}1&\alpha\\0&1\end{pmatrix}
 \begin{pmatrix}1&0\\-\alpha^{-1}&1\end{pmatrix}
 \begin{pmatrix}1&\alpha\\0&1\end{pmatrix}
 \begin{pmatrix}0&-1\\1&0\end{pmatrix}.
\]
Direct block multiplication proves this equality. Lift the finitely
many entries in each column of $\alpha$ and $\alpha^{-1}$ arbitrarily
through $V\to D$, choosing zero lifts of zero entries. Replace the
two off-diagonal matrices in the displayed product by these lifts.
Each elementary block matrix is invertible over $V$, regardless of
whether the lifted off-diagonal matrices are inverses. Their product
is an automorphism $\Gamma$ of $V^{(I\sqcup I)}$ reducing to
$\operatorname{diag}(\alpha,\alpha^{-1})$.

The direct sum $M_\alpha\oplus M_{\alpha^{-1}}$ is precisely the
module patched with that block diagonal gluing map. Applying
$\Gamma^{-1}$ to its $V$-component and the identity to its
$k$-component identifies it with
\[
 \{(b,c)\in V^{(I\sqcup I)}\oplus k^{(I\sqcup I)}:
                       \bar b=\bar c\}=R^{(I\sqcup I)}.
\]
The last equality is coordinatewise; all supports are finite.
Thus $M_\alpha$ is a direct summand of a free module, as claimed.
No finite-rank assumption or infinite version of Nakayama's lemma
has entered the argument.
\end{proof}

\begin{saresult}{B.022}{Proved}{The complete equal-residue-field cusp family}
Let $k=\mathbb F_q$, let $m\ge1$, put
\[
 V=k[[t]],\qquad K=k((t)),\qquad R=k+t^mV.
\]
Then $R$ is a one-dimensional Noetherian local domain with fraction
field $K$, and $A=\operatorname{Int}(R)$ is a projective $R$-module.
Consequently all its multivariable tensor maps are isomorphisms
and it has the unique compatible plethory structure.
In particular $R=\mathbb F_2[[t^2,t^3]]$ satisfies these assertions.
The conjectured nonflatness for this example recorded in the opened
Elliott paper \cite[introduction, p.2]{ElliottPresentations2011}
is false: projectivity implies flatness.
\end{saresult}
\begin{proof}
The maximal ideal is $t^mV$, and the residue field is $k$;
an element with nonzero constant term has its inverse in $R$.
The ring is finite free over $k[[t^m]]$, with basis
$1,t^{m+1},\ldots,t^{2m-1}$ (for $m=1$, just $1$).
Thus it is Noetherian and one-dimensional, and it is a domain
inside $V$. Also $t^mV\subset R$ implies
$V\subset\operatorname{Frac}(R)$, hence its fraction field is $K$.
For $m=2$, the allowed exponents are exactly $0,2,3,4,\ldots$,
which gives the asserted power-series-ring identification.

First compute $B=\operatorname{Int}(R,V)\subset K[X]$ as a
$V$-module. List $k$ as $c_0=0,c_1,\ldots,c_{q-1}$. For the
base-$q$ expansion $n=\sum_{j\ge0}n_jq^j$ put
\[
 a_n=c_{n_0}+\sum_{j\ge1}c_{n_j}t^{m+j-1},\quad
 d_n=\prod_{i<n}(a_n-a_i),\quad
 b_n(X)=\frac{\prod_{i<n}(X-a_i)}{d_n}.
\]
All nodes belong to $R$ and are distinct. Congruence modulo
$t^r$ is determined by the first digit for $1\le r\le m$,
and by the first $r-m+1$ digits for $r>m$. Counting factors
congruent to a point $x\in R$ gives
\[
 v_t\!\left(\prod_{i<n}(x-a_i)\right)\ge
 m\left\lfloor\frac nq\right\rfloor+
 \sum_{j\ge2}\left\lfloor\frac n{q^j}\right\rfloor
 =v_t(d_n).
\]
The equality at $x=a_n$ follows because its residue class has
exactly the displayed floor counts. Hence every $b_n\in B$,
$b_n(a_n)=1$, and $b_n(a_i)=0$ for $i<n$. Successive Newton
interpolation at these nodes shows
$B=\bigoplus_{n\ge0}Vb_n$: at each step the next coefficient
is a value in $V$ and the diagonal value is exactly $1$.

Choose $Q=q^e\ge m$. For every $F\in B$, its values lie in $V$,
so $F^Q$ has values in $k+t^QV\subset R$ and belongs to $A$.
Moreover $F^Q-F$ has values in $tV$, hence belongs to $tB$.
Thus $A+tB=B$. Also
\[
 A\cap tB=t^mB:
\]
a value in $R\cap tV$ lies in $t^mV$, and division of the
polynomial by $t^m$ is still a polynomial in $K[X]$ with values
in $V$. Conversely $t^mB\subset A\cap tB$ is immediate.
It follows that $C=A/t^mB$ maps isomorphically, as a $k$-vector
space, onto $B/tB$.

Multiplication defines a map
\[
 \alpha:C\otimes_k(V/t^mV)\longrightarrow B/t^mB.
\]
Choose the $k$-basis of $C$ corresponding to the reductions of
the $b_n$. Relative to these bases, $\alpha$ is a column-finite
matrix congruent to the identity modulo $t$. Write it as $1+N$
with every entry of $N$ divisible by $t$. Then $N^m=0$, and
$1-N+N^2-\cdots+(-N)^{m-1}$ is an inverse. Thus $\alpha$ is an
isomorphism by a finite nilpotent calculation, not an appeal to
Nakayama for infinitely generated modules.

The ring $R$ is the fiber product $V\times_{V/t^mV}k$, and $A$
is exactly the module patched from the free $V$-module $B$ and
the free $k$-module $C$ by this $\alpha$. Indeed a polynomial in
$B$ has reduction in $C\subset B/t^mB$ precisely when its values
lie in $R$; this is also immediate from the definition of $C$ and
$t^mB\subset A$. B.021 makes $A$ projective. B.019, applied with
coefficient field $K$, gives all the tensor isomorphisms and the
plethory. Flatness follows directly because tensoring with a
direct summand of a free module preserves injections.
\end{proof}

\begin{saresult}{B.023}{Proved}{Finite residue extensions and arbitrary conductor gaps}
Let $k\subset l$ be finite fields of cardinalities $q$ and
$Q=q^f$, let $m\ge1$, and put
\[
 V=l[[t]],\qquad K=l((t)),\qquad R=k+t^mV.
\]
Then $R$ is a one-dimensional Noetherian local domain and
$\operatorname{Int}(R)$ is projective over $R$. All multivariable
tensor maps are isomorphisms, so its unique compatible plethory
exists. In particular this proves flatness also for
$\mathbb F_2+t\mathbb F_4[[t]]$, the second conjectured nonflat
example recorded in \cite[introduction, p.2]{ElliottPresentations2011}.
Neither example is the two-branch local domain of B.016.
\end{saresult}
\begin{proof}
Locality, the maximal ideal $t^mV$, and the fraction field $K$
are verified as in B.022. If $1,e_1,\ldots,e_{f-1}$ is a
$k$-basis of $l$, then $R$ is free of rank $mf$ over $k[[t^m]]$,
with basis
\[
 1;\quad e_jt^m\ (1\le j<f);\quad
 e_jt^{m+r}\ (0\le j<f,\ 1\le r<m),\qquad e_0=1.
\]
This follows by collecting the coefficients in each residue class
of exponents modulo $m$. Hence $R$ is Noetherian. For a nonzero
element $x\in R$, multiplication by $x$ on this finite free module
has nonzero determinant. Its characteristic equation puts that
determinant in $xR$. Thus any nonzero prime of $R$ meets
$k[[t^m]]$ nontrivially and lies over its maximal ideal; since $R$
is local and integral over this ring it has dimension one.

Put $B=\operatorname{Int}(R,V)$. The Newton construction of B.022
uses mixed digits here. List the elements of $k$ as $c_0,\ldots,c_{q-1}$
and those of $l$ as $d_0,\ldots,d_{Q-1}$, both starting with zero.
Write
\[
 n=n_0+q\sum_{j\ge1}n_jQ^{j-1},\qquad
 0\le n_0<q,\quad 0\le n_j<Q,
\]
and set
\[
 a_n=c_{n_0}+\sum_{j\ge1}d_{n_j}t^{m+j-1},\qquad
 b_n(X)=\frac{\prod_{i<n}(X-a_i)}{\prod_{i<n}(a_n-a_i)}.
\]
Exactly the same residue-class counting gives the denominator
valuation
\[
 m\left\lfloor\frac nq\right\rfloor+
 \sum_{j\ge1}\left\lfloor\frac n{qQ^j}\right\rfloor,
\]
and gives at least this valuation at every point of $R$.
The diagonal Newton values are $1$, so
$B=\bigoplus_{n\ge0}Vb_n$.

We carefully distinguish polynomials from their residue functions.
Evaluation gives an injection
\[
 D_0=B/tB\longrightarrow l^R:
\]
its kernel before quotienting consists exactly of polynomials with
values in $tV$, namely $tB$. Thus every $z\in D_0$ satisfies
$z^Q=z$. The map $\sigma:z\mapsto z^q$ is an automorphism of
this ring, is semilinear for the $q$-power automorphism of $l$,
and satisfies $\sigma^f=1$. Put $C_0=D_0^\sigma$.

Here is a direct finite-field descent proof:
\[
 C_0\otimes_k l\xrightarrow{\sim}D_0.
\]
Every finitely generated $l$-subalgebra of $D_0$ is a quotient of
$l[X_1,\ldots,X_r]/(X_i^Q-X_i)$, a finite product of copies of
$l$ by the Chinese remainder theorem. It is therefore itself a
finite product of copies of $l$. On such a product $\sigma$ is
coordinatewise $q$-power, since it is the actual power operation
on elements; it fixes each idempotent and its fixed ring is the
corresponding product of copies of $k$. Scalar extension gives
the asserted isomorphism on each finite subalgebra. These
subalgebras are stable under $\sigma$ and form a filtered union
of $D_0$. Taking that union proves the displayed descent statement.
In particular no unverified assertion about Frobenius acting on
polynomial coefficients has been substituted for its action on
residue polynomial functions.

Let $A=\operatorname{Int}(R)$. Its image in $D_0$ is exactly
$C_0$. Indeed its residue values lie in $k$, which is the fixed
field of $q$-power. Conversely, if $F\in B$ has residue function
in $C_0$, choose $N=Q^e\ge m$. Then $F^N$ has values in
$k+t^NV\subset R$ and has the same residue function as $F$.
The kernel of $A\to D_0$ is $A\cap tB=t^mB$, by the same
valuewise argument as in B.022. Hence
$C=A/t^mB\simeq C_0$ as $k$-algebras.

The multiplication map
\[
 C\otimes_k(V/t^mV)\longrightarrow B/t^mB
\]
reduces modulo $t$ to the proved descent isomorphism.
Choose bases of the free modules on both sides. Their reduced
matrix is a column-finite invertible matrix over $l$; lift it
and its inverse by constant coefficients in $V/t^mV$ to change
one basis. The resulting matrix is $1+N$ with all entries in
$t(V/t^mV)$, and $N^m=0$. Its finite geometric inverse proves
that the multiplication map itself is an isomorphism. This is
valid in infinite rank because all columns and all sums in each
matrix product remain finite.

Finally $R=V\times_{V/t^mV}k$. The polynomial module $A$ is the
pullback of the free $V$-module $B$ and the free $k$-module $C$
along the last isomorphism. B.021 makes it projective over $R$;
B.019 gives the tensor isomorphisms and the plethory. This proves
the whole family and both stated special cases.
\end{proof}

\begin{saresult}{B.024}{Proved}{Every finite-conductor subring of one discrete valuation ring}
Let $V$ be a discrete valuation ring with fraction field $K$,
uniformizer $\pi$, and finite residue field of cardinality $Q$.
Let $R\subset V$ be a unital subring containing $\pi^mV$ for some
$m\ge1$. Then $R$ is a one-dimensional Noetherian local domain,
$V$ is its finite normalization, and
$A=\operatorname{Int}(R)$ is projective over $R$. Consequently
all multivariable tensor maps are isomorphisms and $A$ has the
unique compatible plethory. This includes incomplete and mixed
characteristic rings, and arbitrary conductor gaps rather than
only the pure gaps of B.023.
\end{saresult}
\begin{proof}
Put $J=\pi^mV$ and $R_0=R/J\subset D=V/J$. The ring $D$ is finite.
Its finitely many representatives generate $V$ over $R$, since
$J\subset R$. If $x\in R$ has valuation zero, multiplication by
its residue is injective on the finite ring $R_0$, hence surjective.
There is $y\in R$ with $xy=1+j$, $j\in J$. The inverse of $1+j$
in $V$ lies in $1+J\subset R$, so $x$ is a unit of $R$.
Thus $R$ is local with maximal ideal $R\cap\pi V$.
Also $V\subset\operatorname{Frac}(R)$ because $\pi^mV\subset R$.

For any nonzero ideal $H\subset R$, choose $h\in H$ of least
valuation. Then $H\subset hV$ and $JhV\subset H$.
The quotient $hV/JhV$ is finite, while $JhV$ is a finite $R$-module
because $V$ is. Finitely many representatives therefore generate
$H$ over $R$. This proves that $R$ is Noetherian. Every nonzero
prime has finite domain quotient, since it contains some $JhV$;
that quotient is a field. The nonzero maximal ideal gives dimension
one. Finally $V$ is finite, hence integral, over $R$ and is
integrally closed in $K$, so it is exactly the normalization.

Set $B=\operatorname{Int}(R,V)$. It is free over $V$ by an explicit
greedy Newton construction. Choose $a_0\in R$. Given distinct
$a_0,\ldots,a_{n-1}$, choose $a_n\in R$ minimizing the finite value
\[
 v_\pi\left(\prod_{i<n}(a_n-a_i)\right).
\]
Such a minimum exists: $R$ is infinite, so some value is finite,
and finite valuations are nonnegative integers. Put
\[
 b_n(X)=\frac{\prod_{i<n}(X-a_i)}{\prod_{i<n}(a_n-a_i)}.
\]
Minimality gives $b_n(R)\subset V$, and $b_n(a_n)=1$.
Successive evaluation at the nodes proves
$B=\bigoplus_{n\ge0}Vb_n$, exactly as in B.022.

We next compute $B/JB$ completely. Let $E_N$ be the finite set of
congruence classes modulo $\pi^NV$ met by $R$, and let
$\mathcal F(R,D)$ be the union of the function rings $D^{E_N}$
under pullback along refinement. Evaluation induces
\[
 B/JB\xrightarrow{\sim}\mathcal F(R,D).
\]
Injectivity is immediate: a polynomial with values in $J$ is
$\pi^m$ times one in $B$. Every polynomial in $B$ is constant
modulo $J$ on a sufficiently fine congruence partition. Indeed,
if its coefficients have valuations at least $-h$, then
$v(F(x)-F(y))\ge v(x-y)-h$ for $x,y\in V$.

For surjectivity we construct all congruence-class indicators.
The polynomial
\[
 \delta(X)=\frac{X^Q-X}{\pi}
\]
is integer-valued on $V$. If $v(x-y)=r\ge1$, factor its numerator
difference to obtain
\[
 v(\delta(x)-\delta(y))=r-1:
\]
the other factor reduces to $Q\bar x^{Q-1}-1=-1$ in the residue
field. Hence the residues of
$X,\delta(X),\ldots,\delta^{\circ(N-1)}(X)$ distinguish precisely
the classes modulo $\pi^N$. For a chosen class, let $c_j\in V$
represent its $j$th residue digit and put
\[
 H(X)=\prod_{j=0}^{N-1}
       \left(1-(\delta^{\circ j}(X)-c_j)^{Q-1}\right).
\]
Modulo $\pi$, this is its characteristic function on $V$.
A sufficiently high $Q$-power makes it the characteristic function
modulo $\pi^m$. Explicitly $H^{Q^m}$ suffices: when $H\in\pi V$
its power has valuation at least $Q^m\ge m$; when $H=1+z$,
$v(z)\ge1$, each $Q$-power increases $v((1+z)^Q-1)$ by at least
one. For the latter assertion, factor $Q$ as a power of the
residue characteristic $p$ and use the binomial formula for
$(1+z)^p$: every intermediate coefficient is divisible by $p$,
and the final term has valuation $p\,v(z)\ge v(z)+1$.
The argument also covers characteristic $p$, where these
intermediate coefficients vanish. These indicators, multiplied
by arbitrary representatives in $V$ and added finitely, prove
surjectivity.

Put $C=A/JB$. The ideal $JB=\pi^mB$ is contained in $A$; no
identification of $JA$ with $JB$ is assumed.
Its image is exactly $\mathcal F(R,R_0)\subset\mathcal F(R,D)$.
One implication follows from values in $R$. Conversely a
polynomial whose residue function has values in $R_0$ has all
its actual values in $R$, since $J\subset R$; and the preceding
indicator construction lifts every such function. Thus
\[
 C=A/JB\simeq\mathop{\mathrm{colim}}_N R_0^{E_N}.
\]
This is a free $R_0$-module. To see it directly, select one child
in $E_{N+1}$ over each class in $E_N$. Restriction to the selected
children retracts the pullback
$R_0^{E_N}\to R_0^{E_{N+1}}$, with free kernel on the unselected
children. The union is the direct sum of the first finite free
module and these finite free kernels. Its rank is countably
infinite, since the numbers $|E_N|$ are unbounded for infinite $R$.

Tensor product commutes with these unions and with their finite
free terms. Consequently multiplication gives an isomorphism
\[
 C\otimes_{R_0}D\xrightarrow{\sim}B/JB.
\]
The ring square is $R=V\times_D R_0$, and $A$ is the module
pulled back from $B$ and $C$ along this isomorphism. Both modules
are free of countable rank, so B.021 applies and makes $A$
projective. B.019 proves the tensor and plethory conclusions.
In particular the proof did not assume that $JA=JB$ before the
patching argument, or use completeness or equal characteristic.
\end{proof}

\begin{saresult}{B.025}{Proved}{The branch criterion for finite normalization in dimension one}
Let $(R,\mathfrak m)$ be a one-dimensional Noetherian local domain
with finite residue field, fraction field $K$, and finite
normalization $T\subset K$. Then the following are equivalent:
\begin{enumerate}
\item $T$ is local;
\item every map $\operatorname{Int}(R)^{\otimes_R n}
\to\operatorname{Int}(R^n)$ is an isomorphism;
\item the polynomial identities of 153/004 give unique coaddition
and comultiplication tensors;
\item $F(XY)$ is a sum of terms $G(X)H(Y)$ with
$G,H\in\operatorname{Int}(R)$ for every $F\in\operatorname{Int}(R)$.
\end{enumerate}
If $T$ is not local, the proof gives an explicit higher-variable
polynomial obstruction, which rules out condition 4 by finite
iteration. Thus the criterion does not merely assert failure of
projectivity or of a sufficient condition.
\end{saresult}
\begin{proof}
The finite normalization is a one-dimensional Noetherian normal
semilocal domain. Its maximal ideals
$\mathfrak p_0,\ldots,\mathfrak p_s$ all contract to
$\mathfrak m$, and all their residue fields are finite.
Its localizations are discrete valuation rings, by the opened
normal-domain facts used in B.020. If $s=0$, it is a discrete
valuation ring. A nonzero conductor ideal exists: clear the
denominators of finitely many $R$-module generators of $T$.
In a discrete valuation ring it contains $\pi^mT$ for some $m$.
B.024 therefore proves condition 2. The implications
$2\Rightarrow3\Rightarrow4$ are immediate from B.009.

Suppose now $s\ge1$. We verify the conductor and approximation
steps explicitly. For each $i$, choose $\pi_i\in T$ congruent
to a valuation-one element modulo $\mathfrak p_i^2$, and
congruent to $1$ modulo every other maximal ideal. The Chinese
remainder theorem supplies it. Thus
$v_i(\pi_i)=1$ and $v_j(\pi_i)=0$ for $j\ne i$.
Localizing at all maximal ideals shows
$\mathfrak p_i=\pi_iT$; no general claim that a nonprincipal
ideal has a chosen generator is being used.
Write the nonzero conductor as
\[
 I=(R:T)=\prod_i\mathfrak p_i^{c_i}=aT,
 \qquad a=\prod_i\pi_i^{c_i}.
\]
The equality follows by localizing its two sides at the finitely
many discrete valuation rings. Every $c_i$ is positive. Otherwise
some conductor element lies outside $\mathfrak p_i$, hence outside
$\mathfrak m$ in $R$, so it is a unit of $R$. The conductor would
then contain $1$ and force $T=R$, contrary to $s\ge1$.
In particular $a\in I\subset R$ and $v_0(a)=c_0$.

Set $r=c_0+1$ and choose $N>r c_0$.
The unit $\pi_1$ has finite order $L$ in
$(T/\mathfrak p_0^N)^*$, since this quotient is finite.
For $n\ge1$ put
\[
 u_n=a\pi_1^{Ln}\in I\subset R.
\]
Then $u_n\equiv a\pmod{\mathfrak p_0^N}$ and
$v_1(u_n)\longrightarrow\infty$. For every finite family
$\mathcal H\subset\operatorname{Int}(R)$, sufficiently large $n$
gives
\[
 h(u_n)-h(0)\in\mathfrak p_1\cap R=\mathfrak m
 \qquad(h\in\mathcal H).
\]
Indeed every nonconstant term of each of the finitely many
polynomials tends to zero for $v_1$, and both
values lie in $R$.

Use the digit-indicator construction in B.024 at the discrete
valuation ring $T_{\mathfrak p_0}$ to produce a polynomial $H$
which is $1$ modulo $\mathfrak p_0$ on the class of $a^r$ modulo
$\mathfrak p_0^N$, and $0$ on the zero class. These classes are
distinct because $v_0(a^r)=r c_0<N$. In that construction the
residue representatives can be chosen in $T$, and its only
denominators are powers of $\pi_0$. Thus
$H\in T[\pi_0^{-1}][X]$, it is integer-valued in
$T_{\mathfrak p_0}$, and its coefficients are integral at every
other maximal ideal of $T$. Consequently $H(R)\subset T$.
Put
\[
 F=H-H(0),\qquad g=aF.
\]
Then $F(R)\subset T$, $F(0)=0$, and
$g(R)\subset aT=I\subset R$. Also
\[
 v_0(g(u_n^r))=c_0\quad\hbox{for every }n,
\]
since $u_n^r\equiv a^r\pmod{\mathfrak p_0^N}$ and $F$ has
unit residue on this class.

The polynomial $g(X_1\cdots X_r)$ is not in the image of the
$r$-fold tensor map. Otherwise write it as a finite sum of
products $\prod_{j=1}^r h_{ij}(X_j)$ and choose $n$ large for
all these factors in the displayed difference condition.
Its iterated difference on the $2^r$ points with coordinates
$0,u_n$ would be
\[
 \sum_i\prod_{j=1}^r(h_{ij}(u_n)-h_{ij}(0))\in\mathfrak m^r.
\]
The other expression gives exactly $g(u_n^r)$, because $g(0)=0$.
Every element of $\mathfrak m^r$ has $v_0\ge r>c_0$, a
contradiction. Finally, if condition 4 held, repeated binary
expansion, first separating $X_1$ from $X_2\cdots X_r$ and
then expanding every remaining factor, would put
$g(X_1\cdots X_r)$ in that image. The expansions are finite at
each of the $r-1$ steps, so no coherence or injectivity assumption
is needed for this implication. Thus condition 4 fails.
This proves all the equivalences and the negative direction.
\end{proof}

\begin{saresult}{B.026}{Proved}{All coefficient extensions for the one-branch local class}
Let $R\subset V$ satisfy B.024, with fraction field $K$, and let
$R\subset S$ be any $R$-torsion-free unital extension. Put
$L=S_{\mathrm{rat}}\subset K$ as in B.014. Then the identities
on 153/004 determine unique coaddition and comultiplication tensors
if and only if
\[
 L=R\quad\hbox{or}\quad L=K.
\]
In both cases the compatible plethory exists. A proper intermediate
coefficient overring $R\subsetneq L\subset K$ has a noninjective
tensor map. Thus even in this one-branch class the answer depends
on the specified coefficient ring, and does not grow monotonically
with that ring.
\end{saresult}
\begin{proof}
B.014 reduces the problem to $R\subset L\subset K$.
Any overring $L$ not contained in $V$ is $K$. Indeed, take
$b\in L$ of negative valuation and choose $N$ so large that
$c=b^{-N}\in\pi^mV\subset R$. Then $c^{-1}=b^N\in L$.
For every $x\in K$, a sufficiently high power of $c$ sends $x$
into $\pi^mV\subset R$. Hence $K=R[c^{-1}]\subset L$.
Every remaining overring is contained in the finite $R$-module
$V$, so it is finite over the Noetherian ring $R$.

Suppose $R\subsetneq L\subset V$. The conductor
$J=(R:L)$ contains $\pi^mV$, so $R/J$ is finite. Choose one
representative $r_i$ of each of its classes and put
$P(T)=\prod_i(T-r_i)$. It is monic and $P(R)\subset J$.
Therefore $P L[T]\subset\Omega(R,L)$, and monic division gives
\[
 \Omega(R,L)=E\oplus PL[T],\qquad
 E=\{H\in L[T]:\deg H<\deg P,\ H(R)\subset R\}.
\]
In particular $LP$ is a direct $R$-summand. Injectivity of
$\theta_2$ would force multiplication
$L\otimes_RL\to L$ to be injective, by restricting to this
summand, and hence an isomorphism.

A finite unital extension of a local ring with this property and
injective structure map is the identity. Here is a proof. After
base change to $k=R/\mathfrak m$, the finite-dimensional algebra
$C=L/\mathfrak mL$ still has injective multiplication
$C\otimes_kC\to C$. It is nonzero by Nakayama's elementary
finite-module argument. Choose a linear functional on $C$ taking
$1$ to $1$. Applying it to $a\otimes1=1\otimes a$ shows $C=k$.
Thus $L=R+\mathfrak mL$, so the finite module $L/R$ equals
$\mathfrak m(L/R)$ and is zero. For completeness, the finite-module
argument used twice says that if generators $x_i$ satisfy
$x_i=\sum a_{ij}x_j$ with $a_{ij}\in\mathfrak m$, the adjugate
matrix of $1-(a_{ij})$ annihilates the generators by a determinant
congruent to $1$ modulo $\mathfrak m$, hence a unit. This proves
$L=R$, a contradiction. Therefore the tensor map is not injective.

If $L=R$, the answer is the usual polynomial plethory.
If $L=K$, it is B.024. These cases exhaust all overrings and
prove the equivalence.

For a fully explicit noninjective example take
$R=\mathbb F_2+t^2\mathbb F_2[[t]]$ and $L=V=\mathbb F_2[[t]]$.
With $q(T)=T^2-T$ one has
\[
 \Omega(R,V)=R\oplus RT\oplus q(T)V[T].
\]
Division by $q$ proves this decomposition; the remainder has
coefficients in $R$ by evaluating at $0,1$. The tensor
\[
 z=(tq)\otimes q-q\otimes(tq)
\]
has zero polynomial image but is nonzero in the split summand
$Vq\otimes_RVq$. After tensoring this summand with $R/\mathfrak m$
it maps to $\tau-\sigma$ in
$\mathbb F_2[\tau,\sigma]/(\tau^2,\sigma^2)$, which is nonzero.
The attached exact program checks all $16$ elements of this finite
tensor algebra and its multiplication-to-the-diagonal map.
By contrast, with coefficient field $K$ the same base ring works
by B.022. This verifies the asserted dependence on coefficients.
\end{proof}


\begin{saresult}{B.027}{Proved}{Localization of every multivariable polynomial module in dimension one}
Let $R$ be a one-dimensional Noetherian domain, $K$ its fraction
field, and $R\subset L\subset K$ an arbitrary overring. For
$n\ge1$ put
\[
 \Omega_n(R,L)=\{F\in L[X_1,\ldots,X_n]:F(R^n)\subset R\}.
\]
For every maximal ideal $\mathfrak m$ one has, as submodules of
$K[X_1,\ldots,X_n]$,
\[
 \Omega_n(R,L)_{\mathfrak m}
   =\Omega_n(R_{\mathfrak m},L_{\mathfrak m}).
\]
In particular the tensor maps in B.009 localize to the corresponding
maps for $(R_{\mathfrak m},L_{\mathfrak m})$. No polynomial
localization property is assumed in this statement.
\end{saresult}
\begin{proof}
We first prove the required density of inputs. If $0\ne d\in R$,
then either $d$ is a unit in $R_{\mathfrak m}$ or
\[
 \sqrt{dR_{\mathfrak m}}=\mathfrak mR_{\mathfrak m}.
\]
Indeed every nonzero prime of the one-dimensional local domain
$R_{\mathfrak m}$ is its maximal ideal. Because that ideal is
finitely generated, some power
$\mathfrak m^NR_{\mathfrak m}$ is contained in
$dR_{\mathfrak m}$. Also
\[
 R\longrightarrow R_{\mathfrak m}/\mathfrak m^NR_{\mathfrak m}
\]
is surjective: the ring $R/\mathfrak m^N$ has the unique maximal
ideal $\mathfrak m/\mathfrak m^N$, so every denominator outside
$\mathfrak m$ is already invertible in this quotient.
It follows that every element of $R_{\mathfrak m}$ can be
approximated by an element of $R$ modulo $dR_{\mathfrak m}$.
The same assertion is immediate if $d$ is a unit.

Take $F\in\Omega_n(R,L)$ and choose $d\ne0$ with
$dF\in R[X_1,\ldots,X_n]$. Given
$x\in R_{\mathfrak m}^n$, choose $r\in R^n$ with
$x_j-r_j\in dR_{\mathfrak m}$ for each $j$.
The difference of the values of the polynomial $dF$ is a sum
of terms divisible by these coordinate differences. Thus
\[
 dF(x)-dF(r)\in dR_{\mathfrak m},
 \qquad F(x)-F(r)\in R_{\mathfrak m}.
\]
As $F(r)\in R$, we have $F(x)\in R_{\mathfrak m}$.
Dividing $F$ by any element outside $\mathfrak m$ preserves this
conclusion and puts its coefficients in $L_{\mathfrak m}$.
This proves the first inclusion.

Conversely take
$G\in\Omega_n(R_{\mathfrak m},L_{\mathfrak m})$.
Clear the finitely many coefficient denominators outside
$\mathfrak m$: there is $t\in R\setminus\mathfrak m$ with
$F=tG\in L[X_1,\ldots,X_n]$. Choose $d\ne0$ with
$dF\in R[X_1,\ldots,X_n]$ and form the ideal
\[
 H=(dF(r):r\in R^n)\subset R.
\]
It is finitely generated because $R$ is Noetherian, and
$H_{\mathfrak m}\subset dR_{\mathfrak m}$ because $F(r)$ lies
in $R_{\mathfrak m}$. For each of finitely many generators
$h_i$ choose $u_i\notin\mathfrak m$ with $u_i h_i\in dR$.
Their product $u$ satisfies $uH\subset dR$. Therefore
$uF(R^n)\subset R$, so $uF\in\Omega_n(R,L)$ and
\[
 G=\frac{uF}{ut}\in\Omega_n(R,L)_{\mathfrak m}.
\]
This proves the reverse inclusion with a uniform denominator,
not a separate denominator for each value. Localization commutes
with finite tensor products; all polynomial multiplication maps
are the evident ones inside $K[X_1,\ldots,X_n]$. This gives the
last assertion.
\end{proof}

\begin{saresult}{B.028}{Proved}{The global finite-normalization criterion}
Let $R$ be a one-dimensional Noetherian domain with fraction field
$K$ and finite normalization $T$ in $K$. The following are
equivalent:
\begin{enumerate}
\item for every maximal ideal $\mathfrak m$ with finite residue
field, $T_{R\setminus\mathfrak m}$ is local, equivalently exactly
one maximal ideal of $T$ lies above $\mathfrak m$;
\item all maps $\operatorname{Int}(R)^{\otimes_R n}
\to\operatorname{Int}(R^n)$ are isomorphisms;
\item the polynomial coaddition and comultiplication of 153/004
have unique tensor representatives;
\item $F(XY)$ has a finite expansion
$\sum_iG_i(X)H_i(Y)$ with $G_i,H_i\in\operatorname{Int}(R)$
for every $F\in\operatorname{Int}(R)$.
\end{enumerate}
When these hold, $\operatorname{Int}(R)$ has its unique compatible
plethory. Thus, for fraction-field coefficients, B.2(c) is completely
settled in dimension one with finite normalization. Only the
maximal ideals with finite residue field enter the criterion.
\end{saresult}
\begin{proof}
For any maximal ideal $\mathfrak m$, localization of the finite
normalization is the finite normalization of $R_{\mathfrak m}$.
Here is the integrality check. It is finite, hence integral, over
$R_{\mathfrak m}$ and is integrally closed as a localization of
$T$. Conversely if $x\in K$ is integral over $R_{\mathfrak m}$,
write a monic equation with its finitely many denominators outside
$\mathfrak m$. Choose their common multiple $s\notin\mathfrak m$.
Multiplying the equation for $x$ by $s^n$ gives a monic equation
for $sx$ over $R$, because each coefficient of $(sx)^i$ is the
original coefficient times $s^{n-i}$. Thus $sx\in T$ and
$x\in T_{R\setminus\mathfrak m}$. The maximal ideals of this
localized integral extension correspond exactly to those of $T$
above $\mathfrak m$.

Assume condition 1. If $R/\mathfrak m$ is infinite, the local
Vandermonde argument of B.017 gives
\[
 \operatorname{Int}(R_{\mathfrak m})=R_{\mathfrak m}[X],
 \qquad
 \operatorname{Int}(R_{\mathfrak m}^n)
       =R_{\mathfrak m}[X_1,\ldots,X_n].
\]
For the multivariable statement apply the same invertible
Vandermonde matrix successively in each variable. All local
tensor maps are therefore isomorphisms. If the residue field
is finite, condition 1 and B.025 give the same conclusion.
By B.027 the kernel and cokernel of each global tensor map
localize to zero at every maximal ideal, so they are zero.
Indeed a nonzero element of any module has a proper annihilator;
a maximal ideal containing it gives a localization in which
that element remains nonzero. This proves condition 2.
The implications $2\Rightarrow3\Rightarrow4$ follow from B.009.

Finally assume condition 4 and fix a maximal ideal
$\mathfrak m$ with finite residue field. B.027 writes any
$G\in\operatorname{Int}(R_{\mathfrak m})$ as $F/s$, with
$F\in\operatorname{Int}(R)$ and $s\notin\mathfrak m$.
An expansion of $F(XY)$, divided by $s$, gives an expansion of
$G(XY)$ over $R_{\mathfrak m}$. B.025 then forces its finite
normalization to be local. Thus condition 1 holds.
The plethory identities follow from the all-variable tensor
isomorphisms exactly as in B.009; no independent gluing of
operations or an unstated localization theorem is needed.
\end{proof}

\begin{saresult}{B.029}{Proved}{Coefficient extensions over the global one-branch class}
Let $R$ and $T$ satisfy B.028, and suppose the equivalent conditions
of that result hold. For any $R$-torsion-free unital extension
$R\subset S$, put $L=S_{\mathrm{rat}}\subset K$ as in B.014.
Then the following are equivalent:
\begin{enumerate}
\item $\Omega(R,S)$ has uniquely determined polynomial
coaddition and comultiplication tensors;
\item all its multivariable tensor maps are isomorphisms;
\item for every maximal ideal $\mathfrak m$ with finite residue
field, one has $L_{\mathfrak m}=R_{\mathfrak m}$ or
$L_{\mathfrak m}=K$.
\end{enumerate}
Whenever these hold the compatible plethory exists. The choices
$R_{\mathfrak m}$ and $K$ may differ at different maximal ideals.
\end{saresult}
\begin{proof}
B.014 identifies the polynomial module with $\Omega(R,L)$.
At a maximal ideal with infinite residue field it localizes,
by B.027 and the Vandermonde argument, to the ordinary polynomial
module over $R_{\mathfrak m}$; this works for any overring
$L_{\mathfrak m}$. At a maximal ideal with finite residue field,
the normalized local ring is a single discrete valuation ring
and its nonzero conductor contains a power of its maximal ideal.
Thus B.026 applies. Its proof gives two facts: the local tensor
maps are all isomorphisms for $L_{\mathfrak m}=R_{\mathfrak m}$
or $K$, and the second tensor map is noninjective in every other
case. A global injective tensor map stays injective on localization,
so condition 1 implies condition 3. Conversely condition 3 makes
all localized tensor maps isomorphisms; B.027 and the kernel and
cokernel argument in B.028 make them global isomorphisms.
The remaining implication and the plethory assertion are B.009.
\end{proof}


\begin{saresult}{B.030}{Proved}{The finite-extension obstruction with a cofinite conductor}
Let $R\subset S$ be an injective unital extension of nonzero
commutative rings, with $R$ local and $S$ finite as an $R$-module.
Suppose its conductor
\[
 J=\{r\in R:rS\subset R\}
\]
has finite quotient $R/J$. Then the second polynomial tensor map
for $\Omega(R,S)$ is injective if and only if $S=R$.
Consequently unique polynomial coaddition and comultiplication
occur exactly in that case.
\end{saresult}
\begin{proof}
Choose one representative of each class of $R/J$ and let $P$ be
the monic product of the corresponding linear factors. Then
$P(R)\subset J$, so $PS[X]\subset\Omega(R,S)$.
Division by the monic polynomial $P$ shows
\[
 \Omega(R,S)=E\oplus PS[X],
 \quad E=\{F\in\Omega(R,S):\deg F<\deg P\}.
\]
Thus $SP$ is a direct $R$-summand. Restriction of the tensor map
to $SP\otimes_RSP$ is multiplication
$S\otimes_RS\to S$, followed by multiplication by the monic
polynomial $P(X)P(Y)$. Injectivity of the tensor map would make
$S\otimes_RS\to S$ injective and hence an isomorphism.

Let $k$ be the residue field of $R$. The finite-dimensional
$k$-algebra $C=S\otimes_Rk$ is nonzero by the finite-module
Nakayama argument given explicitly in B.026. Its multiplication
would be an isomorphism $C\otimes_kC\to C$, so
$(\dim_k C)^2=\dim_k C$ and $C=k$.
It follows that the finite module $S/R$ satisfies
$S/R=\mathfrak m(S/R)$ and is zero by the same argument.
Thus $S=R$. Conversely for $S=R$ the module is $R[X]$ and all
polynomial tensor maps are isomorphisms.
\end{proof}

\begin{saresult}{B.031}{Proved}{All local branches and all torsion-free coefficient extensions}
Let $(R,\mathfrak m)$ be a one-dimensional Noetherian local domain
with finite residue field, fraction field $K$, and finite
normalization $T$. Let $R\subset S$ be any $R$-torsion-free
unital extension, and put $L=S_{\mathrm{rat}}$.
Then the polynomial coaddition and comultiplication tensors for
$\Omega(R,S)$ are uniquely determined if and only if
\[
 L=R\quad\hbox{or}\quad (T\hbox{ is local and }L=K).
\]
In these cases every multivariable tensor map is an isomorphism
and the unique compatible plethory exists. If $T$ has more than
one maximal ideal and $L\not\subset T$, even the existence of
multiplicative tensors fails.
\end{saresult}
\begin{proof}
Reduce to $\Omega(R,L)$ by B.014. If $T$ is local this is B.026,
using the nonzero conductor furnished by finite normalization.
Suppose $T$ has at least two maximal ideals. The notation and
conductor computation of B.025 apply:
\[
 \mathfrak p_i=\pi_iT,\qquad
 I=(R:T)=aT,\qquad v_i(a)=c_i>0.
\]
Every $T/\mathfrak p_i^{c_i}$ is finite, so $T/I$ and $R/I$
are finite. If $R\subsetneq L\subset T$, then $L$ is a finite
$R$-module because $R$ is Noetherian and $T$ is finite.
Its conductor contains $I$, and B.030 makes its second tensor
map noninjective. This handles every proper intermediate
coefficient ring contained in the normalization.

Now suppose $L\not\subset T$ and choose $b\in L\setminus T$.
At least one valuation $v_i(b)$ is negative. Define
\[
 c=\prod_{v_j(b)<0}\pi_j^{-v_j(b)},\qquad d=bc.
\]
Both $c$ and $d$ lie in $T$: all their valuations are nonnegative,
and $T$ is the intersection of its localizations inside $K$.
The latter assertion also follows directly from the ideal of
denominators: if $x\notin T$, the proper ideal
$\{t\in T:tx\in T\}$ is contained in a maximal ideal whose
localization still excludes $x$. At each maximal ideal at least
one of $c,d$ is a unit. Hence $(c,d)=T$, and there are $p,q\in T$
with $pc+qd=1$. Therefore
\[
 \frac1c=p+qb\in T[b],\qquad
 \pi_i^{-1}=(c/\pi_i)c^{-1}\in T[b]
       \quad\hbox{if }v_i(b)<0.
\]
This proves the needed inverse inclusion; negative valuation
alone has not been mistaken for a globally invertible element.

Relabel such an index as $0$ and choose another branch as $1$.
The construction in B.025 produces
\[
 H\in T[\pi_0^{-1}][X],\qquad
 g=a\bigl(H-H(0)\bigr)\in\operatorname{Int}(R),
\]
for which $g(X_1\cdots X_{c_0+1})$ has no tensor expansion even
with factors from the full module $\operatorname{Int}(R)$.
Every coefficient of $g$ lies in $aT[b]\subset R[b]\subset L$,
because $aT\subset R$. Thus $g\in\Omega(R,L)$.
If all multiplicative tensors existed for this smaller module,
repeated binary expansion would give the forbidden finite
$(c_0+1)$-variable expansion, exactly as in B.025.
They therefore do not all exist.

The only remaining possibility for a multibranch normalization
is $L=R$, for which the module is $R[X]$. Combined with the
one-branch case this proves the assertion in full.
\end{proof}

\begin{saresult}{B.032}{Proved}{Complete coefficient classification in dimension one with finite normalization}
Let $R$ be a one-dimensional Noetherian domain, let $K$ be its
fraction field, and suppose its normalization $T$ in $K$ is a
finite $R$-module. For any $R$-torsion-free unital extension
$R\subset S$, put $L=S_{\mathrm{rat}}\subset K$.
The following are equivalent:
\begin{enumerate}
\item the polynomial identities on 153/004 determine unique
coaddition and comultiplication tensors for $\Omega(R,S)$;
\item every multivariable tensor map for $\Omega(R,S)$ is an
isomorphism onto the corresponding integer-valued polynomials;
\item for every maximal ideal $\mathfrak m$ with finite residue
field, either
\[
 L_{\mathfrak m}=R_{\mathfrak m},
 \quad\hbox{or}\quad
 L_{\mathfrak m}=K\ \hbox{ and }\ T_{R\setminus\mathfrak m}
       \hbox{ is local}.
\]
\end{enumerate}
These conditions also give the unique compatible plethory.
There is no restriction at a maximal ideal with infinite residue
field. The criterion includes all coefficient overrings and all
normalization branch counts in the stated class.
\end{saresult}
\begin{proof}
The rational-coefficient reduction of B.014 works in every finite
number of variables by applying its Vandermonde matrices
successively in each variable. Thus all modules and tensor maps
can be computed with $L$ in place of $S$.
B.027 identifies their localizations at every maximal ideal.
At a maximal ideal with infinite residue field the Vandermonde
argument of B.017 gives the usual polynomial module and tensor
isomorphisms. At a maximal ideal with finite residue field,
B.031 gives exactly the alternatives in condition 3.

Both injectivity and the existence of the two substituted
polynomial tensors pass from the global module to its localization:
write a local polynomial as $F/s$ by B.027 and divide a tensor
expansion of $F(X+Y)$ or $F(XY)$ by $s$.
Therefore condition 1 forces all the local alternatives and
implies condition 3. Conversely condition 3 makes every local
tensor map an isomorphism. Its global kernel and cokernel vanish
by the localization test proved in B.028, so condition 2 follows.
B.009 gives condition 1 and all the plethory identities from
condition 2. This proves the full equivalence.
\end{proof}


\begin{saresult}{B.033}{Proved}{Polynomial localization in every Noetherian dimension}
The conclusion of B.027 holds for every Noetherian domain $R$,
without a dimension restriction, and for localization at any
multiplicatively closed subset $U\subset R\setminus\{0\}$:
\[
 U^{-1}\Omega_n(R,L)=\Omega_n(U^{-1}R,U^{-1}L).
\]
The inclusion from left to right holds for every domain, without
Noetherianity.
\end{saresult}
\begin{proof}
We prove the inclusion by induction on the total degree $d$ of
$F\in\Omega_n(R,L)$. A constant such polynomial belongs to $R$.
For $s\in U$ set
\[
 G_s(X_1,\ldots,X_n)
   =s^dF(X_1,\ldots,X_n)-F(sX_1,\ldots,sX_n).
\]
All terms of total degree $d$ cancel; hence $G_s$ has smaller
degree and still takes $R^n$ to $R$. Every tuple in $(U^{-1}R)^n$
can be written $a/s$ with a common $s\in U$ and $a\in R^n$.
The induction hypothesis and
\[
 s^dF(a/s)=G_s(a/s)+F(a)
\]
show that $F(a/s)\in U^{-1}R$. Dividing $F$ by an additional
element of $U$ preserves this property. This proves the inclusion
without any approximation or dimension assumption.

For the converse, take a polynomial with coefficients in
$U^{-1}L$ and values in $U^{-1}R$. Clear its finitely many
coefficient denominators by $t\in U$ to obtain $F\in L[X]$,
where $X$ denotes all $n$ variables. Choose $d\ne0$ with
$dF\in R[X]$. The ideal $(dF(r):r\in R^n)$ is finitely generated
and is contained in $dU^{-1}R$ after localization.
Multiplying by one element $u\in U$ puts every generator in $dR$,
so $uF\in\Omega_n(R,L)$. This is exactly the uniform-denominator
argument of B.027, and proves the converse.
\end{proof}

\begin{saresult}{B.034}{Proved}{The punctured-spectrum coefficient ring and the depth test}
Let $(R,\mathfrak m)$ be a nonfield Noetherian local domain with
fraction field $K$. Define the explicitly specified overring
\[
 R^\sharp=\{x\in K:\mathfrak m^Nx\subset R
                    \text{ for some }N\ge0\}.
\]
Then
\[
 R^\sharp=\bigcap_{\mathfrak p\ne\mathfrak m}R_{\mathfrak p}
 \quad\hbox{inside }K,
 \qquad
 R^\sharp=R\ \Longleftrightarrow\ \operatorname{depth}R\ge2.
\]
If the residue field is finite, and $R\subset L\subset K$ is an
overring, put $E=L\cap R^\sharp$. Then
\[
 \Omega_n(R,L)=\Omega_n(R,E)\quad(n\ge1),
 \qquad
 \Omega_1(R,L)=R[X]\ \Longleftrightarrow\ E=R.
\]
Thus the ring controlling whether nonpolynomial operations can
occur is determined without referring to tensor maps.
\end{saresult}
\begin{proof}
The displayed set is closed under addition and multiplication,
using a common sufficiently high power of $\mathfrak m$.
For $x\in K$ its nonzero denominator ideal
$I_x=\{r\in R:rx\in R\}$ detects membership in every
$R_{\mathfrak p}$. If $x$ belongs to all localizations with
$\mathfrak p\ne\mathfrak m$, either $I_x=R$ or
$\sqrt{I_x}=\mathfrak m$. Noetherianity then gives
$\mathfrak m^N\subset I_x$. Conversely if
$\mathfrak m^Nx\subset R$ and $\mathfrak p\ne\mathfrak m$,
choose $a\in\mathfrak m\setminus\mathfrak p$; the equality
$x=(a^Nx)/a^N$ places $x$ in $R_{\mathfrak p}$.
This proves the intersection formula.

Suppose $a,b\in\mathfrak m$ is a regular sequence. If
$R^\sharp\ne R$, choose $x\notin R$ with
$\mathfrak m x\subset R$: starting with any element of
$R^\sharp\setminus R$, use a least annihilating power and a
suitable element of the preceding power. Then $ax,bx\in R$ and
$b(ax)=a(bx)$. The nonzerodivisor $b$ on $R/aR$ forces
$ax\in aR$; cancellation by nonzero $a$ gives $x\in R$,
a contradiction. Hence depth at least two implies $R^\sharp=R$.
Conversely, if no regular sequence of length two exists, choose
$0\ne a\in\mathfrak m$. Every element of $\mathfrak m$ is a
zerodivisor on $R/aR$. The zerodivisors of this finite Noetherian
module form the union of its finitely many associated primes.
Prime avoidance forces $\mathfrak m$ to be one of those primes.
Thus there is $z\notin aR$ with $\mathfrak m z\subset aR$;
$x=z/a$ belongs to $R^\sharp\setminus R$.
The associated-prime facts used in this short argument are the
opened Stacks Project, Section 10.63, Lemmas 10.63.5 and 10.63.9,
and the depth-zero formulation in Lemma 10.72.5.

For the coefficient assertion take $F\in\Omega_n(R,L)$ and a
prime $\mathfrak p\ne\mathfrak m$. The domain $R/\mathfrak p$
is not a field and is therefore infinite. Choose as many elements
of $R$ with distinct residues modulo $\mathfrak p$ as the degree
in each variable requires. The corresponding Vandermonde
matrices are invertible over $R_{\mathfrak p}$. Applying them
to the values of $F$ proves that every coefficient lies in
$R_{\mathfrak p}$. The intersection formula puts every coefficient
in $E$, proving the module equality. If $E=R$ this is the usual
polynomial ring.

Conversely let $x\in E\setminus R$. Choose $N$ with
$\mathfrak m^Nx\subset R$. The quotient $R/\mathfrak m^N$ is
finite: its finite filtration has finite-dimensional vector-space
quotients over the finite residue field. The monic polynomial
$P$ formed from all of its residue representatives satisfies
$P(R)\subset\mathfrak m^N$. Hence $xP\in\Omega_1(R,L)$, but
its leading coefficient $x$ is not in $R$. This proves the last
equivalence.
\end{proof}

\begin{saresult}{B.035}{Proved}{The higher-dimensional local classification with finite normalization}
Let $(R,\mathfrak m)$ be a Noetherian local domain of dimension
at least two with finite residue field, fraction field $K$, and
finite normalization $T$. Let $R\subset L\subset K$ and set
$E=L\cap R^\sharp$ as in B.034. The following are equivalent:
\begin{enumerate}
\item the polynomial coaddition and comultiplication tensors
for $\Omega(R,L)$ are uniquely determined;
\item $\Omega(R,L)=R[X]$;
\item $E=R$.
\end{enumerate}
When $E\ne R$ and $E\subset T$, the second tensor map fails
injectivity. When $E\not\subset T$, multiplicative tensors fail
to exist for all elements. In particular, for $L=K$ the answer
is exactly $\operatorname{depth}R\ge2$.
\end{saresult}
\begin{proof}
B.034 identifies the module with $\Omega(R,E)$ and proves
$2\Longleftrightarrow3$; ordinary polynomials prove
$2\Rightarrow1$. Suppose first $R\subsetneq E\subset T$.
As $R$ is Noetherian and $T$ is finite, $E$ is a finite
$R$-module. Choose finitely many module generators in $R^\sharp$;
a common power $\mathfrak m^N$ sends all of them into $R$.
Thus the conductor of $R\subset E$ contains $\mathfrak m^N$
and has finite quotient. B.030 makes its second tensor map
noninjective.

It remains to treat $E\not\subset T$. The finite normalization
$T$ is a Noetherian normal semilocal domain. Its maximal ideals
are finite in number and all lie over $\mathfrak m$.
Let $\mathcal D$ be the set of its maximal ideals of height one.
For $\mathfrak p\in\mathcal D$ the ring $T_{\mathfrak p}$ is a
discrete valuation ring. Choose $\pi_{\mathfrak p}$ congruent to
a valuation-one element modulo $\mathfrak p^2$, and congruent to
$1$ at every other maximal ideal. Localizing shows
\[
 \mathfrak p=\pi_{\mathfrak p}T;
\]
this uses a principal ideal only at these height-one maximal
ideals, not at the other branches.

Take $b\in E\setminus T$. Since a Noetherian normal domain is
the intersection of its height-one valuation rings, some
height-one prime has negative valuation on $b$. Any such prime
must be maximal. Indeed, for a nonmaximal height-one prime
$\mathfrak q$ of $T$, its contraction in $R$ differs from
$\mathfrak m$ by integrality, and B.034 gives
$b\in R_{\mathfrak q\cap R}\subset T_{\mathfrak q}$.
Therefore all negative valuations of $b$ are among
$\mathcal D$. Put
\[
 c=\prod_{\substack{\mathfrak p\in\mathcal D\\v_{\mathfrak p}(b)<0}}
          \pi_{\mathfrak p}^{-v_{\mathfrak p}(b)},\qquad d=bc.
\]
Then $c,d\in T$ by the height-one intersection formula.
At a maximal ideal occurring in the product, $d$ is a unit;
at every other maximal ideal, $c$ is a unit. Thus $(c,d)=T$,
and the Bezout argument of B.031 proves
\[
 \pi_{\mathfrak p}^{-1}\in T[b]
       \quad\text{whenever }v_{\mathfrak p}(b)<0.
\]
Fix one such branch $\mathfrak p_0$ and abbreviate its uniformizer
as $\pi$. There is another maximal ideal $\mathfrak p_1$ of
$T$: otherwise $T$ would be a one-dimensional local domain,
contrary to $\dim T=\dim R\ge2$ for an integral extension.
No dimension assumption on $T_{\mathfrak p_1}$ is made.

The nonzero conductor $I=(R:T)$ is contained in $\mathfrak m$,
since otherwise $T=R$ would be local. Choose any nonzero $a\in I$
and write $c_0=v_0(a)>0$, $r=c_0+1$, and $N>r c_0$.
The digit construction of B.024 supplies
\[
 H\in T[\pi^{-1}][X]
\]
whose residue at $\mathfrak p_0$ is $1$ on the class of $a^r$
modulo $\mathfrak p_0^N$, and $0$ on the zero class.
It is integer-valued at $T_{\mathfrak p_0}$; at every other
maximal ideal its coefficients are integral because $\pi$ is
a unit. Thus $H(R)\subset T$. Set
\[
 g=a(H-H(0)).
\]
Its values lie in $I\subset R$, and its coefficients lie in
$aT[b]\subset R[b]\subset L$. Hence $g\in\Omega(R,L)$.

We prove the finite-difference condition needed for the
obstruction, allowing the second branch to have higher dimension.
Given any finite family $\mathcal H\subset\operatorname{Int}(R)$,
choose a common nonzero denominator $e\in R$ for all its
polynomial coefficients. Put
\[
 e'=e/\pi^{v_0(e)}\in T.
\]
The membership follows from the height-one intersection formula,
and $e'\notin\mathfrak p_0$.
The ideals $e'\mathfrak p_1$ and $\mathfrak p_0^N$ are comaximal,
so there is $x\in e'\mathfrak p_1$ with
$x\equiv1\pmod{\mathfrak p_0^N}$. Let $u=ax\in I\subset R$.
Then $u\equiv a\pmod{\mathfrak p_0^N}$, while
\[
 u/e\in\mathfrak p_1T_{\mathfrak p_1}.
\]
The latter assertion uses that $\pi$ is a unit at
$\mathfrak p_1$ and $x/e'\in\mathfrak p_1$.
Each nonconstant term of $h(u)-h(0)$ is therefore in
$\mathfrak p_1T_{\mathfrak p_1}$, and its total value lies in
$R$. Consequently
\[
 h(u)-h(0)\in\mathfrak p_1\cap R=\mathfrak m
        \quad(h\in\mathcal H).
\]
At the chosen valuation branch,
$v_0(g(u^r))=c_0<r$.
If $g(X_1\cdots X_r)$ had a finite tensor expansion, apply this
construction to all its factors. Its iterated difference on
$\{0,u\}^r$ would lie in $\mathfrak m^r$ and hence have
valuation at least $r$, a contradiction. Repeated binary
expansion again shows that multiplicative tensors cannot exist
for every element of $\Omega(R,L)$.
This excludes the final case and proves the theorem.
\end{proof}

\begin{saresult}{B.036}{Proved}{Complete Noetherian finite-normalization classification in every dimension}
Let $R$ be a nonfield Noetherian domain whose normalization $T$
in its fraction field $K$ is finite. Let $R\subset S$ be any
$R$-torsion-free unital extension and put $L=S_{\mathrm{rat}}$.
For a maximal ideal $\mathfrak m$, write $R_m=R_{\mathfrak m}$,
$L_m=L_{\mathfrak m}$ and $T_m=T_{R\setminus\mathfrak m}$.
The polynomial identities of 153/004 give unique coaddition and
comultiplication tensors if and only if the following holds
at every maximal ideal with finite residue field:
\begin{enumerate}
\item if $\dim R_m=1$, then $L_m=R_m$, or $L_m=K$ and $T_m$
is local;
\item if $\dim R_m\ge2$, then
\[
 L_m\cap R_m^\sharp=R_m,
 \qquad
 R_m^\sharp=\{x\in K:(\mathfrak mR_m)^Nx\subset R_m
                         \text{ for some }N\ge0\}.
\]
\end{enumerate}
In every positive case all multivariable tensor maps are
isomorphisms and the unique compatible plethory exists.
For fraction-field coefficients $S=K$, this criterion is:
\begin{enumerate}
\item a single normalization branch at each finite-residue
maximal ideal of height one;
\item depth at least two at every finite-residue maximal ideal
of height at least two.
\end{enumerate}
No hypothesis or restriction is imposed at maximal ideals with
infinite residue field.
\end{saresult}
\begin{proof}
Use B.014, in all variables, to replace $S$ by $L$.
B.033 identifies every localized polynomial module and tensor
map. At infinite-residue maximal ideals use the Vandermonde
argument. At finite-residue maximal ideals of height one use
B.031; at those of larger height use B.035.
These local results prove necessity because injectivity and
existence of the substituted tensors localize, as in B.032.
Conversely their positive cases give all local tensor
isomorphisms. The global kernels and cokernels vanish by the
localization test in B.028, yielding all global isomorphisms.
B.009 gives the unique plethory.
Finally B.034 identifies $R_m^\sharp=R_m$ with depth at least
two. This proves the stated fraction-field specialization.
\end{proof}


\begin{saresult}{B.037}{Proved}{An explicit two-dimensional tensor kernel}
Put
\[
 T=\mathbb F_2[x,y]_{(x,y)},\qquad
 I=(x,y)^2T,\qquad R=\mathbb F_2+I,\qquad K=\mathbb F_2(x,y).
\]
Then $R$ is a two-dimensional Noetherian local domain with residue
field $\mathbb F_2$ and finite normalization $T$. With $q(X)=X^2-X$,
\[
 \operatorname{Int}(R)=R\oplus RX\oplus q(X)T[X].
\]
The nonzero tensor
\[
 z=(xq)\otimes q-q\otimes(xq)
\]
has zero polynomial image. Thus even fraction-field coefficients
do not give uniquely determined polynomial tensors in this
explicit higher-dimensional case, essentially of finite type over
$\mathbb F_2$.
\end{saresult}
\begin{proof}
Let $A=\mathbb F_2[x^2,y^2]_{(x^2,y^2)}$. Then
$T=A\oplus Ax\oplus Ay\oplus Axy$: in characteristic two the
square of every denominator with nonzero constant term belongs
to $A^*$, so this equality holds after the stated localizations.
The ring $R$ is generated as an $A$-module by
\[
 1,\ xy,\ x^3,\ xy^2,\ x^2y,\ y^3.
\]
Indeed an $Ax$ or $Ay$ coefficient must have zero constant term,
and the maximal ideal of $A$ is $(x^2,y^2)A$.
Thus $R$ is a finite $A$-module, and every ideal of $R$ is
finitely generated already as an $A$-module; this proves
Noetherianity. Also $T$ is finite over $R$, and both have
fraction field $K$, since $x=x^3/x^2$ and $y=y^3/y^2$.
Dimension is two by the integral dimension theorem.
The quotient $R/I$ is $\mathbb F_2$, and any element of $R$
outside $I$ has inverse in $1+I$ up to its constant unit;
thus $R$ is local. The normal ring $T$, a localization of a
polynomial unique factorization domain, is its normalization.

Every height-one prime of $T$ is nonmaximal and contracts to a
nonmaximal prime of $R$. Vandermonde interpolation at distinct
residues of $R$ puts every coefficient of an integer-valued
polynomial in the corresponding localization of $T$.
Intersecting these valuation rings gives coefficients in $T$.
Thus $\operatorname{Int}(R)=\Omega(R,T)$.
For $r=a+i$ with $a\in\mathbb F_2$ and $i\in I$, one has
$q(r)=i^2-i\in I$, so $qT[X]\subset\operatorname{Int}(R)$.
Monic division by $q$ leaves a linear remainder; evaluation at
$0,1$ forces both its coefficients into $R$. This proves the
displayed direct sum.

The summand $Tq$ is split over $R$, so $Tq\otimes_RTq$ injects
as a summand of the full tensor module. After base change to
$R/I=\mathbb F_2$ it becomes
\[
 D\otimes_{\mathbb F_2}D,
 \qquad D=\mathbb F_2[x,y]/(x,y)^2.
\]
The image of $z$ is $x\otimes1-1\otimes x$, a nonzero difference
of two of its nine basis vectors. Its polynomial image is zero
because both terms equal $xq(X)q(Y)$.
The attached exact program checks all $512$ elements and
$262144$ multiplication pairs of this finite tensor algebra;
the diagonal map has image of size $8$ and kernel of size $64$.
It also checks the null-polynomial/conductor identity on all
$256$ possible $R$-inputs modulo $(x,y)^4$ and all $1024$
coefficient classes in that quotient.
\end{proof}


\begin{saresult}{B.038}{Proved}{A derivation obstruction to tensor injectivity}
Let $R$ be a domain with fraction field $K$, let $V\subset K$ be
a discrete valuation overring of $R$, and let $D:K\to K$ be a
derivation satisfying $D(R)\subset V$. Suppose there are
$u_n\in R$ with $v(u_n)\to\infty$ and $D(u_n)=1$.
For $A=\operatorname{Int}(R)$ the map
\[
 \lambda:A\longrightarrow V,\qquad F\longmapsto F'(0)
\]
is well-defined. There is an $R$-balanced bilinear map
\[
 \beta:A\times A\longrightarrow K/V,
 \qquad \beta(F,G)=D(F'(0))G'(0)+V.
\]
If $P\in R[X]$, $c\in K$, $cP\in A$ and
$D(c)P'(0)^2\notin V$, then
\[
 (cP)\otimes P-P\otimes(cP)
\]
is a nonzero tensor with zero polynomial image.
\end{saresult}
\begin{proof}
Apply $D$ to $F(u_n)\in R$. The polynomial chain rule gives
\[
 D(F(u_n))=(DF)(u_n)+F'(u_n),
\]
where $DF$ means applying $D$ to the coefficients.
Its constant coefficient is $D(F(0))\in V$.
All nonconstant terms of $(DF)(u_n)$ have nonnegative valuation
for sufficiently large $n$, since their coefficients are fixed
and $v(u_n)\to\infty$. Hence $(DF)(u_n)\in V$ eventually.
The left side lies in $V$, so $F'(u_n)\in V$.
The nonconstant terms of $F'(u_n)-F'(0)$ likewise have
nonnegative valuation eventually. Therefore $F'(0)\in V$.

Additivity of $\beta$ is immediate. For $r\in R$,
\[
 D(rF'(0))G'(0)-rD(F'(0))G'(0)
   =D(r)F'(0)G'(0)\in V.
\]
Thus $\beta(rF,G)=r\beta(F,G)=\beta(F,rG)$, so it induces
an $R$-linear functional on $A\otimes_RA$ with target $K/V$.
On the displayed tensor its value is
\[
 \bigl(D(cP'(0))P'(0)-D(P'(0))cP'(0)\bigr)+V
   =D(c)P'(0)^2+V,
\]
which is nonzero by assumption. Its polynomial image is zero
by cancellation of the two equal products. This proves the
obstruction without any flatness assumption.
\end{proof}

\begin{saresult}{B.039}{Proved}{An explicit nonfinite-normalization counterexample}
Let $k=\mathbb F_2$, $A_0=k[t]_{(t)}$, $n_0=0$ and
$n_i=(i+1)!$ for $i\ge1$. In $k[[t]]$ set
\[
 z=\sum_{i\ge0}t^{n_i},\qquad
 z_i=\frac{z-\sum_{j<i}t^{n_j}}{t^{n_i}},\qquad
 u=t(z-1),\qquad s_i=(z_i-1)^2.
\]
Define
\[
 C=A_0[u,s_0,s_1,\ldots],\qquad
 M=(t,u)C,\qquad R=C_M.
\]
Then $R$ is a one-dimensional Noetherian local domain with
residue field $k$, fraction field $K=k(t,z)$ and normalization
a discrete valuation ring $V$ with parameter $t$ and residue
field $k$. Its normalization is not a finite $R$-module.
Nevertheless, unlike the finite-conductor one-branch cases,
$\operatorname{Int}(R)$ has a noninjective tensor map.
Explicitly, with $q(X)=X^2-X$ and $c=z-1$,
\[
 cq\in\operatorname{Int}(R),\qquad
 (cq)\otimes q-q\otimes(cq)\ne0,
\]
while its polynomial image is zero.
\end{saresult}
\begin{proof}
This is an explicit characteristic-two specialization of the
Akizuki--Reid construction. The known ring-theoretic conclusions
are stated in Reid, \emph{Akizuki's counterexample}, the two
theorems on page 2 and their proofs on pages 2--4; the opened
author LaTeX source is lines 49--212. We verify the hypotheses
and give the short characteristic-two arguments needed below,
including a direct nonfiniteness proof.

First $z$ is transcendental over $k(t)$. Otherwise let
$0\ne P(t,Z)\in k[t,Z]$ vanish at $z$, with $Z$-degree $d$
and coefficient $t$-degrees at most $h$. Its truncation
$z^{(N)}=\sum_{i\le N}t^{n_i}$ gives
\[
 v_t(P(t,z^{(N)}))\ge n_{N+1},\qquad
 \deg_t P(t,z^{(N)})\le h+d n_N.
\]
For all sufficiently large $N$, $n_{N+1}>h+dn_N$, so the
polynomial $P(t,z^{(N)})$ is zero. The distinct elements
$z^{(N)}\in k(t)$ would then give infinitely many roots of the
nonzero polynomial $P(t,Z)$ over that field, a contradiction.
The gap condition $n_i\ge2n_{i-1}+2$ of Reid's theorem also
holds, including $i=1,2$.

Here are direct checks of the local-ring assertions in this
specialization. Put $B=A_0[z_0,z_1,\ldots]$. The equations
\[
 z_i-1=t^{n_{i+1}-n_i}z_{i+1}
\]
make $tB$ the kernel of the residue map $B\to k$.
The local domain $V=B_{tB}$ has principal maximal ideal $tV$
and $\bigcap_n t^nV=0$, because it embeds in $k[[t]]$.
Every nonzero element therefore equals a power of $t$ times
a unit. Every nonzero ideal has an element of least such power,
so is principal; this proves that $V$ is a discrete valuation
ring, without presupposing Noetherianity of $B$.

In characteristic two,
\[
 s_i=t^{2(n_{i+1}-n_i)}(1+s_{i+1}),\qquad u^2=t^2s_0.
\]
Thus every $s_i$ lies in $tC$, and $C/M=k$.
Iterating the first identity shows that, modulo $t^NC$,
each of finitely many specified $s_i$ is represented by an
element of $A_0$. Using the second identity to reduce powers
of $u$ shows that $C/t^NC$ is generated over $A_0/(t^N)$ by
$1,u$. More precisely, with
$u_r=t(z-\sum_{j\le r}t^{n_j})=t^{n_r+1}(z_r-1)$,
every $f\in C$ can be written
\[
 f=X+Yu_r+t^NZ,\qquad X,Y\in A_0,\quad Z\in C,
\]
for any prescribed $r,N$, since $u-u_r\in A_0$.
The radical of $tC$ is $M$, as $C/tC$ is generated by the
square-zero image of $u$. Hence localization at $M$ does not
change $C/t^NC$, which is a finite $A_0/(t^N)$-module.

For $0\ne f\in M$, let $v_t(f)=a$, choose $N>a$ and then
$r$ with $n_r\ge N$. In the preceding expression $X$ has
valuation $a$. Divide by its unit factor in $A_0$ to write
\[
 f=t^a(1+t^{N-a}Z)+Yu_r.
\]
Squaring in characteristic two and using
$u_r^2=t^{2n_r+2}s_r$ gives
$f^2=t^{2a}w$ with $w\in C\setminus M$.
Thus $fR$ contains $t^{2a}$. The same conclusion holds for
any nonzero element of $R$ after clearing its unit denominator.
Every nonzero ideal consequently contains a power of $t$;
its quotient by that principal ideal is a submodule of a
finite $A_0/(t^N)$-module, and is finitely generated. This proves
that $R$ is Noetherian. Its only nonzero prime is $MR$, so its
dimension is one and its residue field is $k$.

Both $R$ and $V$ have fraction field $K$, since $z=1+u/t$.
Every $z_i$ is integral over $C$, as $z_i^2=1+s_i\in C$.
After first localizing $B$ by $C\setminus M$, all its maximal
ideals lie over $MR$ and contain $t$, so the only one is the
extension of $tB$. This localized ring is therefore $B_{tB}=V$
and is integral over $R$. As $V$ is integrally closed in $K$,
it is the normalization.

The derivation $D=t^{-1}\partial/\partial z$ of $K$ preserves
$R$: it kills $A_0$ and all the squares $s_i$, sends $u$ to $1$,
and extends to the localization by the quotient rule.
If $V$ were a finite $R$-module with generators $v_1,\ldots,v_h$,
then the valuations of $D(V)$ would be uniformly bounded below
by the minimum of $0,v(D(v_1)),\ldots,v(D(v_h))$.
Indeed $D(rv_j)=D(r)v_j+rD(v_j)$ and $D(r),r\in R\subset V$.
But $z_i\in V$ and
\[
 D(z_i)=t^{-n_i-1},
\]
whose valuations are unbounded below. This proves nonfiniteness
directly.

The elements $u_r\in R$ satisfy $D(u_r)=1$ and
$v_t(u_r)=n_{r+1}+1\to\infty$, so B.038 applies.
Also $cMR\subset R$, since $ct=u$ and $cu=t s_0$ both lie in
$R$. As $q(R)\subset MR$ by the residue field $k=\mathbb F_2$,
we have $cq\in\operatorname{Int}(R)$.
Finally $D(c)=1/t\notin V$ and $q'(0)=-1$.
The functional of B.038 takes the displayed tensor to
$1/t+V\ne0$, whereas its polynomial image is zero.
This proves the asserted counterexample with all its hypotheses.
\end{proof}


\begin{saresult}{B.040}{Proved}{Every finite-residue Akizuki--Reid family fails tensor uniqueness}
Let $A_0$ be any discrete valuation ring with uniformizer $t$ and
finite residue field of cardinality $q$. Choose units $a_i\in A_0$
and exponents $n_0=0$, $n_{i+1}\ge2n_i+2$, such that
$z=\sum_i a_it^{n_i}\in\widehat A_0$ is transcendental over
$\operatorname{Frac}(A_0)$. Define $z_i,C,M,R$ by Reid's construction:
\[
 z_i=\frac{z-\sum_{j<i}a_jt^{n_j}}{t^{n_i}},\qquad
 C=A_0[t(z-a_0),(z_i-a_i)^2:i\ge0].
\]
Put
\[
 M=(t,t(z-a_0))C,\qquad R=C_M.
\]
Then $R$ is a one-dimensional Noetherian local domain with
nonfinite discrete valuation normalization $V$, as in Reid's
opened theorem. With $c=z-a_0$ and $P=X^q-X$, one has
$cP\in\operatorname{Int}(R)$ and
\[
 (cP)\otimes P-P\otimes(cP)\ne0,
\]
although its polynomial image is zero. This includes mixed
characteristic; it is not restricted to characteristic two.
\end{saresult}
\begin{proof}
The ring assertions, with exactly the hypotheses just stated,
are Reid's two theorems on page 2, with proofs on pages 2--4.
Set $D=t^{-1}\partial/\partial z$ on
$K=\operatorname{Frac}(A_0)(z)$. This derivation kills $A_0$
and sends $t(z-a_0)$ to $1$. On the remaining generators,
\[
 D((z_i-a_i)^2)=2t^{-n_i-1}(z_i-a_i).
\]
Its valuation is at least
$n_{i+1}-2n_i-1\ge1$; if $2=0$ the value is zero.
Thus $D(C)\subset V$. Denominators from $C\setminus M$ are
units of $V$, so the quotient rule gives $D(R)\subset V$.
The elements
\[
 u_i=t\left(z-\sum_{j\le i}a_jt^{n_j}\right)
      =t^{n_i+1}(z_i-a_i)\in R
\]
satisfy $D(u_i)=1$ and $v_t(u_i)=n_{i+1}+1\to\infty$.
Consequently B.038 applies with its weaker, explicitly stated
hypothesis $D(R)\subset V$.

Writing $u=t(z-a_0)$, one has $ct=u$ and
$cu=t(z-a_0)^2\in C$, so $cMR\subset R$.
The polynomial $P$ has values in $MR$ because the residue field
has $q$ elements. Hence $cP\in\operatorname{Int}(R)$.
Finally $D(c)=1/t$ and $P'(0)=-1$; B.038 takes the tensor to
$1/t+V\ne0$ and proves the assertion.
The same bounded-derivative argument as in B.039 reproves
nonfiniteness of $V/R$, since $D(z_i)=t^{-n_i-1}$ has
unboundedly negative valuation.
\end{proof}


\begin{saresult}{B.041}{Proved}{Every infinite valuation domain and every torsion-free extension}
Let $V$ be an infinite valuation domain with fraction field $K$,
and let $V\subset S$ be any $V$-torsion-free unital extension.
Then every polynomial tensor map for $\Omega(V,S)$ is an
isomorphism, and the unique compatible plethory exists.
More precisely, put $L=S_{\mathrm{rat}}$. If the residue field
of $V$ is infinite or its maximal ideal is nonprincipal, then
$\Omega(V,S)=V[X]$. If its nonzero maximal ideal is principal,
with finite residue field, then
\[
 \Omega(V,S)=
 \begin{cases}
 V[X],&L=V,\\
 \operatorname{Int}(V),&L\ne V,
 \end{cases}
\]
and $\operatorname{Int}(V)$ has an explicit free Newton basis.
This covers valuation domains of arbitrary rank and does not
assume Noetherianity. A finite valuation domain is a finite field
and is governed instead by B.013.
\end{saresult}
\begin{proof}
The field case is B.014, so suppose $V$ is not a field.
B.014 reduces $S$ to the overring $L\subset K$.
An infinite residue field gives $\operatorname{Int}(V)=V[X]$
by the invertible Vandermonde argument; hence it gives the
asserted equality for every coefficient overring.

Suppose the maximal ideal is nonprincipal. The ordered value
group $\Gamma$ has no least positive element: a least positive
value would give a generator of that ideal, and conversely.
For every $\gamma>0$ and every positive integer $N$, there is
$\delta>0$ with $N\delta<\gamma$. To verify this without an
Archimedean assumption, choose $0<\epsilon<\gamma$; any positive
$\epsilon$ has a positive $\eta$ with $2\eta\le\epsilon$.
Indeed choose $0<\xi<\epsilon$ and take $\eta=\xi$ if
$2\xi\le\epsilon$, otherwise take $\eta=\epsilon-\xi$.
Iterate until $2^j\ge N$.

Take $F\in\operatorname{Int}(V)$ of degree $d\ge1$. If one
coefficient had value $-\gamma<0$, choose $a\in V$ with value
$\delta>0$ so small that
\[
 \frac{d(d+1)(d+2)}6\,\delta<\gamma.
\]
The $d+1$ points $0,a,a^2,\ldots,a^d$ are distinct. Their
Vandermonde determinant has value exactly
$d(d+1)(d+2)\delta/6$: the difference of distinct positive
powers of $a$ has the smaller of their two values.
Cramer's rule, with all values $F(a^i)$ in $V$, forces every
coefficient to have value at least the negative of this
quantity, contradicting $-\gamma$. Constant polynomials are
already in $V$. Thus again $\operatorname{Int}(V)=V[X]$.

It remains to suppose the maximal ideal is $\pi V$ and the
residue field has $q$ elements. Choose representatives
$c_0=0,c_1,\ldots,c_{q-1}\in V$. For the base-$q$ expansion
$n=\sum_j n_jq^j$ put $a_n=\sum_jc_{n_j}\pi^j$ and
\[
 w_q(n)=\sum_{j\ge1}\lfloor n/q^j\rfloor,\qquad
 b_n(X)=\frac{\prod_{i<n}(X-a_i)}{\prod_{i<n}(a_n-a_i)}.
\]
The denominator has value exactly $w_q(n)v(\pi)$, by the
first differing digit of each pair. Every $x\in V$ has any
prescribed finite number of successive residue digits, since
its maximal ideal is $\pi V$. Among $a_0,\ldots,a_{n-1}$ at
least $\lfloor n/q^j\rfloor$ share the first $j$ digits of
$x$. Summing these finite lower bounds shows that the
numerator has value at least $w_q(n)v(\pi)$.
This argument uses only finitely many digit levels with
nonzero lower bounds, so it remains valid in a higher-rank
value group. Hence $b_n\in\operatorname{Int}(V)$ and
$b_n(a_n)=1$. Successive evaluation at these nodes gives
\[
 \operatorname{Int}(V)=\bigoplus_{n\ge0}Vb_n.
\]
Every $b_n$ has coefficients in $V[\pi^{-1}]$.

If $L\ne V$, choose $b\in L\setminus V$.
The least positive value $v(\pi)$ satisfies
$-v(b)\ge v(\pi)$. Therefore $(\pi b)^{-1}\in V$ and
$\pi^{-1}=(\pi b)^{-1}b\in L$.
The entire Newton basis belongs to $L[X]$, so
$\Omega(V,L)=\operatorname{Int}(V)$. If $L=V$ the module is
$V[X]$. These modules are free over $V$ in every case;
B.018 gives all tensor isomorphisms and the plethory.
\end{proof}


\begin{saresult}{B.042}{Proved}{Every infinite Krull domain with arbitrary torsion-free coefficients}
Let $R$ be an infinite Krull domain with fraction field $K$, and
let $R\subset S$ be an $R$-torsion-free unital extension.
Then $\Omega(R,S)$ is a projective $R$-module, every
multivariable tensor map is an isomorphism, and the unique
compatible plethory exists. No Noetherian hypothesis is needed.
Here the Krull hypotheses used are precisely
$R=\bigcap_{\operatorname{ht}\mathfrak p=1}R_{\mathfrak p}$,
each $R_{\mathfrak p}$ a discrete valuation ring, and every
nonzero element of $R$ lying in only finitely many such primes.
\end{saresult}
\begin{proof}
The field case is B.014. Put $L=S_{\mathrm{rat}}\subset K$.
We verify the two places where B.020 used Noetherianity and
then complete the fractional-ideal decomposition directly.

Every height-one maximal ideal $\mathfrak p$ is invertible.
Choose $b\in R$ with $v_{\mathfrak p}(b)=1$. Besides
$\mathfrak p$, only finitely many height-one primes
$\mathfrak q_i$ contain $b$. Choose
$a_i\in\mathfrak q_i^{v_{\mathfrak q_i}(b)}\setminus\mathfrak p$
and put $a=\prod_i a_i$. Such choices exist because distinct
height-one primes are incomparable; powers of an element of
$\mathfrak q_i\setminus\mathfrak p$ suffice.
Then $x=a/b$ has valuation $-1$ at $\mathfrak p$ and
nonnegative valuation at every other height-one prime.
Consequently $x\mathfrak p\subset R$, by the Krull intersection
property. The ideal $\mathfrak p\mathfrak p^{-1}$ contains
$\mathfrak p$ and $xb=a\notin\mathfrak p$, so equals $R$.
This proves invertibility rather than assuming finite generation.

For each fixed prime power $q$, choose $r_q\in R$ with
$r_q^q-r_q\ne0$; this is possible since $R$ is infinite.
Every height-one maximal ideal with residue field of size $q$
contains this nonzero element, so there are only finitely many
such ideals. In particular the following product is finite:
\[
 I_n=\prod_{\substack{\operatorname{ht}\mathfrak p=1,\,
                       \mathfrak p\text{ maximal}\\
                       |R/\mathfrak p|=q<\infty}}
           \mathfrak p^{-w_q(n)},\qquad
 w_q(n)=\sum_{j\ge1}\lfloor n/q^j\rfloor.
\]
All its factors, and hence $I_n$, are invertible fractional ideals.
An invertible fractional ideal is the intersection of its
height-one localizations: if $x$ belongs to those localizations,
then $xI_n^{-1}$ is integral at every height-one prime and hence
lies in $R$, which implies $x\in I_n$.

The digit Newton argument at a finite-residue height-one maximal
ideal shows that the leading coefficient of an integer-valued
polynomial of degree at most $n$ has valuation at least
$-w_q(n)$. The necessary nodes can be chosen in $R$ by successive
residue digits and a valuation-one element in $R$.
At every other height-one prime, the image of $R$ in its residue
field is infinite, so Vandermonde interpolation gives
nonnegative valuations of all coefficients. The leading
coefficient ideal is therefore contained in $I_n$.

For each prime occurring in the finite product, choose its monic
local digit polynomial $Q_{n,\mathfrak p}\in R[X]$.
The Chinese remainder theorem for the pairwise comaximal
powers $\mathfrak p^{w_q(n)}$ gives a monic polynomial
$P_n\in R[X]$ of degree $n$, congruent coefficientwise to
$Q_{n,\mathfrak p}$ modulo that power at each such prime.
For $c\in I_n$ and $r\in R$, the product $cP_n(r)$ is integral
at every height-one prime, and so belongs to $R$.
Subtracting its leading term successively proves
\[
 \operatorname{Int}(R)=\bigoplus_{n\ge0}I_nP_n,\qquad
 \Omega(R,L)=\bigoplus_{n\ge0}(I_n\cap L)P_n.
\]
The second equality uses that the $P_n$ are monic with
coefficients in $R$, so descending elimination preserves $L[X]$.

Put $J_n=I_n\cap L$ and let $E_n$ be the finite set of primes
occurring in $I_n$. Localization commutes with this finite
intersection of submodules of $K$. At a maximal ideal outside
$E_n$, $J_n$ localizes to $R_{\mathfrak m}$. At
$\mathfrak p\in E_n$, the overring $L_{\mathfrak p}$ of the
rank-one discrete valuation ring $R_{\mathfrak p}$ is either
$R_{\mathfrak p}$ or $K$. Thus localization at every maximal
ideal proves the exact formula
\[
 J_n=\prod_{\substack{\mathfrak p\in E_n\\L_{\mathfrak p}=K}}
                 \mathfrak p^{-w_{|R/\mathfrak p|}(n)}.
\]
No finite-generation assertion about an arbitrary intersection
has been used: equality is proved by localization against the
displayed finite invertible product. Each $J_n$ is therefore
projective, and their direct sum is projective. B.019 proves
all tensor and plethory assertions.
\end{proof}

\begin{saresult}{B.043}{Proved}{Infinite-residue conductor cusps and all their coefficients}
Let $k=\mathbb F_q$ be finite, let $l\supset k$ be any infinite
field, and put
\[
 V=l[[t]],\qquad R=k+t^aV\quad(a\ge1),\qquad K=l((t)).
\]
For every $R$-torsion-free unital extension $R\subset S$, put
$L=S_{\mathrm{rat}}$. Then the polynomial coaddition and
comultiplication tensors are uniquely determined if and only
if $L=R$. In particular fraction-field coefficients fail.
For $A=\operatorname{Int}(R)$, $s\in l\setminus k$,
$Q(X)=X^q-X$ and $F(X)=Q(X)^2/t^a$, an explicit kernel is
\[
 (sF)\otimes F-F\otimes(sF)\ne0.
\]
The ring $R$ is a non-Noetherian local domain with maximal ideal
$t^aV$, residue field $k$, and fraction field $K$.
\end{saresult}
\begin{proof}
The ring assertions except non-Noetherianity follow from its
pullback description and from $t^aV\subset R$: an element with
nonzero constant in $k$ has inverse in $k+t^aV$, and division
by a power of $t^a$ gives every element of $K$.
Its maximal ideal $\mathfrak m=t^aV$ has
$\mathfrak m/\mathfrak m^2$ of infinite dimension over $k$,
since it contains $t^al$ and $[l:k]=\infty$.
Thus $\mathfrak m$ is not finitely generated.

We first treat coefficient field $K$ and construct a functional
which proves that the proposed tensor is nonzero.
For $G\in A$, the polynomial $G(t^aY)$ takes $V$ into
$R\subset V$. Since $V$ has infinite residue field,
$\operatorname{Int}(V)=V[Y]$ by Vandermonde interpolation.
Its reduction modulo $t$ is a polynomial over $l$ taking all
of $l$ into $k$. Such a polynomial is constant: the product of
its differences from all elements of $k$ vanishes at every
point of the infinite field $l$, hence is the zero polynomial;
one factor must be zero. Its constant is $G(0)\bmod t$.
It follows that
\[
 G(t^ay)-G(0)\in R\cap tV=t^aV\quad(y\in V).
\]
Apply the same Vandermonde argument to its quotient by $t^a$.
This defines an $R$-linear map, with $R$ acting through $k$,
\[
 h:A\longrightarrow l[Y],\qquad
 h(G)=\left(\frac{G(t^aY)-G(0)}{t^a}\right)\bmod t.
\]
In particular $h$ kills $\mathfrak m A$.

For $r\in R$ the value $Q(r)$ lies in $t^aV$, so both $F$ and
$sF$ take $R$ into $t^aV\subset R$. Direct expansion gives
\[
 h(F)=Y^2,\qquad h(sF)=sY^2.
\]
The image of the proposed tensor under $h\otimes h$ is
\[
 sY^2\otimes Z^2-Y^2\otimes sZ^2
       \in l[Y]\otimes_k l[Z],
\]
which is nonzero. For example choose a $k$-linear functional
$\ell:l\to k$ with $\ell(1)=0$, $\ell(s)=1$, and apply it
to the first $Y^2$ coefficient and take the second $Z^2$
coefficient. The result is $1$. The polynomial image of the
original tensor is zero by cancellation of its two products.

For completeness consider all coefficient overrings. Any
$L\not\subset V$ is $K$: if $v(b)<0$, choose $m$ with $v(b^{-m})\ge a$; then
$c=b^{-m}\in t^aV\subset R$ and $c^{-1}\in L$. Every element
of $K$ is sent into $R$ by a high power of $c$.
Suppose now $R\subsetneq L\subset V$. The conductor contains
$\mathfrak m$, and $Q(R)\subset\mathfrak m$.
Monic division gives a split summand $LQ$ in $\Omega(R,L)$.
If the second tensor map were injective, multiplication
$L\otimes_RL\to L$ would be an isomorphism. Modulo
$\mathfrak m$, it would give injective multiplication on the
nonzero $k$-algebra $C=L/\mathfrak mL=L/\mathfrak m$.
Choose a $k$-linear functional taking $1$ to $1$.
The relation $c\otimes1=1\otimes c$, forced by injectivity,
then implies that every $c\in C$ belongs to $k$.
Hence $C=k$ and $L=k+\mathfrak m=R$, a contradiction.
This argument does not use a finite-module Nakayama lemma:
$L$ can be infinite over $R$.
The only positive case is therefore $L=R$, which gives the
ordinary polynomial plethory. B.014 transfers the result to
arbitrary torsion-free $S$.
\end{proof}


\begin{saresult}{B.044}{Proved}{Integer-valued polynomials commute with one-dimensional finite-residue completion}
Let $(R,\mathfrak m)$ be a one-dimensional Noetherian local
domain with finite residue field, let $K=\operatorname{Frac}(R)$,
and put $A=\widehat R$ and $Q=K\otimes_R A$.
For every number of variables $n$ there is a canonical isomorphism
\[
 \operatorname{Int}(R^n,R)\otimes_R A
       \simeq \operatorname{Int}(A^n,A)\subset Q[X_1,\ldots,X_n].
\]
It identifies the tensor maps after faithfully flat base change.
In particular a noninjective second tensor map for $A\subset Q$
implies a noninjective second tensor map for $R\subset K$.
\end{saresult}
\begin{proof}
Fix bounds on the degree in each variable and let $N$ be the
number of allowed monomials. Choose a rectangular grid of
distinct points in $R$. The product of its Vandermonde
determinants gives a nonzero $d\in R$ such that every
integer-valued polynomial within these bounds has all
coefficients in $d^{-1}R$. The same grid over $A$ shows that
all coefficients there lie in $d^{-1}A$: flatness makes $d$
a nonzerodivisor of $A$.

The ring $R/dR$ is a finite ring, because it is Artinian local
with finite residue field. Choose a finite set $E\subset R$
of representatives. For $f\in d^{-1}R[X_1,\ldots,X_n]$,
its value modulo $R$ depends only on each input modulo $dR$.
Indeed $df$ has coefficients in $R$, and changing a coordinate
by a multiple of $d$ changes its numerator by a multiple of
$d$. Consequently the integer-valued module in these bounds
is the kernel of the finite evaluation map
\[
 (d^{-1}R)^N\longrightarrow (d^{-1}R/R)^{E^n}.
\]
The natural map $R/dR\to A/dA$ is an isomorphism: the source
has finite length and hence is already complete. Thus exactly
the same finite tests describe the bounded-degree module over
$A$. Flat base change commutes with this kernel and identifies
the two modules. Taking the filtered union of all degree
bounds proves the assertion. Polynomial multiplication in
disjoint variables commutes with these identifications, so
the tensor maps correspond as stated. Faithful flatness
detects nonzero kernels.

For clarity, $Q$ is the total quotient ring of $A$. Choose
$0\ne t\in\mathfrak m$. Every $0\ne r\in R$ has
$t^M\in rR$ for some $M$, so $Q=A[1/t]$. This is a
zero-dimensional Noetherian ring: $A/tA$ is Artinian, and
completion has dimension one. It is Artinian and is a
finite product of local Artinian rings. Every nonzerodivisor
of $A$ becomes a unit in this product, so $A[1/t]$ is indeed
the total quotient ring.
The completion facts invoked here are the flatness and exactness
of finite-module completion (Stacks, Lemma 10.97.2, tag 00MB)
and its local faithful-flat consequence.
\end{proof}

\begin{saresult}{B.045}{Proved}{A generic nilpotent forces a tensor kernel after completion}
Let $(A,\mathfrak m)$ be a one-dimensional complete
Noetherian local ring with finite residue field. Assume that
$A$ is Cohen--Macaulay and that either $A$ has characteristic
$p>0$, or $A$ has characteristic zero and $p$ is a
nonzerodivisor, where $p$ is its residue characteristic.
If $A$ is not reduced, then
\[
 \operatorname{Int}(A)\otimes_A\operatorname{Int}(A)
       \longrightarrow\operatorname{Int}(A^2)
\]
is not injective. Here coefficients belong to the total
quotient ring of $A$. The proof constructs a nonzero tensor
and an explicit balanced linear functional detecting it.
\end{saresult}
\begin{proof}
We first construct the finite lattice in which the argument
will take place, including the needed compatible coefficient
field. There is a complete discrete valuation subring
$V_0\subset A$ such that $A$ is finite and torsion-free over
$V_0$, and such that some nonreduced local factor $Q_0$ of
$Q(A)$ has residue field $F$ separable over
$K_0=\operatorname{Frac}(V_0)$.

In mixed characteristic take the Cohen coefficient ring.
It is a complete discrete valuation ring with uniformizer
$p$, and $p$ is a parameter. The finite-module assertion
follows by lifting finitely many generators modulo $p$ and
successively approximating in the $p$-adic topology;
$A$ is $p$-adically complete because $pA$ is
$\mathfrak m$-primary. Separability holds in characteristic
zero.

Here is the extra parameter choice needed in characteristic
$p$. The roots of $X^{|A/\mathfrak m|}-X$ give a coefficient
field $k\subset A$ (or use the Cohen structure theorem).
Choose a minimal prime $\mathfrak p$ corresponding to a
nonreduced generic factor. Such a factor exists because a
parameter is a nonzerodivisor, so a nonzero nilpotent remains
nonzero in $Q(A)$. The complete one-dimensional domain
$C=A/\mathfrak p$ has finite normalization, a complete
discrete valuation ring $V=l[[u]]$ with $l/k$ finite.
Here finiteness follows from the Nagata property of complete
Noetherian local rings; a finite domain over a complete local
ring is local, because a finite complete algebra decomposes
as a product of complete local algebras. A nonzero conductor
then gives $u^N\in C$ for all sufficiently large $N$.
Choose such an $N$ with $p\nmid N$ and lift $u^N$ to
$t_0\in\mathfrak m$.

We can adjust this lift to a parameter without changing its
image in $C$. Choose
$a\in\mathfrak p$ outside every other minimal prime, by
prime avoidance. At each other minimal prime, the nonzero
images of $a,a^2,a^3,\ldots$ are distinct: $a$ lies in
the maximal ideal of the corresponding local domain.
Thus $t_0+a^j$ belongs to each such prime for at most one
$j$. Choosing $j$ outside this finite exceptional set gives
a parameter $t$ mapping to $u^N$. If there is no other
minimal prime, use $t=t_0$.
Set $V_0=k[[t]]\subset A$. This map is injective, since
a nonzero power series is $t^r$ times a unit and $t$ is a
nonzerodivisor. Successive approximation modulo $t^m$ shows
that $A$ is finite over $V_0$. In the chosen factor the
residue field is $F=l((u))$, a finite separable extension
of $k((u^N))$ since $p\nmid N$.

Since $F/K_0$ is finite separable, it embeds in $Q_0$ as
a coefficient field compatibly with $K_0$. To see this
directly, choose a primitive element of $F/K_0$ and lift
it to $Q_0$. Its separable minimal polynomial has invertible
derivative. Newton correction modulo successive powers of
the nilradical, which is nilpotent, gives an actual root
lifting the primitive element. Its $K_0$-algebra is the
required copy of $F$. Write
\[
 \rho:Q_0\longrightarrow F,\qquad N=\ker\rho,
 \qquad N^h=0\quad(h\ge2).
\]
Let $V\subset F$ be the integral closure of $V_0$ in $F$.
It is a finite complete discrete valuation ring, with
uniformizer $u$ and finite residue field. The same notation
$V=l[[u]]$ applies in the equal-characteristic construction.
The ring
\[
 O=V[\operatorname{pr}_0(A)]\subset Q_0
\]
is finite and torsion-free over $V$, contains $V$, satisfies
$\rho(O)=V$, and spans $Q_0$ over $F$.
Finiteness follows because $\operatorname{pr}_0(A)$ is
finite over $V_0$ and hence integral over $V$.
Let $e$ denote the idempotent of the selected factor in
$Q(A)$. Since $eO$ is a finite $V_0$-module in $Q(A)$,
clearing finitely many denominators gives an integer $c\ge0$
such that
\[
                         u^c eO\subset A.                 \tag{B45.1}
\]
Indeed a sufficiently high power of the uniformizer of
$V_0$ works, and in $V$ that power is a unit times a
power of $u$; the unit preserves $O$.

Put $J=N\cap O$. There is a primitive element
$z_0\in J\cap\operatorname{Ann}_{Q_0}(N)$: start with a
nonzero element in the last nonzero power of $N$, multiply
it into $O$, and divide by the greatest possible power of
$u$ in the finite free $V$-module $O$. Choose a
$V$-linear map $\ell:O\to V$ with $\ell(z_0)=1$ and
extend it $F$-linearly to $Q_0$. Such a map exists because
a primitive vector in a finite free module over a discrete
valuation ring extends to a basis.

There is an integer $b\ge0$, independent of
$G\in\operatorname{Int}(A)$, such that, with
$x_G=\operatorname{pr}_0(G'(0))$,
\[
                  Jx_G\subset u^{-b}O,\qquad
                  \rho(x_G)\in u^{-b}V.                   \tag{B45.2}
\]
Here is a finite interpolation proof. Choose finitely many
$V$-generators $n_1,\ldots,n_s$ of $J$. By (B45.1),
$u^c e n_j\in A$. For any $h$ fixed distinct scalars
$s_1,\ldots,s_h\in V_0$, the values
$G(s_i u^c e n_j)$ lie in $A$ and their selected components
lie in $O$. As a polynomial in the scalar $s$, this
evaluation has degree at most $h-1$, because $n_j^h=0$.
Inverting this one fixed Vandermonde matrix bounds its
linear coefficient $u^c n_jx_G$ in a fixed fractional
lattice $u^{-d}O$. Increasing $b$ gives the first inclusion
for every generator, hence all of $J$. Since $z_0N=0$,
$z_0x_G=z_0\rho(x_G)$. Apply $\ell$ to obtain the second
inclusion from the first.

Write $\lambda(G)=\rho(x_G)$. Define
\[
 \beta(G,H)=\ell(x_G)\lambda(H)\pmod {u^{-2b}V}
                         \quad\hbox{in }F/u^{-2b}V.
                                                               \tag{B45.3}
\]
This additive pairing is $A$-balanced. If $a\in A$ and
$a_0=\operatorname{pr}_0(a)$, the difference between
$\beta(aG,H)$ and $\beta(G,aH)$ equals
\[
 \ell\bigl((a_0-\rho(a_0))x_G\bigr)\lambda(H).
\]
The first factor lies in $u^{-b}V$ by (B45.2), since
$a_0-\rho(a_0)\in J$, and the second does also.
Their product therefore vanishes in the indicated quotient.
The quotient is an $A$-module through $\rho\operatorname{pr}_0$,
whose image is in $V$.

Choose an integer $r>2c+2b$. We construct a polynomial
$g_r\in F[X]$ with
\[
 g_r(0)=0,\quad g_r'(0)=(-u^{-1})^r,\quad
 g_r(V)\subset V,\quad g_r(O)\subset V+u^{-r}J.             \tag{B45.4}
\]
In characteristic $p$, choose a power $D$ of the residue
cardinality with $D\ge h$ and set
$\delta(X)=(X^D-X)/u$, $g_r=\delta^{\circ r}$.
The Frobenius identity and $N^h=0$ give
$\delta(a+n)=\delta(a)-n/u$ for $a\in V,n\in N$;
all assertions follow.

In mixed characteristic choose a power $D$ of the residue
cardinality so large that $D\ge h$ and
\[
 v(D)\ge (r-1)(h-2)+\max_{1\le j<h}v(j).
\]
Again use $g_r=\delta^{\circ r}$. The elementary identity
$\binom Dj=(D/j)\binom{D-1}{j-1}$ gives
$v(\binom Dj)\ge v(D)-v(j)$.
For $0\le i<r$, $a\in V$ and $n\in u^{-i}J$, expand
\[
 \delta(a+n)=\delta(a)+\frac{Da^{D-1}-1}{u}n
          +\sum_{2\le j<h}\frac{\binom Dj a^{D-j}}u n^j.
\]
The linear term belongs to $u^{-(i+1)}J$, and so does
each higher term: its valuation bound is
$v(\binom Dj)-ij-1\ge-i-1$.
Thus $\delta(V+u^{-i}J)\subset V+u^{-(i+1)}J$.
Induction proves (B45.4); its derivative follows from
$\delta(0)=0$ and $\delta'(0)=-u^{-1}$.

Extend polynomials from $Q_0$ to $Q(A)$ by zero on the
other factors, and put
\[
       F_r=u^c e z_0g_r,\qquad H_r=u^{r+c}eg_r.
\]
Both are in $\operatorname{Int}(A)$. For $H_r$ this is
(B45.1) and (B45.4). For $F_r$, multiplication by $z_0$
kills the nilpotent part of $g_r(a)$, leaving
$u^cz_0g_r(\rho(a))\in u^cO$.
The polynomial image of
\[
                   F_r\otimes H_r-H_r\otimes F_r
\]
is zero, because $F_r=u^{-r}z_0H_r$ in $Q(A)[X]$.
Its value under (B45.3) is $u^{2c-r}$: the second term
vanishes since $\rho(z_0)=0$, and the two signs cancel
in the first. This is nonzero modulo $u^{-2b}V$ by the
choice of $r$. The required tensor kernel is proved.

The structural inputs were opened in Stacks, Theorem 10.160.8
(Cohen structure, tag 032A), Lemma 10.160.2 (complete finite
algebras), and Lemma 10.162.8 (complete local rings are Nagata,
tag 032W). The lattice and obstruction arguments above are
given in full and do not invoke flatness of integer-valued
polynomial rings.
\end{proof}

\begin{saresult}{B.046}{Proved}{Complete one-dimensional fraction-field classification without finite normalization}
Let $R$ be an infinite one-dimensional Noetherian domain,
and put $K=\operatorname{Frac}(R)$.
The compatible coaddition and comultiplication tensors for
$\operatorname{Int}(R)$ are uniquely determined if and only
if, at each maximal ideal with finite residue field,
the completion $\widehat{R_{\mathfrak m}}$ is a domain.
Equivalently, at each such ideal the normalization of
$R_{\mathfrak m}$ is finite and local. In every positive
case all multivariable tensor maps are isomorphisms and
the compatible plethory is unique.
\end{saresult}
\begin{proof}
By B.033 the tensor question localizes, so fix a
finite-residue maximal ideal and write its local ring
as $R$. If its completion is nonreduced, B.045 applies:
the completion is one-dimensional Cohen--Macaulay, and
in mixed characteristic the nonzero element $p\in R$
remains a nonzerodivisor by flatness. B.044 then proves
failure of tensor injectivity over $R$.

If the completion is reduced, $R$ is analytically
unramified and its normalization $T$ is finite
(Stacks, Lemma 10.162.10(4), tag 032Y, whose full proof
was opened). B.025 proves that the tensors work exactly
when $T$ is local. This is equivalent to the completion
of $R$ being a domain. In one direction, when $T$ is
local it is a discrete valuation ring, and exact
completion embeds $\widehat R$ in the complete discrete
valuation ring $\widehat T$. Conversely, if
$\widehat R$ is a domain, flatness embeds
$T\otimes_R\widehat R$ into
$K\otimes_R\widehat R$, which is a field. But the
completion of the finite semilocal algebra $T$ is the
product of its nonzero completions at its maximal
ideals. A subring of a field cannot have the nontrivial
idempotents of a product with more than one factor.
Thus $T$ is local.

At maximal ideals of infinite residue field the ordinary
Vandermonde argument gives the polynomial ring, as
already proved. In the positive finite-residue case,
B.024 gives every multivariable tensor isomorphism.
Testing kernels and cokernels at every maximal ideal
proves the global assertion, and B.009 supplies the
unique plethory compatibility.
\end{proof}


\begin{saresult}{B.047}{Proved}{Completion with an arbitrary rational coefficient overring}
In the setting of B.044, let $R\subset L\subset K$ be
any overring and put $B=L\otimes_R A\subset Q$.
For every $n$ there is a natural isomorphism
\[
 \Omega_n(R,L)\otimes_R A\simeq\Omega_n(A,B).
\]
Here $\Omega_n(A,B)$ denotes the polynomials in $B[X_1,
\ldots,X_n]$ taking $A^n$ into $A$. The tensor maps
correspond, and existence of the coaddition and
comultiplication expressions passes to this base change.
\end{saresult}
\begin{proof}
Flatness embeds $B=L\otimes_R A$ in $K\otimes_R A=Q$.
Within a fixed finite rectangular degree bound with $N$
monomials, the coefficient module of $\Omega_n(R,L)$ is
the intersection of the bounded-degree integer-valued
module with $L^N$, inside $K^N$. Flat base change preserves
this intersection, by tensoring the kernel sequence for
the difference map from the direct sum of the two
submodules to $K^N$. B.044 identifies the first module
after base change, and the second becomes $B^N$.
Take the filtered union over degree bounds. These
identifications preserve multiplication in disjoint
variables and polynomial substitution. Finally every
element of the algebraic tensor product is a finite
$A$-linear combination of original polynomials. The
expressions for $F(X+Y)$ and $F(XY)$ therefore extend
by this linearity. No finite generation of $L$ is used.
\end{proof}

\begin{saresult}{B.048}{Proved}{Finite-rank modules over a complete discrete valuation ring}
Let $V$ be a complete discrete valuation ring with
uniformizer $t$ and fraction field $F$. Let $W$ be a
finite-dimensional $F$-vector space, and let $M\subset W$
be a $V$-submodule containing a full $V$-lattice.
There is an $F$-vector subspace
$D=\bigcap_{j\ge0}t^jM$ such that $M/D$ is a finite
free $V$-module. More precisely, after a choice of
$F$-basis, $M$ is a direct sum of finitely many
fractional ideals $t^aV$ and finitely many copies of $F$.
\end{saresult}
\begin{proof}
Choose a full lattice $\Lambda\subset M$ and identify
$(W,\Lambda)$ with $(F^d,V^d)$. Put $E=F/V$.
We first record the elementary duality facts used below.
The divisible module $E$ is injective over the principal
ideal domain $V$. For a direct proof of the extension
property, extend a homomorphism maximally by Zorn's
lemma. To adjoin a new element $x$, its annihilator
relative to the old submodule is an ideal $(a)$.
If $a\ne0$, divisibility solves $ay=f(ax)$; if the
ideal is zero choose any $y$. Sending $x$ to $y$
extends the map, a contradiction unless the domain was
already the whole module.
Also $E$ detects every nonzero torsion module: a nonzero
element has annihilator $(t^s)$, its cyclic submodule
maps nontrivially to $t^{-s}V/V$, and injectivity extends
this map to the whole module.

An endomorphism of $E$ is multiplication by a unique
element of $V$. Indeed its value on $t^{-s}+V$ is
$a_s t^{-s}+V$ for $a_s\in V/t^sV$, and compatibility
is exactly $a_{s+1}\equiv a_s\pmod {t^s}$. Completeness
identifies the compatible sequence with an element of
$V$. Consequently every homomorphism $E^d\to E$
is induced by a vector in $V^d$.

Let
\[
 M^\#=\{\phi\in W^*: \phi(M)\subset V\}
                                      \subset\Lambda^*.
\]
Then
\[
 M=\{x\in W: \phi(x)\in V\text{ for every }
                                      \phi\in M^\#\}.       \tag{B48.1}
\]
Only the reverse inclusion needs proof. If $x\notin M$,
the nonzero torsion element $x+M$ in $W/M$ is detected
by a homomorphism to $E$. Composing with
$W/\Lambda\simeq E^d$ and using the preceding
endomorphism computation gives an $F$-linear form
$\phi$ integral on $M$ but not on $x$.

The submodule $M^\#\subset\Lambda^*$ is finite free,
since $V$ is Noetherian. Smith reduction of its matrix
inside the finite free module $\Lambda^*$ puts it in
the form
$t^{a_1}Ve_1^*\oplus\cdots\oplus t^{a_s}Ve_s^*$
with zero remaining coordinates. Formula (B48.1)
then gives
\[
 M=t^{-a_1}Ve_1\oplus\cdots\oplus t^{-a_s}Ve_s
           \oplus Fe_{s+1}\oplus\cdots\oplus Fe_d.
\]
Its intersection of all $t^j$ multiples is exactly the
last vector-space summand, proving every assertion.
\end{proof}

\begin{saresult}{B.049}{Proved}{Epimorphic overrings of a complete one-dimensional order}
Let $V$ be a complete discrete valuation ring with finite
residue field, let $A$ be a local commutative $V$-algebra
finite and torsion-free over $V$, and put
$Q=A\otimes_V\operatorname{Frac}(V)$.
For any overring $A\subset B\subset Q$, the multiplication
map $B\otimes_A B\to B$ is an isomorphism if and only if
\[
                            B=A+eQ
\]
for an idempotent $e\in Q$. Moreover injectivity of the
second polynomial tensor map for $\Omega(A,B)$ forces
this ring-epimorphism condition.
\end{saresult}
\begin{proof}
For the last assertion choose any $b\in B$. There is
a nonzerodivisor $d\in A$ with $db\in A$; a power
of the uniformizer of $V$ suffices. The finite ring
$A/dA$ has a monic null polynomial, for instance the
product of $X-a$ over a complete set of representatives.
Call it $P$. Then $P,bP\in\Omega(A,B)$. The tensor
\[
                       bP\otimes P-P\otimes bP
\]
has zero polynomial image, hence is zero if the tensor
map is injective. Taking the coefficient of degree
$\deg P$ in both factors gives
$b\otimes1=1\otimes b$ in $B\otimes_A B$.
This for every $b$ makes multiplication an isomorphism,
with inverse $b\mapsto b\otimes1$.

Now assume multiplication is an isomorphism.
Apply B.048 to $B\subset Q$, using the full finite
free lattice $A$. The divisible part
$D=\bigcap_j t^jB$ is an ideal of $B$ and a vector
space over $\operatorname{Frac}(V)$. It is in fact
an ideal of $Q$, since $Q=A[1/t]$. The quotient
$C=B/D$ is finite over $V$, hence finite over $A$.
The composite $A\to C$ is a ring epimorphism because
$A\to B$ and $B\to C$ are. If $C\ne0$, its fiber
$C/\mathfrak m_A C$ is nonzero by Nakayama and is a
finite-dimensional algebra over $k=A/\mathfrak m_A$.
The multiplication isomorphism on this fiber implies
$(\dim_k C/\mathfrak m_A C)^2=\dim_k C/\mathfrak m_A C$,
so the dimension is one and the unit generates it.
Nakayama applied to $C/\operatorname{im}(A)$ shows
that $A\to C$ is surjective. The same conclusion is
automatic when $C=0$. Therefore $B=A+D$.

Put $I=A\cap D$. It is a finite $V$-module. From
$B/A\simeq D/I$ and the epimorphism condition,
\[
              B\otimes_A(D/I)=B\otimes_A(B/A)=0.
\]
As $D$ is already a $Q$-module,
$B\otimes_A D\simeq Q\otimes_QD=D$.
Right exactness now gives $D=BI=I+DI$.
The quotient $D/DI$ is a vector space over
$\operatorname{Frac}(V)$ generated as a $V$-module
by the image of the finite module $I$. Such a vector
space must be zero: multiplication by $t$ is
surjective, and Nakayama applies to a finite
$V$-module. Hence $D=DI\subset D^2\subset D$.
An idempotent ideal in the Artinian ring $Q$ is
$eQ$: in each local Artinian factor a proper ideal
is nilpotent, so an idempotent ideal is either zero
or the whole factor. This proves necessity.

Conversely, for $B=A+eQ$, write $f=1-e$. Then
$B=eQ\times fA$. The map $A\to fA$ is a quotient,
so $fA\otimes_A fA\simeq fA$. Localization gives
$eQ\otimes_A eQ\simeq eQ$, and the cross terms
$eQ\otimes_A fA$ and $fA\otimes_A eQ$ vanish,
because after inverting $t$ they involve the
orthogonal factors $eQ$ and $fQ$. Multiplication
is therefore an isomorphism on the whole product.
\end{proof}

\begin{saresult}{B.050}{Proved}{All one-dimensional Noetherian domains and every torsion-free coefficient extension}
Let $R$ be an infinite one-dimensional Noetherian domain,
let $K=\operatorname{Frac}(R)$, and let $R\subset S$
be an $R$-torsion-free unital extension. Put
$L=S_{\mathrm{rat}}\subset K$ as in B.014.
The coaddition and comultiplication tensors are
uniquely determined if and only if at each maximal
ideal $\mathfrak m$ with finite residue field either
\[
 \begin{gathered}
 L_{\mathfrak m}=R_{\mathfrak m},\qquad\hbox{or}\\
 L_{\mathfrak m}=K\quad\hbox{and}\quad
 \widehat{R_{\mathfrak m}}\hbox{ is a domain}.
 \end{gathered}                                            \tag{B50.1}
\]
Equivalently, in the second alternative the local
normalization is finite and local. In every positive
case all tensor maps are isomorphisms and the
compatible plethory is unique. There is no
finite-normalization assumption on $R$.
\end{saresult}
\begin{proof}
B.014 reduces to $L$ and B.033 localizes the question.
At infinite-residue maximal ideals Vandermonde
interpolation gives the ordinary polynomial ring.
Fix a finite-residue local ring $R$, put
$A=\widehat R$, $Q=K\otimes_R A$, and
$B=L\otimes_R A\subset Q$. B.047 transfers the
question to $\Omega(A,B)$. The completion $A$ is
one-dimensional Cohen--Macaulay. As in the first
paragraphs of B.045, it is finite and torsion-free
over a complete discrete valuation subring with
finite residue field.

If the tensor map is injective, B.049 forces
$B=A+eQ$ for an idempotent $e$. If $e=0$, then
$B=A$ and faithful flatness gives $L=R$.
Suppose $e\ne0$. If any selected local Artinian
factor $Q_0$ is nonreduced, the two polynomials
$F_r,H_r$ constructed in B.045 have all coefficients
in that selected factor $e_0Q\subset eQ\subset B$.
They therefore belong to $\Omega(A,B)$, and the same
balanced functional proves a nonzero tensor kernel
there. Thus every selected factor must be a field.

Suppose there is some other generic factor besides
a selected field $Q_0$, and let $e_0$ be its
idempotent. Let $V\subset Q_0$ be the integral
closure of the complete discrete valuation subring.
It is a complete discrete valuation ring with finite
residue field, and the projection of $A$ lies in $V$.
Choose its uniformizer $u$ and $c\ge1$ with
$u^c e_0V\subset A$. Put $x=u^c e_0\in A$
and $r=c+1$.
The ideal $I_0=A\cap e_0Q$ is proper, since
$e_0\ne1$, and hence lies in $\mathfrak m_A$.
For every $h\in\Omega(A,B)$ the difference
$h(x)-h(0)$ has zero components outside $Q_0$;
as both values belong to $A$, it lies in $I_0$.

There is an integer-valued polynomial $H$ on $V$
such that $H(u^{cr})-H(0)$ is a unit. For example
the finite congruence-class indicator construction
in B.024 separates $0$ and $u^{cr}$ modulo a power
$u^M$ with $M>cr$. Define the polynomial, supported
on $Q_0$,
\[
                    G=u^c e_0\bigl(H-H(0)\bigr).
\]
Its coefficients lie in $e_0Q\subset B$, it takes
$A$ into $u^ce_0V\subset A$, and $G(0)=0$.
Moreover its selected valuation at $x^r$ is exactly
$c$. If multiplicative tensors existed for every
polynomial in $\Omega(A,B)$, iterating the binary
expressions would give an $r$-fold expression for
$G(X_1\cdots X_r)$. Apply the difference of evaluation
at $x$ and at zero in each variable. Every term on
the tensor side is in $\mathfrak m_A^r$, whereas
the other side is $G(x^r)$.
The projection of $\mathfrak m_A$ lies in the
maximal ideal of $V$, since this valuation ring
dominates the selected reduced local branch.
Thus $\mathfrak m_A^r$ has valuation at least
$r=c+1$, contradicting the value $c$ just found.
So this case fails even existence.

The only remaining possibility with $e\ne0$ is
that $Q$ itself is one field and $B=Q$.
Because $A\hookrightarrow Q$, this is exactly
the case that $A$ is a domain and $B=Q$.
Faithful flatness applied to $K/L$ gives $L=K$.
This proves necessity of (B50.1).

Conversely $L=R$ gives the ordinary polynomial ring.
If $L=K$ and the completion is a domain, B.046
gives all tensor isomorphisms. The same theorem
identifies this hypothesis with finite local
normalization. Applying these local conclusions
at every maximal ideal proves sufficiency globally
by testing kernels and cokernels. The unique
plethory compatibility is B.009.
\end{proof}

\begin{saresult}{B.051}{Proved}{Completion in every Noetherian dimension}
Let $(R,\mathfrak m)$ be a Noetherian local domain with
fraction field $K$, let $A=\widehat R$, $Q=K\otimes_R A$,
and let $R\subset L\subset K$ be any overring.
Put $B=L\otimes_R A\subset Q$. For every $n$,
\[
                    \Omega_n(R,L)\otimes_R A
                         \simeq\Omega_n(A,B).              \tag{B51.1}
\]
The tensor maps and existence of the two polynomial
co-operations pass through this identification. There is
no dimension or finite-normalization hypothesis.
\end{saresult}
\begin{proof}
If the residue field is infinite, a grid of points with
distinct residues gives an invertible Vandermonde
matrix over both local rings, so the integer-valued
modules are the ordinary polynomial rings. Assume
the residue field is finite.

For $j\ge0$ define
\[
 C_j=\{x\in K:\mathfrak m^jx\subset R\},\qquad
 C'_j=\{x\in Q:(\mathfrak m A)^jx\subset A\}.
\]
Flatness gives $C_j\otimes_R A=C'_j$.
Indeed, choose finitely many generators $a_1,\ldots,a_s$
of $\mathfrak m^j$ and tensor the kernel description
\[
 C_j=\ker\bigl(K\longrightarrow(K/R)^s,
                       x\longmapsto(a_ix\bmod R)_i\bigr).
\]

Every coefficient of an integer-valued polynomial in
$K[X_1,\ldots,X_n]$ lies in some $C_j$, by the
punctured-spectrum Vandermonde argument of B.034.
We verify the corresponding assertion over $A$ even
though $A$ need not be a domain and $Q$ need not be
Artinian. Each coefficient $x\in Q$ has a denominator
$d\in R\setminus\{0\}$, which is a nonzerodivisor
of $A$ by flatness. If $\mathfrak p$ is a nonmaximal
prime of $A$, its contraction to $R$ is not
$\mathfrak m$, since otherwise it contains
$\mathfrak m A$, the maximal ideal of $A$.
The image of $R$ modulo this contraction is infinite.
Choose a rectangular Vandermonde grid of $R$-points
whose determinant $v$ is outside $\mathfrak p$.
The values of the polynomial are in $A$, so Cramer's
rule gives $vx\in A$. Consequently the ideal
\[
                         I_x=\{a\in A:ax\in A\}
\]
is not contained in any nonmaximal prime. If it is
proper, its radical is the maximal ideal, and finite
generation gives $(\mathfrak m A)^j\subset I_x$
for some $j$. Thus $x\in C'_j$.
Taking a common $j$ handles all coefficients of a
given polynomial.

Fix this $j$ and a finite rectangular degree bound
with $N$ allowed monomials. Let $E_j\subset R$
represent the finite ring $R/\mathfrak m^j$.
For coefficients in $C_j$, the value modulo $R$
depends only on the inputs modulo $\mathfrak m^j$:
a monomial difference has a factor in that ideal,
which sends each coefficient into $R$.
The desired bounded module is therefore
\[
 \ker\bigl(C_j^N\longrightarrow(C_j/R)^{E_j^n}\bigr).
\]
Since $R/\mathfrak m^j\simeq A/(\mathfrak m A)^j$,
the same finite tests apply over $A$. Tensoring the
kernel gives exactly the bounded module there.
Take the filtered union over $j$ and the degree
bounds to prove (B51.1) for $L=K$.
For arbitrary $L$, intersect the bounded coefficient
module with $L^N$; flatness preserves this finite
intersection, as in B.047. Tensor compatibility and
existence follow from multiplication and finite
$A$-linear combination exactly as in B.047.
\end{proof}

\begin{saresult}{B.052}{Proved}{Epimorphic enlargements supported at the closed point}
Let $(A,\mathfrak m)$ be a complete Noetherian local
ring, containing a complete discrete valuation ring
$V$ with uniformizer $t$, with $t\in\mathfrak m$
a nonzerodivisor on $A$. Assume the residue field
of $V$ maps isomorphically to $A/\mathfrak m$.
Let $W\subset A$ be a multiplicative set of
nonzerodivisors containing $t$, and put $Q=W^{-1}A$.
Suppose
\[
 A\subset B\subset Q,\qquad
 \text{every element of }B/A\text{ is killed by a power
 of }\mathfrak m.
\]
Then there is a finite-dimensional
$F=\operatorname{Frac}(V)$ vector space $D\subset B$
which is an ideal of $Q$, such that $B/D$ is a finite
$A$-module and $A\cap D$ is a finite $V$-module.
If $A\to B$ is a ring epimorphism, then
\[
                         B=A+eQ,\qquad D=eQ                 \tag{B52.1}
\]
for an idempotent $e\in Q$, and $eQ$ is finite-dimensional
over $F$. If the common residue field is finite,
injectivity of the polynomial tensor map for
$\Omega(A,B)$ implies this epimorphism condition.
\end{saresult}
\begin{proof}
Put $T=B/A$. Since every class is killed by a power
of $\mathfrak m$, it is killed by a power of $t$,
and $B\subset A[1/t]$.
The $t$-torsion subgroup $T[t]$ embeds in $A/tA$:
send the class of $a/t$ to $a\bmod tA$.
Its image is contained in the submodule of $A/tA$
killed elementwise by powers of $\mathfrak m$.
The ascending chain of annihilator submodules
$(0:_{A/tA}\mathfrak m^j)$ stabilizes, since
$A/tA$ is Noetherian. The stable module is killed
by one power of $\mathfrak m$, and hence has finite
length. In particular $T[t]$ is a finite-dimensional
vector space over $V/tV$.

Let $E=F/V$. By the injectivity of $E$ proved in
B.048, an embedding $T[t]\hookrightarrow E[t]^s$
extends to a map $T\to E^s$. This extension is
injective: any nonzero kernel element, being killed
by a power of $t$, has a nonzero multiple killed
by $t$, contradicting injectivity on $T[t]$.
The same double-annihilator argument as in B.048
classifies submodules of $E^s$. Explicitly, annihilate
$T$ by a submodule of
$\operatorname{Hom}_V(E^s,E)=V^s$; cogeneration by
$E$ shows that taking the annihilator again gives
$T$. Smith reduction of that finite submodule of
$V^s$ shows
\[
                            T\simeq E^r\oplus T_0,
\]
where $T_0$ has finite length over $V$.
Its maximal divisible submodule $T_{\mathrm{div}}$
is $E^r$ and is canonical, hence stable under $A$.

The ring $A$ is also $t$-adically complete and
separated. To see completeness, a $t$-adic Cauchy
sequence is $\mathfrak m$-adically Cauchy; its limit
has the required congruences because each ideal
$t^jA$ is $\mathfrak m$-adically closed.
Closedness and separatedness are the Krull
intersection theorem applied to finite modules.
There is consequently a natural isomorphism
\[
 \operatorname{Hom}_V(E,A[1/t]/A)\simeq A.                  \tag{B52.2}
\]
Indeed the image of $t^{-j}+V$ is
$a_j/t^j+A$ with $a_j\in A/t^jA$; the homomorphism
condition is the compatibility of these residues,
and completeness gives a unique $a\in A$.

Choose the $r$ coordinate injections of
$T_{\mathrm{div}}\simeq E^r$ into $A[1/t]/A$.
By (B52.2) they are multiplication by elements
$a_1,\ldots,a_r\in A$. Put
$D=Fa_1+\cdots+Fa_r$.
These vectors are $F$-linearly independent:
a relation, after scaling to coefficients in
$V$ with at least one unit, would give a nonzero
kernel on $E^r$.
Their fractional multiples lie in $B$ because
their classes lie in $T$. Moreover
\[
              A\cap D=Va_1+\cdots+Va_r,\qquad
              D/(A\cap D)=T_{\mathrm{div}}.
\]
The first equality follows from injectivity of
the coordinate map $E^r\to T$.
Since $B/(A+D)\simeq T_0$ is finite over $V$,
$B/D$ is a finite $A$-module.

Krull intersection on this finite module shows
$\bigcap_j t^j B=D$. Thus $D$ is an ideal of
$B$, and in particular stable under $A$.
Every $w\in W$ acts injectively on the finite
$F$-vector space $D$, because it is a unit of $Q$.
It therefore acts bijectively. Hence $D$ is stable
under $W^{-1}A=Q$, so it is an ideal of $Q$.

Suppose now $A\to B$ is a ring epimorphism.
The finite $A$-algebra $B/D$ is an epimorphic image
of $A$. If it is zero the map from $A$ is surjective;
otherwise the residue-field dimension argument
and Nakayama proof of B.049 show that $A\to B/D$
is surjective. Thus $B=A+D$. With $I=A\cap D$,
the same tensor computation gives
\[
 B\otimes_A(B/A)=0,\qquad B\otimes_A D=D,\qquad
                         D=BI=I+DI.
\]
Here $B\otimes_A D=D$ follows by localizing at
$W$ on the already $Q$-module $D$.
The quotient $D/DI$ is an $F$-vector space generated
over $V$ by the finite module $I$, so it vanishes.
Therefore $D=D^2$.
Since $D$ is finite-dimensional over $F\subset Q$,
it is finitely generated as an ideal of $Q$.
A finitely generated idempotent ideal is generated
by an idempotent: write generators $d_i$ as
$d_i=\sum a_{ij}d_j$ with $a_{ij}\in D$;
the adjugate identity says $\det(1-(a_{ij}))D=0$.
Write this determinant as $1-e$, with $e\in D$.
Then $(1-e)e=0$, so $e^2=e$, and $D=eQ$.
This proves (B52.1).

Finally suppose the residue field is finite and
the polynomial tensor map is injective. For any
$b\in B$ choose $j$ with $\mathfrak m^j b\subset A$.
A monic null polynomial of the finite ring
$A/\mathfrak m^j$ gives $P,bP\in\Omega(A,B)$.
Top-coefficient extraction from
$bP\otimes P-P\otimes bP$ gives
$b\otimes1=1\otimes b$ exactly as in B.049.
Thus $B\otimes_A B\simeq B$.
\end{proof}

\begin{saresult}{B.053}{Proved}{The higher-dimensional local classification without finite normalization}
Let $(R,\mathfrak m)$ be a Noetherian local domain
of dimension at least two with finite residue field,
let $K=\operatorname{Frac}(R)$ and $R\subset L\subset K$.
Put
\[
 R^\sharp=\{x\in K:\mathfrak m^jx\subset R
                       \text{ for some }j\},\qquad
 E=L\cap R^\sharp.
\]
The polynomial coaddition and comultiplication
tensors are uniquely determined if and only if
$E=R$, equivalently $\Omega(R,L)=R[X]$.
No assumption on normalization is needed.
\end{saresult}
\begin{proof}
The coefficient argument of B.034 gives
$\Omega(R,L)=\Omega(R,E)$, and $E=R$ is precisely
the polynomial-ring case, which is positive.
For necessity put $A=\widehat R$,
$Q=K\otimes_R A$ and $B=E\otimes_R A\subset Q$.
B.051 identifies the integer-valued modules and
their tensor maps. Every element of $B/A$ is
killed by a power of the maximal ideal of $A$.

There is a complete coefficient discrete valuation
subring $V\subset A$ to which B.052 applies.
In characteristic $p$ use the coefficient field
$k\subset A$ and $V=k[[t]]$, where
$0\ne t\in\mathfrak m\subset R$; flatness makes
$t$ a nonzerodivisor, so this power-series map
is injective. In mixed characteristic use the
Cohen coefficient ring with uniformizer $p$;
$p\in R$ is again a nonzerodivisor. In both
cases the uniformizer belongs to the
multiplicative set $W=R\setminus\{0\}$ defining
$Q=W^{-1}A$.

If the tensor map is injective, B.052 yields
$B=A+eQ$, where $eQ$ is finite-dimensional over
$F=\operatorname{Frac}(V)$ and is contained in
$A[1/t]$. If $e=0$, faithful flatness gives
$E=R$, as required. Suppose $e\ne0$ and choose
a local Artinian factor $Q_0=e_0Q$ of the
finite-dimensional algebra $eQ$.
The image $A_0=e_0A\subset Q_0$ is a quotient
of $A$, hence a complete Noetherian local ring,
and $A_0[1/t]=Q_0$.
It is finite over $V$: apply B.048 to this
submodule of the finite-dimensional $F$-space
$Q_0$. A vector-space summand would contradict
$t$-adic separatedness of the finite $A$-module
$A_0$. Thus $A_0$ is a one-dimensional complete
local order, with the same residue field as $A$.

We detail why the selected-factor arguments remain
valid although the whole ring $A$ need not be
finite over $V$. If $Q_0$ is nonreduced, apply
the coefficient-field and finite-order construction
in B.045 to the one-dimensional ring $A_0$.
It first chooses a new coefficient discrete valuation
ring $V_0\subset A_0$ and a coefficient field
$F_1\subset Q_0$ whose embedding agrees with
$\operatorname{Frac}(V_0)\subset Q_0$. This compatibility
is supplied by the separable primitive-element and
nilpotent Newton-lifting argument of B.045.
Let $V_1\subset F_1$ be the finite integral closure
of $V_0$, with uniformizer $u$. The ring
$O=V_1[A_0]\subset Q_0$ is finite over $V_1$,
contains $A_0$, and spans $Q_0$ over $F_1$.
It is also finite over $A_0$: a finite set of
$V_0$-module generators of $V_1$ generates $O$
as an $A_0$-module, because the two copies of
$V_0$ agree. Equivalently, these generators are
integral over $V_0\subset A_0$. Since $A_0$ is
finite over the original $V$, so is $O$.
No inclusion $V_1\subset A_0$ is asserted.
Since $Q_0\subset eQ\subset A[1/t]$, a single
power $t^N$ sends $e_0O$ into $A$.
Multiplication by $t$ is invertible on $Q_0$;
on the finite free $V_1$-module $O$ its cokernel
is torsion. Hence, for some $c$, one has
\[
                        u^ce_0O\subset t^Ne_0O\subset A.
                                                               \tag{B53.1}
\]
All the subsequent steps of B.045 use only the
projection $A\to A_0\subset O$, the inclusion
(B53.1), and nilpotent inputs supported on $Q_0$.
They give polynomials $F_r,H_r$ supported on
$Q_0$, hence with coefficients in
$e_0Q\subset B$, taking all of $A$ into $A$.
The very same pairing (B45.3), with scalar action
through $A\to A_0\to V_1$, is $A$-balanced and
detects their nonzero kernel tensor.
Thus a selected nonreduced factor is impossible.

If $Q_0$ is a field, its valuation ring is a
finite complete discrete valuation overring of
$A_0$. The same denominator-clearing argument
gives $u^ce_0V_1\subset A$ for some $c\ge1$.
One must have $e_0\ne1$. Otherwise $Q=Q_0$
would be finite-dimensional over $F$, and B.048
applied to the $t$-adically separated submodule
$A\subset Q=A[1/t]$ would make $A$ finite
over $V$, contradicting $\dim A=\dim R\ge2$.
The supported-input finite-difference proof of
B.050 now applies word for word: with
$x=u^ce_0$, every $h(x)-h(0)$ lies in the
proper ideal $A\cap e_0Q\subset\mathfrak m_A$,
while a supported congruence indicator $G$ has
$G(0)=0$ and valuation exactly $c$ at
$x^{c+1}$. Iterating multiplicative tensor
expressions would put this value in
$\mathfrak m_A^{c+1}$, a contradiction.
Thus a selected field factor is also impossible
if all multiplicative expressions exist.

Both possibilities contradict the joint
injectivity-and-existence hypothesis. Hence
$e=0$, $B=A$, and $E=R$ by faithful flatness.
This proves necessity and the theorem.
\end{proof}

\begin{saresult}{B.054}{Proved}{Complete Noetherian-domain classification with all torsion-free coefficients}
Let $R$ be a nonfield Noetherian domain with
fraction field $K$, and let $R\subset S$ be an
$R$-torsion-free unital extension. Put
$L=S_{\mathrm{rat}}$.
The polynomial coaddition and comultiplication
tensors are uniquely determined exactly under
the following conditions, imposed only at
maximal ideals $\mathfrak m$ with finite residue field:
\[
\begin{gathered}
 \dim R_{\mathfrak m}=1:\quad
 L_{\mathfrak m}=R_{\mathfrak m}\quad\hbox{or}\\
 \bigl(L_{\mathfrak m}=K\text{ and }
        \widehat{R_{\mathfrak m}}\text{ is a domain}\bigr);\\[2mm]
 \dim R_{\mathfrak m}\ge2:\quad
 L_{\mathfrak m}\cap R_{\mathfrak m}^{\sharp}=R_{\mathfrak m}.
\end{gathered}                                                   \tag{B54.1}
\]
There is no finite-normalization or finite bound on
Krull dimension. In every positive case all
multivariable tensor maps are isomorphisms and
the compatible plethory is unique.
For $S=K$, the first condition says finite local
normalization with one branch, and the second
is exactly $\operatorname{depth}R_{\mathfrak m}\ge2$.
\end{saresult}
\begin{proof}
Use the rational-coefficient reduction B.014 and
the localization theorem B.033. Infinite-residue
local rings give the ordinary polynomial ring.
At finite-residue maximal ideals, B.050 gives
the complete dimension-one criterion and B.053
gives the complete higher-dimensional criterion.
In the higher-dimensional positive case the
module is the polynomial ring, so every tensor
map is an isomorphism. In dimension one this
was proved in B.050. Test kernels and cokernels
at every maximal ideal and apply B.009 for
the compatible plethory. Finally B.046 identifies
the dimension-one completion condition with
finite local normalization, and B.034 identifies
$R_{\mathfrak m}^{\sharp}=R_{\mathfrak m}$ with
depth at least two. These arguments are local
and do not require a bound on the dimensions
of all maximal localizations simultaneously.
\end{proof}

\begin{saresult}{B.055}{Proved}{Localization and completion with arbitrary coefficient extensions}
Let $R$ be a Noetherian commutative ring and $R\subset S$
an arbitrary unital inclusion. No torsion-free hypothesis is imposed.
For every multiplicative subset $U\subset R$ and every $n$,
\[
 U^{-1}\Omega_n(R,S)
       =\Omega_n(U^{-1}R,U^{-1}S).                         \tag{B55.1}
\]
If $(R,\mathfrak m)$ is local, set $A=\widehat R$ and
$S_A=S\otimes_R A$. Then
\[
 \Omega_n(R,S)\otimes_R A=\Omega_n(A,S_A).                \tag{B55.2}
\]
When the residue field is finite, every coefficient of an
integer-valued polynomial lies in the ring
\[
\begin{gathered}
 C(R,S)=\bigcup_{j\ge0}(R:_S\mathfrak m^j),\\
 (R:_S\mathfrak m^j)=\{s\in S:\mathfrak m^js\subset R\}.
\end{gathered}
                                                               \tag{B55.3}
\]
These rings commute with completion, and $C(R,S)/R$ is
elementwise killed by powers of $\mathfrak m$.
All tensor maps and the existence conditions commute with
the identifications above.
\end{saresult}
\begin{proof}
The degree induction in B.033 proves the forward localization
inclusion over arbitrary rings: the identity
$s^dF(X)-F(sX)$ cancels the terms of highest total degree,
and $s$ is invertible after localization. For the converse,
clear the finitely many coefficient denominators to obtain
$F\in S[X_1,\ldots,X_n]$ taking $R^n$ into $U^{-1}R$.
If $s_1,\ldots,s_N$ are its coefficients, the module
$M=R+Rs_1+\cdots+Rs_N\subset S$ is finite over $R$.
The classes of all values $F(r)$ generate a submodule of
the Noetherian module $M/R$. That submodule is finitely
generated and vanishes after localization. One element
$u\in U$ therefore kills all its generators, and
$uF(R^n)\subset R$. This proves (B55.1), without rational
coefficient denominators or any torsion assumption.

For completion, infinite residue field gives ordinary
polynomial rings by the unit Vandermonde argument. Assume
the residue field is finite. Write $C_j=(R:_S\mathfrak m^j)$.
These are submodules containing $R$, and
$C_jC_k\subset C_{j+k}$, so their union is a ring.
Choose generators $a_1,\ldots,a_s$ of $\mathfrak m^j$.
Flatness of completion applied to
\[
 C_j=\ker\bigl(S\longrightarrow(S/R)^s,
                   x\longmapsto(a_ix\bmod R)_i\bigr)
\]
identifies $C_j\otimes_R A$ with
$(A:_{S_A}(\mathfrak m A)^j)$.

For any coefficient $x$ of $F\in\Omega_n(R,S)$, and any
nonmaximal prime $\mathfrak p$ of $R$, the domain
$R/\mathfrak p$ is infinite. A rectangular grid of
distinct residues gives a Vandermonde determinant
$v\notin\mathfrak p$, and Cramer's rule gives $vx\in R$.
Thus the ideal $\{r\in R:rx\in R\}$ is either $R$ or
has radical $\mathfrak m$. Noetherianity then puts $x$
in some $C_j$. This argument also works over $A$:
a nonmaximal prime of $A$ cannot contract to
$\mathfrak m$, so the same $R$-grids give the required
determinants. It does not require $R$ or $A$ to be a domain.
If there are no nonmaximal primes, the argument says that
every proper denominator ideal has radical $\mathfrak m$,
which is exactly what is needed.

For fixed $j$ and a finite rectangular degree bound, the
integer-valued polynomials with coefficients in $C_j$
are the kernel of finitely many evaluations modulo $R$,
at representatives of the finite ring $R/\mathfrak m^j$.
Changing an input by $\mathfrak m^j$ changes every
monomial value by that ideal, and hence changes the
polynomial value by an element of $R$. Since
$R/\mathfrak m^j=A/(\mathfrak m A)^j$, the same finite
tests apply after completion. Tensor the finite evaluation
kernel, then take the filtered union over $j$ and the
degree bounds. This proves (B55.2) and (B55.3), including
the assertion about completion of the coefficient ring.
Polynomial multiplication and substitution give the stated
compatibility; existence for an $A$-linear combination
follows by the same finite linear combination of tensor
expressions. Faithful flatness detects failure of
injectivity and of the requisite image conditions.
\end{proof}

\begin{saresult}{B.056}{Proved}{Complete local epimorphic algebras supported at one point}
Let $(A,\mathfrak m)$ be a complete Noetherian local ring
with finite residue field. Let $A\subset C$ be a unital
ring epimorphism, and suppose every element of $C/A$ is
killed by a power of $\mathfrak m$. Define
\[
 D=\bigcap_{j\ge0}\mathfrak m^j C,\qquad I=A\cap D.
\]
Then
\[
 C=A+D,\qquad D=\mathfrak mD=D^2=eC
                                                               \tag{B56.1}
\]
for an idempotent $e\in C$, and there is an isomorphism
of $A$-algebras
\[
 C\simeq D\times A/I,\qquad
 a\longmapsto(ea,a+I).                                    \tag{B56.2}
\]
The ideal $D$ is finitely generated over $C$. No finiteness
assumption on $C$ itself is needed.
\end{saresult}
\begin{proof}
First, every map $A\to C/\mathfrak m^jC$ is surjective.
It is a ring epimorphism from the finite local ring
$A/\mathfrak m^j$. After reduction modulo its nilpotent
maximal ideal, an epimorphic algebra over the residue
field is either zero or that field: the linear-functional
argument of B.013 applied to $b\otimes1=1\otimes b$
proves this even for an infinite-dimensional algebra.
Thus $C/\mathfrak m^jC$ equals the image of $A$ plus
its maximal-ideal multiple. Iterating this equality
until the maximal ideal vanishes proves surjectivity.

Fix $c\in C$. For each $j$ the set
\[
 X_j=\{a\in A:c-a\in\mathfrak m^jC\}
\]
is nonempty and is a coset of an ideal containing
$\mathfrak m^j$. These are nested closed subsets of
the compact ring $A=\varprojlim A/\mathfrak m^j$.
For an elementary proof of the needed compactness,
choose $a_j\in X_j$, successively pass to infinite
subsequences with constant residue modulo
$\mathfrak m,\mathfrak m^2,\ldots$, and use completeness
on the resulting compatible residues. Their limit
belongs to every $X_j$. Hence $c-a\in D$ for some
$a\in A$, proving $C=A+D$.

The epimorphism condition gives
$C\otimes_A(C/A)=0$. Since $D/I=C/A$, right exactness
shows that $C\otimes_A I\to C\otimes_A D$ is surjective.
Compose with multiplication onto $D$ to get $D=CI$.
The ideal $I$ is finitely generated over $A$, so $D$
is a finitely generated ideal of $C$.

We next prove $D=\mathfrak mD$. There is an integer
$N\ge1$ for which
\[
                       I\cap\mathfrak m^N A\subset\mathfrak mI.
                                                               \tag{B56.3}
\]
Indeed the descending subspaces
$((I\cap\mathfrak m^jA)+\mathfrak mI)/\mathfrak mI$
of the finite-dimensional vector space $I/\mathfrak mI$
stabilize. Their intersection is zero by Krull
intersection on the finite module $A/\mathfrak mI$,
so their stable value is zero.
For $i\in I\subset D\subset\mathfrak m^NC$, write
$i=a+d$ with $a\in\mathfrak m^NA$ and
$d\in\mathfrak m^ND$, using $C=A+D$.
Then $a\in I\cap\mathfrak m^NA\subset\mathfrak mI$,
so $i\in\mathfrak mD$. Since $D=CI$, this gives
$D=\mathfrak mD$.

Also $D=CI=I+DI$. The quotient $D/DI$ is consequently
a finite $A$-module, generated by the image of $I$,
and it equals its $\mathfrak m$-multiple. Nakayama
gives $D=DI$, hence $D=D^2$.
The determinant argument of B.052 applies to this
finitely generated idempotent ideal and gives $D=eC$.
Finally $(1-e)C=(1-e)A$ because $C=A+D$; the kernel
of $A\to(1-e)A$ is precisely $I$. Orthogonal
idempotents now give (B56.2).
\end{proof}

\begin{saresult}{B.057}{Proved}{The torsion obstruction over a completed domain}
Let $(R,\mathfrak m_R)$ be a nonfield Noetherian local
domain with finite residue field, let $K$ be its fraction
field, and put $A=\widehat R$, $Q=K\otimes_R A$.
Let $A\subset C$ be any unital extension such that
$C/A$ is elementwise killed by powers of the maximal
ideal of $A$. The polynomial coaddition and
comultiplication tensors for $\Omega(A,C)$ are uniquely
determined if and only if
\[
 C=A\quad\text{or}\quad
 \bigl(\dim R=1,\ A\text{ is a domain},\ C=Q\bigr),       \tag{B57.1}
\]
where the second identification is as $A$-algebras.
In both positive cases all multivariable tensor maps
are isomorphisms. In particular no nonzero effective
base torsion survives the uniqueness and existence
requirements.
\end{saresult}
\begin{proof}
For each $c\in C$, choose $j$ with
$\mathfrak m_A^jc\subset A$. A monic null polynomial
for the finite ring $A/\mathfrak m_A^j$ gives $P,cP$
in $\Omega(A,C)$. If the tensor map is injective,
the zero-image tensor $cP\otimes P-P\otimes cP$
and top-coefficient extraction give
$c\otimes1=1\otimes c$ in $C\otimes_A C$.
Thus $A\to C$ is a ring epimorphism, and B.056 applies.
Use its notation $C=A+D$, $D=eC$, $I=A\cap D$.

Set $W=R\setminus\{0\}$. A nonzero element of
$\mathfrak m_R$ kills every class of $C/A$ after
localization, so $W^{-1}C=W^{-1}A=Q$.
Let $T$ be the kernel of $C\to Q$, and let $B$ be
its image. Each element $z\in T$ is killed by a
power of $\mathfrak m_A$: some such power sends
$z$ into $A$, where $W$ acts without torsion by
flatness of completion. Since $D=\mathfrak m_A^jD$
for all $j$, one has $TD=0$. It follows that
$D\cap T=0$, because $e$ acts as the identity on
$D$. We may therefore identify $D$ with its image
in $Q$.

Choose the complete coefficient discrete valuation
ring $V\subset A$ used in B.053: in characteristic
$p$ take $k[[t]]$ for $0\ne t\in\mathfrak m_R$,
and in mixed characteristic take the Cohen ring
with uniformizer $p$. Write $F=\operatorname{Frac}(V)$.
The map $A\to B$ is epic, as a composite of epic
maps. B.052 gives its canonical finite-dimensional
$F$-vector ideal $D'\subset Q$, with
$B=A+D'$ and $B/D'$ finite over $A$.
The image of $D$ in $B/D'$ is a finite $A$-module
and equals its maximal-ideal multiple, so it is
zero by Nakayama. Conversely $B/D$ is a quotient
of $A$, hence finite over $A$. The image there
of the $t$-divisible module $D'$ lies in every
$t^j(B/D)$, hence is zero by Krull intersection.
Consequently $D=D'$ is an ideal of $Q$, finite
dimensional over $F$, and contained in $A[1/t]$.

Let $J=A\cap D$ now mean the intersection of the
images inside $Q$, to distinguish it from $I$,
which was the intersection inside $C$. Then
$I\subset J$, and
\[
                              J/I\simeq T.                \tag{B57.2}
\]
Explicitly, for $a\in J$, subtract its unique lift
in $D\subset C$ from $a\in A\subset C$; this
gives every element of $T$, with kernel $I$.
B.052 says $J$ is a finite $V$-lattice in $D$.
Its quotient $J/I$ is torsion, hence killed by
one power of $t$. Therefore $I$ is also a full
finite $V$-lattice in $D$.

Suppose $D\ne0$ and select a local Artinian factor
$Q_0=e_0Q$ of the finite-dimensional algebra $D$.
The order $A_0=e_0A$ is a complete local quotient
of $A$, and $A_0[1/t]=Q_0$, since
$Q_0\subset A[1/t]$. B.048 and $t$-adic
separatedness show that $A_0$ is finite over $V$,
exactly as in B.053.

If $Q_0$ is nonreduced, use the compatible
coefficient field and finite order $O\subset Q_0$
constructed in B.045. The corrected integrality
argument in B.053 proves that $O$ is finite over
the original $V$: the normalized coefficient DVR
is finite integral over the newly chosen DVR
inside $A_0$, with compatible scalar embeddings.
As $I$ is a full lattice in $D$, some $t^N e_0O$
is contained in $I$. If $u$ is the new uniformizer,
finite generation of $O$ and invertibility of
$t$ on $Q_0$ give
\[
                              u^ce_0O\subset I             \tag{B57.3}
\]
for some $c$. The construction of B.045 now gives
its two polynomials supported on $Q_0\subset D$
with all their values in $I\subset A$.
They belong to $\Omega(A,C)$. Its balanced
derivative functional still applies: arbitrary
polynomials in $\Omega(A,C)$ have their values
in $A$, and projection to $Q_0$ puts those values
in $A_0\subset O$. The input and valuation
bounds use only (B57.3). The same nonzero tensor
kernel is therefore obtained. Thus every selected
factor must be a field.

Assume next that $I$ is a proper ideal of $A$.
For a selected field factor, let $V_1$ be the
finite discrete valuation normalization of $A_0$,
with uniformizer $u$. Clearing its finite lattice
into $I$ gives, after increasing $c$ if necessary,
\[
                     u^ce_0V_1\subset I,\qquad c\ge1.
\]
Put $x=u^ce_0\in I$. By (B56.2), the second
component of any $H\in\Omega(A,C)$ is a
polynomial over $A/I$. Its evaluations at $x$
and $0$ are equal. Since the actual values belong
to $A$, this proves
\[
                         H(x)-H(0)\in I.                 \tag{B57.4}
\]
Set $r=c+1$. An integer-valued congruence indicator
on $V_1$, as constructed in B.024, distinguishes
$0$ and $u^{cr}$. Thus there is a polynomial
$G$ supported on $Q_0$ with all values in
$u^ce_0V_1\subset I$, with $G(0)=0$, and with
selected valuation $v(G(x^r))=c$.
If every binary multiplicative expression existed,
iteration would express $G(X_1\cdots X_r)$ as
a sum of products of one-variable elements of
$\Omega(A,C)$. Taking the difference at $x$ and
$0$ in each variable and using (B57.4) would
put $G(x^r)$ in $I^r$. But $I\subset\mathfrak m_A$,
the map $A\to A_0$ is local, and $V_1$ dominates
$A_0$. Every element of $I^r$ has selected
valuation at least $r=c+1$, a contradiction.
Therefore $D\ne0$ forces $I=A$.

If $D=0$, then $C=A$. If $I=A$, then $1\in D$,
so $C=D$ and its image in $Q$ is the whole of
$Q$. Thus $C=Q$, finite-dimensional over $F$.
B.048 applied to the separated submodule
$A\subset Q=A[1/t]$ makes $A$ finite over $V$;
therefore $\dim R=\dim A=1$.
Nonreduced factors have already been excluded.
If $Q$ had two field factors, the supported
finite-difference argument of B.050, with the
proper ideal $A\cap e_0Q$, would exclude
multiplicative existence. Hence $Q$ is a field,
and its subring $A$ is a domain. This proves
the necessity of (B57.1).

Conversely $C=A$ gives the polynomial ring. In
the second case $Q$ is the fraction field of
the complete one-dimensional domain $A$, whose
normalization is finite and local. B.024 and
B.050 give every multivariable tensor isomorphism
and the required compatible co-operations.
\end{proof}

\begin{saresult}{B.058}{Proved}{Every Noetherian domain with arbitrary coefficient extensions}
Let $R$ be a nonfield Noetherian domain with fraction
field $K$, and let $R\subset S$ be an arbitrary unital
inclusion, allowing base torsion in $S$. For each maximal
ideal $\mathfrak m$ with finite residue field define
the actual coefficient subring of $S_{\mathfrak m}$
\[
 C_{\mathfrak m}
   =\{s\in S_{\mathfrak m}:\mathfrak m^jR_{\mathfrak m}s
              \subset R_{\mathfrak m}\text{ for some }j\}.
\]
The polynomial coaddition and comultiplication tensors
are uniquely determined if and only if, for every such
$\mathfrak m$,
\[
\begin{gathered}
 C_{\mathfrak m}=R_{\mathfrak m},\quad\text{or}\\
 \dim R_{\mathfrak m}=1,\qquad
 \widehat{R_{\mathfrak m}}\text{ is a domain},\qquad
 C_{\mathfrak m}=K.                                      \tag{B58.1}
\end{gathered}
\]
The last equality means the canonical copy of $K$
extending $R_{\mathfrak m}$ inside $S_{\mathfrak m}$.
No torsion-free or finite-normalization hypothesis is
imposed. In every positive case all multivariable
tensor maps are isomorphisms and the compatible
plethory is unique. Together with B.013--B.014, this
settles every Noetherian domain base and every
commutative coefficient extension.
\end{saresult}
\begin{proof}
B.055 permits localization and completion for this
arbitrary extension. At an infinite-residue maximal
ideal the unit Vandermonde argument gives the ordinary
polynomial ring. At a finite-residue maximal ideal,
replace $S_{\mathfrak m}$ by $C_{\mathfrak m}$, which
contains every possible coefficient by B.055.
After completion the latter ring has quotient over
the base killed elementwise by maximal-ideal powers,
so B.057 gives exactly its two positive possibilities.

Faithful flatness descends them. If
$C_{\mathfrak m}\otimes A=A$, then
$(C_{\mathfrak m}/R_{\mathfrak m})\otimes A=0$,
so $C_{\mathfrak m}=R_{\mathfrak m}$.
If $C_{\mathfrak m}\otimes A=K\otimes A$, the
kernel of $C_{\mathfrak m}\to K\otimes C_{\mathfrak m}$
vanishes after faithful flat base change, so it was
zero. Its image lies in $K$, since every coefficient
is sent into $R_{\mathfrak m}$ by a nonzero element
of a maximal-ideal power. Now
$(K/C_{\mathfrak m})\otimes A=0$ proves
$C_{\mathfrak m}=K$ as the canonical subring.
Thus the completed conditions descend precisely
to (B58.1), and conversely these local conditions
give the positive cases of B.057.

The local positive tensor isomorphisms imply the
global ones by testing kernels and cokernels at
every maximal ideal. The existence identities
and their compatibility follow as in B.009.
In the torsion-free case
$C_{\mathfrak m}=S_{\mathrm{rat},\mathfrak m}
\cap R_{\mathfrak m}^{\sharp}$, so this criterion
recovers B.054 while allowing arbitrary additional
torsion that never occurs in an integer-valued
polynomial coefficient.
\end{proof}

\begin{saresult}{B.059}{Proved}{A nilpotent detector without a coefficient-field section}
Let $(A,\mathfrak m)$ be a complete one-dimensional
Noetherian Cohen--Macaulay local ring with finite
residue field, in arbitrary characteristic. Write
$Q=Q(A)$ for its total quotient ring. If $A$ is
nonreduced, the polynomial tensor map for
$\operatorname{Int}(A)\subset Q[X]$ is not injective.
This includes generic local factors of characteristic
$p^a$ for $a>1$ and mixed-characteristic rings in
which $p$ is a zero divisor.

The same detector applies to a selected nonreduced
local Artinian ring $Q_0$ as an idempotent factor
of a coefficient extension of an ambient complete local
ring, provided
its image order $A_0$ is a complete one-dimensional
Cohen--Macaulay quotient and a finite order in $Q_0$
has a nonzero conductor into the ambient ring.
\end{saresult}
\begin{proof}
Choose a regular parameter $t\in\mathfrak m$.
Then $Q=A[1/t]$ is Artinian. A nonzero nilpotent
of $A$ stays nonzero there, so some local factor
$Q_0=e_0Q$ is nonreduced. Let $N$ be its nilradical,
let $N^h=0$ with $h\ge2$, and write $F=Q_0/N$.
The image $A_0=e_0A$ is a complete local quotient,
and $A_0[1/t]=Q_0$. Its reduced quotient is a
complete one-dimensional domain with fraction field
$F$. Its normalization $V\subset F$ is a finite
complete discrete valuation ring with finite residue
field, by the complete-local Nagata theorem and the
finite-algebra decomposition cited in B.045.

We construct a finite order $O\subset Q_0$ mapping
onto $V$, without embedding $F$ in $Q_0$.
Take finite integral generators $v_i$ of $V$ over
the reduced order and arbitrary lifts $y_i\in Q_0$.
Lift a monic integral equation for $v_i$ to $A_0$.
Its value at $y_i$ lies in $N$, so the $h$th power
of that monic polynomial vanishes at $y_i$.
Thus each $y_i$ is integral over $A_0$, and
\[
 O=A_0[y_1,\ldots,y_s],\qquad J=O\cap N,
 \qquad O/J=V                                             \tag{B59.1}
\]
is a finite local $A_0$-algebra. It is Noetherian
and complete. Choose $u\in O$ lifting a uniformizer
$\bar u$ of $V$. Then $u$ is invertible in $Q_0$,
regular on $O$, and a parameter of $O$; moreover
$O[1/u]=Q_0$. Clearing the finitely many denominators
of $e_0O\subset Q$ and then comparing the two
parameter ideals gives an integer $c\ge0$ with
\[
                             u^ce_0O\subset A.            \tag{B59.2}
\]
Explicitly, some $t^Ne_0O\subset A$, and the finite
length of $O/t^NO$ gives $u^cO\subset t^NO$.

The socle $\operatorname{Ann}_{Q_0}(N)$ has its
canonical $F$-vector-space structure: multiplication
by two lifts of the same residue differs by an
element of $N$, which annihilates it. Choose a
nonzero line $Fz\subset N\cap\operatorname{Ann}(N)$.
Its intersection with $O$ is a finite torsion-free
rank-one module over $V=O/J$, hence is generated
by some $z_0$:
\[
                           Fz_0\cap O=Vz_0.               \tag{B59.3}
\]
Here finiteness follows because the intersection is
an $O$-submodule of the Noetherian module $O$, and
its action factors through $V$. It has rank one
because multiplication by a power of $u$ puts any
chosen nonzero socle element into $O$.

For $G\in\operatorname{Int}(A)$ set
$x_G=\operatorname{pr}_0(G'(0))$. There is one
integer $b\ge0$, independent of $G$, for which
\[
             Jx_G\subset u^{-b}O,\qquad
             \rho(x_G)\in\bar u^{-b}V,                   \tag{B59.4}
\]
where $\rho:Q_0\to F$ is the residue map.
To prove this, choose finite $O$-generators $n_j$
of $J$, and evaluate $G$ at
$s_i u^ce_0n_j\in A$, for $h$ fixed scalars
$s_i\in O$ whose residues in $F$ are distinct.
For example, use $1,u,u^2,\ldots,u^{h-1}$;
their pairwise differences are units of $Q_0$.
The evaluation as a polynomial in $s$ has degree
at most $h-1$, since $n_j^h=0$. One fixed
Vandermonde inverse has entries in $u^{-d}O$
for some $d$. It therefore bounds the linear
coefficient $u^cn_jx_G$ uniformly. Increasing
$b$ proves the first assertion for every generator,
hence all of $J$. As $z_0x_G=z_0\rho(x_G)$,
(B59.3) gives the second assertion.

Define an $O$-submodule of $Q_0$ by
\[
 M=\{x\in Q_0:Jx\subset O,\ \rho(x)\in V\}.
\]
It contains $O$, and $J(M/O)=0$. Thus $M/O$
is canonically a $V=O/J$-module, without a
coefficient-field section. Set
\[
 y_G=u^b x_G,\qquad \lambda_G=\rho(y_G)\in V,
 \qquad \beta(G,H)=(y_G+O)\lambda_H\in M/O.               \tag{B59.5}
\]
This pairing is additive and $A$-balanced. Indeed
the action of $a\in A$ on $M/O$ factors through
$A\to A_0\to V$, while
$y_{aG}=a_0y_G$ and $\lambda_{aH}=\rho(a_0)\lambda_H$.
No additive or multiplicative splitting of $O\to V$
has been used.

Choose $r>2b+2c$, and put $q=|V/\bar uV|$.
We can choose $D=q^s\ge h$ so that
\[
 \binom Dj\in u^{(r-1)(h-2)}O\qquad(2\le j<h).           \tag{B59.6}
\]
For completeness, let $p$ be the residue
characteristic. Since $p\in\sqrt{uO}$, some
$p^a$ lies in $uO$, and arbitrarily high powers
of $p$ lie in any prescribed $u$-power ideal.
The integer identity
$\binom Dj=(D/j)\binom{D-1}{j-1}$ gives
$v_p(\binom Dj)\ge s v_p(q)-v_p(j)$.
Taking $s$ large proves (B59.6), whether $p$
is regular, nilpotent, or a zero divisor in $O$.

Let $\delta(X)=(X^D-X)/u$ and
$g_r=\delta^{\circ r}$. Reduction modulo $J$
shows $\delta(O)\subset O+u^{-1}J$, since
$a^D-a$ is divisible by $\bar u$ for every
$a\in V$. For $0\le i<r$, expand at
$a+n$ with $a\in O$, $n\in u^{-i}J$:
\[
 \delta(a+n)=\delta(a)+\frac{Da^{D-1}-1}{u}n
       +\sum_{2\le j<h}\frac{\binom Dj a^{D-j}}u n^j.
\]
The linear term lies in $u^{-(i+1)}J$.
Each higher term does too, by (B59.6), since
$(r-1)(h-2)\ge i(j-1)$. Consequently
\[
 g_r(O)\subset O+u^{-r}J,\qquad
 g_r(0)=0,\qquad g_r'(0)=(-u^{-1})^r.                    \tag{B59.7}
\]
Extend supported polynomials by zero on the other
factors of $Q$, and set
\[
                    F_r=u^ce_0z_0g_r,\qquad
                    H_r=u^{r+c}e_0g_r.
\]
Both take $A$ into $A$ by (B59.2)--(B59.7):
$H_r$ clears the indicated nilpotent denominator,
and $z_0$ kills the nilpotent part for $F_r$.
Their antisymmetric tensor has zero polynomial
image because $F_r=u^{-r}z_0H_r$ in $Q[X]$.
Its value under (B59.5) is
\[
                              z_0u^{2b+2c-r}+O.           \tag{B59.8}
\]
The other term has $\lambda_{F_r}=0$. The displayed
class is nonzero by (B59.3), because
$r>2b+2c$. Thus the tensor is nonzero.

For the selected-factor variant, all interpolation
values are projected from the ambient ring into
$A_0$, and the conductor (B59.2) ensures that
every interpolation input and every value of
$F_r,H_r$ belongs to the ambient ring. The same
balanced pairing acts on its entire integer-valued
module. If the conductor lands in an ideal of
the ambient ring, all values of the two supported
polynomials land in that ideal as well. This
verifies precisely the extra scope stated above.
\end{proof}

\begin{saresult}{B.060}{Proved}{The selected algebra is a one-dimensional punctured spectrum}
Under the hypotheses of B.056, suppose $D\ne0$,
and let $A_0=eA\subset D$ be the image order.
Then $A_0$ is a complete one-dimensional
Cohen--Macaulay local ring, and
\[
                         D=Q(A_0).                       \tag{B60.1}
\]
The ideal $I=A\cap D$ maps injectively into $A_0$
as an ideal containing a nonzerodivisor of $A_0$.
If $I$ is proper in $A$, it is also proper in $A_0$.
This conclusion has no characteristic restriction.
\end{saresult}
\begin{proof}
The image $A_0=eA$ is a quotient of the complete
Noetherian local ring $A$, hence is itself complete
Noetherian local. It embeds in $D$, and $D/A_0$
is elementwise killed by powers of its maximal
ideal $\mathfrak m_0$, since $D/I=C/A$ is.
Moreover $\mathfrak m_0D=D$.
Choose generators $f_1,\ldots,f_s$ of
$\mathfrak m_0$. Localization gives
$D_{f_i}=(A_0)_{f_i}$, and the opens
$\operatorname{Spec}(D_{f_i})$ cover
$\operatorname{Spec}(D)$ because $(f_1,\ldots,f_s)D=D$.
These canonical identifications agree on overlaps.
They prove
\[
 \operatorname{Spec}(D)
       \simeq\operatorname{Spec}(A_0)\setminus\{\mathfrak m_0\}.
                                                               \tag{B60.2}
\]

The opened primary theorem Stacks, Lemma 31.17.1
(tag 0BCR) \cite{StacksAffinePuncture}, says that the punctured spectrum of
a Noetherian local ring is affine exactly in
dimensions at most one. Its hypotheses apply
to $A_0$, and (B60.2) supplies affineness.
Since $D\ne0$, this spectrum is nonempty, so
$\dim A_0=1$.
Also $H^0_{\mathfrak m_0}(A_0)=0$: an element
annihilated by $\mathfrak m_0^j$ annihilates
$D=\mathfrak m_0^jD$, and therefore annihilates
its unit $e$, so is zero. Thus the one-dimensional
local ring $A_0$ is Cohen--Macaulay.

Choose a regular parameter $t\in\mathfrak m_0$.
Its basic open is the entire punctured spectrum,
so (B60.2) identifies $D=(A_0)[1/t]=Q(A_0)$.
The map $I\to A_0$ is injective because $ei=i$
for $i\in I$. Its image is an ideal, since
multiplication by $ea\in A_0$ is multiplication
by $a\in A$. The class of $e\in D/I$ is killed
by some maximal-ideal power, so $t^Ne\in I$
for some $N$. This is a nonzerodivisor of $A_0$.
If $I=A_0$, then $e\in I\subset A$; a nonzero
idempotent in the local ring $A$ is $1$.
It follows that $I=A$. This proves the final
assertion by contraposition.

The cited lemma was opened with its complete
elementary proof in Stacks, Section 31.17,
lines 20--34. No appeal to an unchecked local
duality or nonvanishing theorem is needed.
\end{proof}

\begin{saresult}{B.061}{Proved}{Complete local bases with zero divisors}
Let $(A,\mathfrak m)$ be any complete Noetherian
local ring with finite residue field, and let
$A\subset C$ be a unital extension with $C/A$
elementwise killed by powers of $\mathfrak m$.
The polynomial coaddition and comultiplication
tensors are uniquely determined if and only if
\[
 C=A\quad\text{or}\quad
 \bigl(A\text{ is a one-dimensional domain},
                         C=\operatorname{Frac}(A)\bigr).
                                                               \tag{B61.1}
\]
Every positive case has all multivariable tensor
isomorphisms. Thus an actual zero-divisor local
base has only the ordinary-polynomial positive case.
\end{saresult}
\begin{proof}
The monic null-polynomial argument of B.057
forces $A\to C$ to be epic if the polynomial
tensor map is injective. Apply B.056.
If $D=0$, then $C=A$. Otherwise B.060 gives
the complete one-dimensional Cohen--Macaulay
order $A_0=eA$ and $D=Q(A_0)$, with $I$
containing a regular parameter power.

Every selected nonreduced local factor of $D$
is excluded by B.059. To check its conductor
hypothesis explicitly, take the finite order
$O$ of that proof. Since $I[1/t]=D$ and
$e_0O$ is finite over $A_0$, some $t^Ne_0O$
lies in $I$. Replacing that parameter power
by a high power of the lifted uniformizer
gives $u^ce_0O\subset I\subset A$.
Thus B.059 constructs a nonzero tensor kernel
inside the original $\Omega(A,C)$, not merely
inside a polynomial module for a quotient.
Consequently all factors of $D$ are fields.

If $I$ is proper, B.060 makes it proper in
$A_0$ as well. Choose a field factor and its
finite complete discrete valuation normalization
$V_1$, with uniformizer $u$. There is a
conductor $u^ce_0V_1\subset I$, with $c\ge1$.
For $x=u^ce_0\in I$ and any $H\in\Omega(A,C)$,
the component of $H$ in $(A/I)[X]$ gives
$H(x)-H(0)\in I$. The congruence indicator
and $(c+1)$-fold difference proof of B.057
therefore applies unchanged: it gives a value
of valuation $c$ which a tensor expression
would put in $I^{c+1}$, where its valuation
is at least $c+1$. Thus proper $I$ is impossible.

It follows that $I=A$, so $e=1$ and
$C=D=Q(A)$; in particular $A=A_0$ is
one-dimensional Cohen--Macaulay. Multiple
field factors are excluded by the supported
finite-difference proof of B.050, using the
proper ideal $A\cap e_0Q(A)$ of one factor.
Nonreduced factors have already been excluded.
Hence $Q(A)$ is a field, and $A$ is a domain.
This proves necessity. Sufficiency is the
ordinary polynomial case or the complete
one-dimensional domain case of B.024 and
B.050, which also prove all tensor isomorphisms.
\end{proof}

\begin{saresult}{B.062}{Proved}{All Noetherian base rings and all coefficient extensions}
Let $R$ be any nonzero Noetherian commutative ring
and let $R\subset S$ be any unital commutative
extension. For each finite-residue maximal ideal
$\mathfrak m$ form
\[
 C_{\mathfrak m}
   =\{s\in S_{\mathfrak m}:\mathfrak m^jR_{\mathfrak m}s
          \subset R_{\mathfrak m}\text{ for some }j\}.
\]
The polynomial coaddition and comultiplication
tensors are uniquely determined exactly when,
at every such maximal ideal,
\[
\begin{gathered}
 C_{\mathfrak m}=R_{\mathfrak m},\quad\text{or}\\
 R_{\mathfrak m}\text{ is a one-dimensional domain},\\
 \widehat{R_{\mathfrak m}}\text{ is a domain},\qquad
 C_{\mathfrak m}=\operatorname{Frac}(R_{\mathfrak m}).     \tag{B62.1}
\end{gathered}
\]
The fraction field is required as the canonical
subring extending $R_{\mathfrak m}$ inside
$S_{\mathfrak m}$. No reducedness, domain,
torsion-free, or finite-normalization assumption
is imposed on the original inclusion. In every
positive case all multivariable tensor maps
are isomorphisms and the compatible plethory
is unique.
\end{saresult}
\begin{proof}
Use B.055 to localize, restrict to the coefficient
ring $C_{\mathfrak m}$, and complete. Infinite
residue field gives the ordinary polynomial ring.
At a finite-residue maximal ideal, B.061 gives
exactly the completed positive cases.
Faithful flatness descends the equality with the
local base. In the other case the completed
base is a domain, so its embedded local base
is a domain too, and dimension is one.
The fraction-field coefficient ring descends
by the same kernel and quotient argument as
in B.058. This proves precisely (B62.1).
Conversely those conditions give the positive
local cases. The tensor maps are isomorphisms
globally by testing kernels and cokernels at
all maximal ideals. Substitution, counits and
composition transport through these isomorphisms
as in B.009, giving the unique compatible plethory.
The zero ring, if allowed, has only its zero
unital extension and satisfies the trivial version.
\end{proof}

\begin{saresult}{B.063}{Proved}{Every valuation domain with arbitrary coefficient extensions}
Let $(V,\mathfrak m)$ be a nonfield valuation domain,
of arbitrary rank, and let $V\subset S$ be any
unital commutative extension. If the residue field
is infinite, every such extension has the unique
standard polynomial co-operations. If the residue
field is finite, the polynomial coaddition and
comultiplication tensors are uniquely determined
if and only if
\[
                         \operatorname{Ann}_S(\mathfrak m)=0.
                                                               \tag{B63.1}
\]
Every positive case has all multivariable tensor
isomorphisms. This removes the torsion-free
restriction from B.041; base torsion away from
the maximal-ideal socle is permitted.
\end{saresult}
\begin{proof}
Infinite residue field is the unit Vandermonde
case. Suppose the residue field is $k=\mathbb F_q$.
First let $\mathfrak m$ be nonprincipal. Its value
group has no least positive element. The elementary
ordered-group argument in B.041 gives, for every
$\gamma>0$ and integer $N$, a positive $\delta$
with $N\delta<\gamma$.
For a polynomial of degree at most $d\ge1$, choose
$a\in\mathfrak m$ of value $\delta$ and the grid
$0,a,a^2,\ldots,a^d$. Its Vandermonde determinant
has value $d(d+1)(d+2)\delta/6$. Cramer's rule
shows that this determinant sends every coefficient
of an integer-valued polynomial into $V$.
Taking $\delta$ small enough for each prescribed
$r\in\mathfrak m$ and using divisibility in a
valuation domain proves that every coefficient
belongs to
\[
                         C=(V:_S\mathfrak m).
\]
This is a ring because $\mathfrak m^2=\mathfrak m$:
factor an element of positive value into two
elements of smaller positive value.

Localizing $S$ at $V\setminus\{0\}$ maps each
$c\in C$ into the fraction field $K$, since
$rc\in V$ for any nonzero $r\in\mathfrak m$.
That rational image must lie in $V$: a negative
value is contradicted by multiplying by an element
of sufficiently small positive value. Subtracting
this element of $V$ leaves an element annihilated
by $\mathfrak m$, because its products with
$\mathfrak m$ are both in $V$ and zero after
localization. Thus, as $V$-modules,
\[
                    C=V\oplus T,\qquad
                    T=\operatorname{Ann}_S(\mathfrak m).
\]
Writing $Q(X)=X^q-X$, monic division over the
$k$-vector space $T$ gives
\[
                    \Omega(V,S)=V[X]\oplus Q(X)T[X].     \tag{B63.2}
\]
Indeed the second summand consists exactly of
the $T$-coefficient polynomials vanishing on $k$.
If $0\ne z\in T$, the tensor
\[
                  zQ\otimes Q-Q\otimes zQ
\]
has zero polynomial image and is nonzero: its
two terms lie in distinct direct summands of
the tensor square in (B63.2), and $Q$ is monic
in the free module $V[X]$. Hence uniqueness
fails. If $T=0$, the polynomial module is
$V[X]$, proving this case.

Now suppose $\mathfrak m=\pi V$. Its least positive
value is $v(\pi)$. Each quotient $V/\pi^nV$ has
$q^n$ elements. Choosing any finite grid with
distinct residues modulo a sufficiently high
$\pi$-power makes its Vandermonde determinant
a unit times a power of $\pi$. Thus every
integer-valued coefficient lies in
\[
             C=\{s\in S:\pi^ns\in V\text{ for some }n\}.
                                                               \tag{B63.3}
\]
If $\operatorname{Ann}_S(\pi)=0$, this ring embeds
in $V[1/\pi]\subset K$. Every element of
$V[1/\pi]\setminus V$ has value a negative
integer multiple of $v(\pi)$, so an overring
containing one such element contains $\pi^{-1}$.
Consequently $C=V$ or $C=V[1/\pi]$. The former
gives ordinary polynomials; B.041 gives all
tensor isomorphisms in the latter case.

For necessity, put
\[
                         P=\bigcap_{n\ge0}\pi^nV.
\]
An element outside $P$ is a unit times a
nonnegative power of $\pi$: its value is bounded
by a finite multiple of the least positive
value, so is exactly such an integer multiple.
It follows that $P$ is prime and $V/P$ is a
discrete valuation ring. The ideal $P$ remains
an ideal of $C$, and $PC=P$. To see this, if
$\pi^Nc=a\in V$ and $p\in P$, write
$p=\pi^{N+j}b_j$ for every $j$. Then
$pc=\pi^jb_ja\in\pi^jV$ for every $j$, as required.

Write $\bar V=V/P$, $\bar C=C/P$. Reduction of
coefficients gives a surjection
\[
 \Omega(V,C)\longrightarrow\Omega(\bar V,\bar C)
 \quad\text{with kernel }P[X]=P\Omega(V,C).              \tag{B63.4}
\]
Indeed lift the finitely many coefficients of
an integer-valued polynomial over $\bar C$.
Its values on $V$ lie in the preimage of
$\bar V$, which is $V$; hence the lift is
integer-valued. The kernel and last equality
follow from $PC=P$.

Tensor injectivity descends through (B63.4),
despite the possible nonflatness of $V/P$.
Given a tensor in the reduced kernel, lift it
to $z\in\Omega(V,C)^{\otimes2}$. Its polynomial
image belongs to $P[X,Y]$. That polynomial
has a tensor preimage formed from ordinary
monomials with coefficients in $P$, and this
preimage reduces to zero. Subtract it from $z$.
Injectivity upstairs forces the resulting tensor,
and thus the original reduced tensor, to vanish.
The two existence conditions descend as well,
since every reduced polynomial lifts.

Apply B.058 to the discrete valuation ring
$\bar V$; every element of $\bar C/\bar V$ is
killed by a power of $\pi$. It follows that
$\bar C=\bar V$ or $\bar C=\operatorname{Frac}(\bar V)$.
The first equality forces $C=V$. In the second
case lift an inverse of $\pi$: one has
$\pi c=1+p$ for some $p\in P$. Since $1+p$
is a unit of $V$, $\pi$ is a unit of $C$,
and (B63.3) gives $C=V[1/\pi]$.
Both possibilities have no $\pi$-torsion.
But every element of $\operatorname{Ann}_S(\pi)$
belongs to $C$, proving its vanishing.
Since $\operatorname{Ann}_S(\mathfrak m)=
\operatorname{Ann}_S(\pi)$, this proves necessity.

The multivariable coefficient assertion used in
the positive cases follows from rectangular
Vandermonde grids, or by substituting widely
separated powers of one variable to keep all
monomials distinct. Thus the positive tensor
isomorphisms compute the full multivariable
integer-valued modules, including the original
coefficient extension $S$.
\end{proof}

\begin{saresult}{B.064}{Proved}{A universal coefficient-ring obstruction without Noetherian hypotheses}
For any unital inclusion $R\subset S$ and $d\ge1$,
let $D_d(R)$ be the ideal generated by all
Vandermonde determinants on $d+1$ elements of $R$,
and put
\[
 C_d=(R:_S D_d(R)),\qquad C=\bigcup_{d\ge1}C_d.           \tag{B64.1}
\]
Then $C$ is a subring of $S$. It is exactly the
set of coefficients that occur in integer-valued
polynomials, and also exactly the set of their
possible leading coefficients, with zero included. Every multivariable
integer-valued polynomial has its coefficients in
$C$. If the second polynomial tensor map is
injective, the inclusion $R\to C$ is a ring
epimorphism. If $C=R$, all tensor maps are the
ordinary polynomial isomorphisms.
\end{saresult}
\begin{proof}
The ideals $D_d$ decrease with $d$, since a
larger Vandermonde product is a multiple of
one on a subset of its nodes. Thus the $C_d$
increase, and $C_1=R$ because the nodes $0,1$
have determinant one.
Cramer's rule says that every coefficient of
an integer-valued polynomial of degree at most
$d$ lies in $C_d$.

There is an explicit monic polynomial $P_d$
whose every value belongs to $D_d$. For $d=1$
take $P_1=1$. For $d\ge2$, take the Vandermonde
product on the polynomial nodes
\[
                     0,\ 1,\ X,\ X^2,\ldots,X^{d-1}.
\]
Every factor with positive degree is monic,
and the only degree-zero factor is $1$; hence
$P_d$ is monic over every base ring.
For any $c\in C_d$, the polynomial $cP_d$
takes $R$ into $R$ and has leading coefficient
$c$. This proves both coefficient characterizations.
Sums of leading coefficients are obtained after
multiplying the corresponding polynomials by
powers of $X$ to equalize their degrees; products
come from multiplication of polynomials.
Therefore their set $C$ is a subring.

For a multivariable polynomial, substitute
$X_i=T^{N_i}$ with sufficiently separated
positive exponents so that its finitely many
monomials remain distinct. The resulting
one-variable polynomial is integer-valued on
$R$, and has exactly the same coefficients.
Thus all coefficients lie in $C$.

For $c\in C_d$, both $P_d$ and $cP_d$ belong
to $\Omega(R,S)$. The tensor
$cP_d\otimes P_d-P_d\otimes cP_d$ has zero
polynomial image. If the tensor map is
injective it vanishes. Extract the top
coefficient in each factor, now as maps
$\Omega(R,S)\to C$, to get
$c\otimes1=1\otimes c$ in $C\otimes_R C$.
This for every $c\in C$ makes multiplication
$C\otimes_R C\to C$ an isomorphism.
Finally $C=R$ gives $\Omega_n=R[X_1,\ldots,X_n]$
by the coefficient statement, proving the
last assertion.
\end{proof}

\begin{saresult}{B.065}{Proved}{Every absolutely flat base with arbitrary coefficients}
Let $R$ be an absolutely flat commutative ring,
meaning that every $R$-module is flat, and let
$R\subset S$ be any unital commutative extension.
With the explicit Vandermonde ideals of B.064,
the following are equivalent:
\[
\begin{gathered}
 \text{the polynomial co-operations are uniquely determined};\\
 \theta_2\text{ is injective};\qquad
 (R:_S D_d(R))=R\text{ for every }d;\\
 \operatorname{Ann}_{S/R}(D_d(R))=0\text{ for every }d;\qquad
 \Omega(R,S)=R[X].                                       \tag{B65.1}
\end{gathered}
\]
Every positive case has all polynomial tensor
isomorphisms. In particular, for any Boolean
ring, finite or infinite, the positive
extensions are exactly $S=R$.
\end{saresult}
\begin{proof}
By B.064, injectivity makes $R\to C$ an
injective ring epimorphism. Such an epimorphism
over an absolutely flat base is an isomorphism.
Indeed $C/R$ is flat, so tensoring
$0\to R\to C\to C/R\to0$ with any module
$M$ keeps $M\to C\otimes_R M$ injective.
But the epimorphism condition gives
$C\otimes_R(C/R)=0$. Take $M=C/R$ to conclude
$C/R=0$.
Thus injectivity implies $C=R$. B.064 shows
that this is equivalent to the coefficient
conditions and to $\Omega=R[X]$, and gives
all positive tensor isomorphisms. The
annihilator condition is the same assertion
in the quotient module $S/R$.

For a Boolean ring, $x^2=x$ for every $x$.
Every prime quotient is the field with two
elements, so every three-node Vandermonde
determinant vanishes in every prime quotient.
The ring is reduced, hence $D_2(R)=0$.
Consequently $C_2=S$, and the criterion is
exactly $S=R$. Boolean rings are absolutely
flat: every finitely generated ideal is
generated by an idempotent, and the usual
finite-equation criterion for flatness then
shows that every module is flat. More
explicitly, for any finitely generated ideal
$I=eR$, the map $I\otimes_R M\to M$ identifies
with the inclusion $eM\hookrightarrow M$;
the ideal criterion proves flatness, as in the opened
Stacks Lemma 10.39.5, tag 00HD. The equivalent
reduced zero-dimensional characterization of absolute
flatness is Stacks Lemma 15.106.5, tag 092F
\cite{StacksAbsoluteFlat}.
\end{proof}

\begin{saresult}{B.066}{Proved}{Complete ideal criterion for faithful polynomial evaluation}
For every unital inclusion $R\subset S$, let $D_d(R)$
be the explicit Vandermonde ideal of B.064. Evaluation
\[
 \Omega(R,S)\longrightarrow\operatorname{App}(R,R)
\]
is injective if and only if
\[
                  \operatorname{Ann}_S(D_d(R))=0
                  \quad\text{for every }d\ge1.           \tag{B66.1}
\]
There are no finiteness, reducedness, domain or
torsion-free hypotheses. Thus B.2(d), for the
polynomial analyser of any ring inclusion, has an
explicit ring-ideal criterion rather than only a
degreewise evaluation-matrix restatement.
\end{saresult}
\begin{proof}
If a polynomial $F\in S[X]$ of degree at most $d$
vanishes on $R$, Cramer's rule on any $d+1$ nodes
shows that their Vandermonde determinant annihilates
every coefficient of $F$. Hence each coefficient
belongs to $\operatorname{Ann}_S(D_d(R))$.
Under (B66.1) all coefficients are zero; a constant
null polynomial is already zero by evaluation.

Conversely, let $0\ne a\in\operatorname{Ann}_S(D_d(R))$.
The explicitly constructed monic polynomial $P_d$
of B.064 has $P_d(r)\in D_d(R)$ for every $r\in R$.
Thus $aP_d$ is a null polynomial, belongs to
$\Omega(R,S)$, and is nonzero because its top
coefficient is $a$. Evaluation is not injective.
This proves both directions, including arbitrary
zero divisors in the coefficient ring.
\end{proof}

\begin{saresult}{B.067}{Proved}{A positive non-Noetherian local inclusion with a nonflat analyser}
There is an explicit proper inclusion $R\subset S$
of zero-dimensional local $\mathbb F_2$-algebras
such that every nilpotent squares to zero,
$R$ is non-Noetherian, and every map
\[
 \Omega(R,S)^{\otimes_R n}\longrightarrow\Omega_n(R,S)
\]
is an isomorphism. Nevertheless $\Omega(R,S)$
is not flat over $R$, and the coefficient ring
of B.064 is the proper ring $S$.
Thus flatness of the analyser or of its coefficient
ring cannot be imposed as a necessary condition
in the remaining non-Noetherian classification.
\end{saresult}
\begin{proof}
Index an infinite forest of full binary trees by
$v=(j,w)$, where $j\ge1$ is a root index and
$w$ is a finite binary word. Write $v0,v1$ for
the children. Define
\[
 S=\mathbb F_2[x_v,y_v:v]/
       (x_v^2,y_v^2,\ y_v-x_vy_{v0}y_{v1}:v).             \tag{B67.1}
\]
This presentation has the following explicit
injective direct-limit description. At level
$n$, take the squarefree commutative polynomial
algebra on the variables $x_v$ of depth less
than $n$ and $y_v$ of depth exactly $n$.
Pass to level $n+1$ by replacing each boundary
variable $y_v$ by the product
$x_vy_{v0}y_{v1}$, retaining the old $x$-variables.
Different boundary variables use disjoint new
variable sets. Hence different squarefree
monomials map to different squarefree monomials,
and the map is injective. Its direct limit is
precisely (B67.1); every relation has been
imposed at a finite level, and every finite
polynomial expression is represented at one.
In particular every $y_v$ is nonzero.

Give $x_v$ degree $-1$ and $y_v$ degree $1$.
All transition maps and relations are homogeneous.
Set
\[
                         R=\bigoplus_{d\le0}S_d.          \tag{B67.2}
\]
This is a subring containing $1$. It is proper,
since each $y_v$ is a nonzero element of degree
$1$. Both rings have residue field $\mathbb F_2$.
Indeed their augmentation ideals consist of
elements whose square is zero: in characteristic
two the square of a finite sum is the sum of
the squares, and every nonconstant squarefree
monomial has square zero. Every element with
constant coefficient one is a unit, with
inverse itself. Thus each augmentation ideal
is the unique prime and maximal ideal.
Write $N$ for that ideal in $R$.

For each $k\ge1$, a product of $k$ nonzero
degree-zero monomials $x_vy_{v0}$ from distinct
tree components is nonzero and belongs to
$N^k$. Hence $N$ is not nilpotent. A finitely
generated ideal generated by square-zero
elements is nilpotent, by the pigeonhole
principle. Therefore $N$ is not finitely
generated, and $R$ is non-Noetherian.

The inclusion is a ring epimorphism. Both
$x_vy_{v0}$ and $x_vy_{v1}$ have degree zero
and belong to $R$, so in $S\otimes_R S$ one has
\[
\begin{aligned}
 y_v\otimes1
 &=y_{v0}\otimes x_vy_{v1}\\
 &=x_vy_{v0}\otimes y_{v1}
 =1\otimes y_v.
\end{aligned}                                             \tag{B67.3}
\]
The same identity is automatic for $x_v\in R$.
Since these elements generate $S$, the two
maps $S\to S\otimes_R S$ agree, proving
$S\otimes_R S\simeq S$.

We need a precise conductor calculation:
\[
                        (R:_S N^k)=R\qquad(k\ge1).        \tag{B67.4}
\]
If $c\notin R$, its positive-degree part is
nonzero and involves only finitely many tree
components. Choose $k$ fresh components and
let $a$ be the product of the degree-zero
monomials $x_vy_{v0}$ in them. Then $a\in N^k$,
and multiplying the positive-degree part of
$c$ by $a$ does not kill it. To check this
last assertion without a domain assumption,
each finite-stage algebra is the tensor
product over $\mathbb F_2$ of its separate
tree-component algebras. A product involving
disjoint components is a pure tensor of
nonzero vectors, hence is nonzero; injectivity
of the transition maps preserves this fact.
Thus $ac$ still has a nonzero positive-degree
part, so is outside $R$. This proves (B67.4).

Put $Q(X)=X^4+X^2=X^2(X-1)^2$. For every
$r=a+n\in R$, with $a\in\mathbb F_2$ and
$n\in N$, one has $n^2=0$, so $Q(r)=0$.
Monic division by $Q$ reduces an arbitrary
$S$-polynomial to degree less than four.
For such a remainder $H$, the identity
\[
                    H(a+n)=H(a)+H'(a)n
\]
shows that it preserves $R$ exactly when
$H(0),H(1),H'(0),H'(1)$ belong to $R$,
using (B67.4) with $k=1$.
Hermite interpolation at the two points $0,1$
is invertible over $\mathbb F_2$: it is the
Chinese remainder isomorphism for the
comaximal ideals $(X^2)$ and $((X-1)^2)$.
Therefore those four conditions say exactly
that all four coefficients of $H$ lie in $R$.
We have proved the split module formula
\[
                  \Omega(R,S)=R[X]_{<4}\oplus Q(X)S[X].  \tag{B67.5}
\]

The multivariable version is equally explicit.
Use in each variable the monic basis
\[
 b_d(X)=X^d\ (d<4),\qquad
 b_d(X)=Q(X)X^{d-4}\ (d\ge4).
\]
A tensor product of these bases is an
$S$-basis of $S[X_1,\ldots,X_n]$.
Every basis term with at least one index
$d_i\ge4$ vanishes on $R^n$, so its coefficient
may be arbitrary in $S$. The remaining
remainder has multidegree less than four in
each variable. For each $a\in\{0,1\}^n$,
take finite differences in any subset $J$
of the variables with increments $n_j\in N$.
Since each $n_j^2=0$, the difference is
\[
                      (\partial_J H)(a)\prod_{j\in J}n_j.
\]
If $H$ preserves $R$, it follows that
$(\partial_JH)(a)N^{|J|}\subset R$: the products of the arbitrary increments $n_j$ generate $N^{|J|}$ as an $R$-ideal.
Equation (B67.4) gives $(\partial_JH)(a)\in R$;
the case $J=\varnothing$ is direct evaluation.
Tensoring the four-coordinate Hermite
isomorphism proves that every coefficient
of this low-degree remainder lies in $R$.
Thus $\Omega_n$ has coefficient module $R$
on the all-low-index basis terms and $S$
on all other basis terms.

The tensor power of (B67.5) has exactly these
same coefficient modules: a product of $R$
summands gives $R$, while a product containing
at least one $S$ summand gives $S$, by (B67.3).
The polynomial tensor map identifies these
summands coefficient by coefficient. Hence
every $\theta_n$ is an isomorphism.
For additional explicitness, the two identities
\[
\begin{aligned}
 Q(X+Y)&=Q(X)+Q(Y),\\
 Q(XY)&=Q(X)Q(Y)+Q(X)Y^2+X^2Q(Y)
\end{aligned}
\]
give the co-operations on $cQ$ for every
$c\in S$; multiplication and ordinary
polynomials give them on all of (B67.5).
Composition and every compatibility identity
follow by the tensor isomorphisms as in B.009.

Finally $S$ is not flat over $R$. If it were,
the local inclusion, whose residue fiber is
nonzero, would be faithfully flat; equivalently
its map of one-point spectra is surjective.
A faithfully flat epimorphism is an isomorphism,
contrary to $y_v\notin R$. The two facts are
the opened Stacks Lemmas 10.39.16 (tag 00HQ)
and 10.107.7 (tag 04VU). Since (B67.5)
contains a split copy of $S$, the analyser
is not flat either. Its coefficient ring is
all of $S$, since $cQ$ has leading coefficient
$c$ for every $c$.

The attached exact program \texttt{check\_B\_forest.py}
checks $1890$ transition monomials and $464100$
monomial products, the complete $256$-relation
tensor presentation of one finite branching level,
Hermite ranks $4,16,64$, and the two displayed
polynomial identities. The finite level has tensor
dimension $20$: the parent-variable difference
vanishes while a leaf-variable difference remains
nonzero. Thus the certificate explicitly does
not claim that a finite truncation is epic.
The infinite assertions follow from the proof above.

The epic tensor shuffle is analogous to the
Lazard example reproduced in Stacks Section
110.36, tag 06RH, which was opened. The
bounded-square branching construction, its
injective direct limit, and all integer-valued
and tensor assertions have been proved here;
\cite{StacksEpicExamples}; no priority claim is made and no omitted
detail of that source is used.
\end{proof}

\begin{saresult}{B.068}{Proved under stated hypotheses}{Regular Vandermonde ideals and arbitrary coefficient extensions}
Let $R$ be a nonzero commutative ring, let $Q$ be its total quotient ring,
and let $D_d(R)$ be the Vandermonde ideals of B.064. Suppose that, for every
$d\ge1$, the ideal $D_d(R)$ contains a nonzerodivisor of $R$, and
\[
                         (R:_Q D_d(R))=R.                 \tag{B68.1}
\]
Then for every unital inclusion $R\subset S$ the following are equivalent:
\[
\begin{gathered}
 \text{the polynomial co-operations exist uniquely};\\
 \theta_2\text{ is injective};\\
 \Omega(R,S)\longrightarrow\operatorname{App}(R,R)
       \text{ is injective};\\
 \operatorname{Ann}_S(D_d(R))=0\quad\text{for every }d;\\
 \Omega(R,S)=R[X].                                         \tag{B68.2}
\end{gathered}
\]
Every positive case has all multivariable polynomial tensor isomorphisms.
The hypotheses hold for every infinite domain with
$\operatorname{Int}(R,\operatorname{Frac}R)=R[X]$.
They also hold for $R=A[Z,W]$ over every nonzero commutative ring $A$,
without finiteness or reducedness assumptions. More generally they hold
for $R=A[Z]$ whenever each $D_d(A)$ contains a nonzerodivisor; in particular,
this includes every infinite domain $A$.
\end{saresult}
\begin{proof}
Write $U$ for the multiplicative set of nonzerodivisors in $R$.
The inclusion $R\subset S$ induces an inclusion $Q\subset U^{-1}S$:
if an element of $R$ becomes zero, a nonzerodivisor of $R$ annihilates it,
which is impossible unless it was zero.
For $F=\sum_{i=0}^d a_iX^i\in\Omega(R,S)$, Cramer's rule gives
$D_d(R)a_i\subset R$ for each coefficient. Choose a nonzerodivisor
$v\in D_d(R)$; it may be a finite linear combination of Vandermonde
products, since the Cramer conclusion holds on their whole ideal. The image of $a_i$ in $U^{-1}S$ equals
$(va_i)/v\in Q$. The rationalized polynomial is therefore in
$\operatorname{Int}(R,Q)$. By B.064 and (B68.1),
$\operatorname{Int}(R,Q)=R[X]$.
Rationalization consequently gives an $R$-linear retraction
\[
 \rho:\Omega(R,S)\longrightarrow R[X],\qquad
 \Omega(R,S)=R[X]\oplus\mathcal N,                         \tag{B68.3}
\]
where $\mathcal N$ is exactly the module of null polynomials on $R$.
Indeed an element in the kernel has each coefficient killed by some
member of $U$, hence all its coefficients by a common such member.
Its values belong to $R$, so regularity makes all values zero.
Conversely a null polynomial has $D_d(R)$ annihilating every coefficient,
and the chosen regular element makes its rationalization zero.

Suppose $\mathcal N\ne0$. Select a nonzero coefficient $a$ of a null
polynomial of degree at most $d$. Then $D_d(R)a=0$. The monic polynomial
$P_d$ of B.064 gives a nonzero null polynomial $aP_d$. The tensor
\[
             aP_d\otimes P_d-P_d\otimes aP_d               \tag{B68.4}
\]
has zero polynomial image. It is nonzero: the two terms lie in distinct
summands $\mathcal N\otimes_R R[X]$ and $R[X]\otimes_R\mathcal N$ of the
tensor square of (B68.3), and extracting the top coefficient of the
ordinary polynomial $P_d$ sends the first term to $aP_d\ne0$.
Thus injectivity forces $\mathcal N=0$. The other implications now follow
from B.066, from (B68.3), and from the ordinary polynomial tensor
isomorphisms. This proves (B68.2).

For an infinite domain, each $D_d(R)$ is nonzero by choosing distinct
nodes. B.064 identifies (B68.1) with the stated ordinary
integer-valued-polynomial condition, proving that assertion.

For the polynomial examples, first assume every $D_d(A)$ contains a
nonzerodivisor of $A$, and put $R=A[Z]$. Such an element remains regular
in $R$, so it also belongs to $D_d(R)$ and proves the first hypothesis.
If $F\in\operatorname{Int}(R,Q(R))$ has degree $d$, evaluating on nodes
in $A$ and taking a finite $A$-linear combination of the corresponding
Cramer identities gives $v a_i\in A[Z]$ for some regular $v\in D_d(A)$.
Thus all coefficients belong to
\[
                      A[1/v][Z]\subset Q(R).
\]
Choose an integer $N$ strictly larger than the $Z$-degree of every
coefficient polynomial $a_i$. The degree intervals occupied by
$a_i(Z)Z^{Ni}$ are pairwise disjoint. Since $F(Z^N)\in A[Z]$, comparison
of these coefficients in the injective polynomial-ring inclusion
$A[Z]\subset A[1/v][Z]$ proves that every coefficient of every $a_i$
lies in $A$. Hence $F\in A[Z][X]$. B.064 gives (B68.1).

Finally, for arbitrary nonzero $A$, every $D_d(A[W])$ contains a monic
polynomial in $W$: for $d\ge2$ take the determinant on
$0,1,W,W^2,\ldots,W^{d-1}$, and for $d=1$ take $1$.
A monic polynomial over any nonzero commutative ring is a
nonzerodivisor. Applying the preceding argument with $A[W]$ in place
of $A$ proves the two-variable assertion. No localization of the
integer-valued module, flatness of $S$, or cancellation by zero divisors
has been used.
\end{proof}

\begin{saresult}{B.069}{Proved under stated hypotheses}{Power-generated maximal ideals and surjective Frobenius}
Let $R$ be a nonzero commutative ring. Suppose that, at each maximal ideal
$\mathfrak m$ with finite residue field, the local maximal ideal
$\mathfrak n=\mathfrak mR_{\mathfrak m}$ satisfies
\[
             \mathfrak n=(a^N:a\in\mathfrak n)
             \quad\text{for every integer }N\ge1.          \tag{B69.1}
\]
Then for every unital inclusion $R\subset S$ the following conditions
are equivalent:
\[
\begin{gathered}
 \text{the polynomial co-operations exist uniquely};\\
 \theta_2\text{ is injective};\qquad C=R;\\
 (R:_S D_d(R))=R\quad\text{for every }d\ge1;\\
 \operatorname{Ann}_{S/R}(D_d(R))=0\quad\text{for every }d\ge1;\\
 \Omega(R,S)=R[X],                                        \tag{B69.2}
\end{gathered}
\]
where $C$ is the actual coefficient ring of B.064. Every positive case
has all multivariable polynomial tensor isomorphisms.

In particular this class includes every commutative $\mathbb F_p$-algebra
with surjective Frobenius, with arbitrary nilpotents and without
Noetherian assumptions. Surjectivity, rather than injectivity, of
Frobenius is what this assertion requires.
\end{saresult}
\begin{proof}
We first compute the determinant ideals locally, and do not assume
localization of integer-valued polynomial modules. For any multiplicative
set $U$, one has
\[
                   D_d(U^{-1}R)=U^{-1}D_d(R).              \tag{B69.3}
\]
The forward inclusion follows by giving a finite tuple of nodes a common
denominator $u\in U$: its Vandermonde determinant is the determinant of
nodes in $R$, divided by $u^{d(d+1)/2}$. The reverse inclusion is immediate.

Let $(A,\mathfrak n)$ be a local ring satisfying (B69.1), with residue
field of cardinality $q$. When $d<q$, lifts of $d+1$ distinct residue
elements have unit Vandermonde determinant, so $D_d(A)=A$. When $d\ge q$,
any $d+1$ residues repeat, so $D_d(A)\subset\mathfrak n$.
For the reverse inclusion use the polynomial $P_d$ of B.064. At any
$a\in\mathfrak n$ its explicit factors give
\[
 P_d(a)=a^{e_d}u_a,\qquad
 e_d=\frac{d(d-1)(d+1)}6,\qquad u_a\in A^\times.            \tag{B69.4}
\]
Indeed every difference between powers is a power of $a$ times a factor
$a^j-1$, which is a unit. The exponent is the sum of the exponents from
the differences with zero and the smaller exponents from pairs of
positive powers. Thus $a^{e_d}\in D_d(A)$ for every $a\in\mathfrak n$.
Equation (B69.1), with $N=e_d\ge1$, proves
\[
 D_d(A)=
 \begin{cases}A,&d<q,\\ \mathfrak n,&d\ge q.\end{cases}     \tag{B69.5}
\]
Also (B69.1) with $N=2$ gives $\mathfrak n^2=\mathfrak n$.
At an infinite-residue local ring all $D_d(A)$ equal $A$, by distinct
residue nodes, without any hypothesis on its maximal ideal.

Suppose now that $\theta_2$ is injective. B.064 makes $R\to C$ an
injective ring epimorphism. Localizing this inclusion and the epimorphism
at any maximal ideal $\mathfrak m$ gives $A=R_{\mathfrak m}\subset
B=C_{\mathfrak m}$. For each $b\in B$, some $D_d(A)$ multiplies $b$ into
$A$, because $b$ comes from one $C_d$ and (B69.3) applies. If the residue
field is infinite this already says $B=A$.

In the finite-residue case, (B69.5) says
$\mathfrak nB\subset A$. As an $A$-ideal it contains $\mathfrak n$, so it
equals $\mathfrak n$ or $A$. The latter is impossible. Otherwise write
$1=\sum_{i=1}^s a_i b_i$ with $a_i\in\mathfrak n$, $b_i\in B$. For every
$a\in\mathfrak n$ the coefficients $ab_i$ lie in $A$, so
$a=\sum_i a_i(ab_i)$ proves that $\mathfrak n$ is finitely generated.
But $\mathfrak n^2=\mathfrak n$ and Nakayama's elementary finite-generator
lemma then give $\mathfrak n=0$, contradicting $1\in\mathfrak nB$.
Hence
\[
                         \mathfrak nB=\mathfrak n.         \tag{B69.6}
\]
Base change of the epimorphism by $A\to A/\mathfrak n$ is therefore an
injective epimorphism from the field $k=A/\mathfrak n$ into
$B/\mathfrak n$. Such an inclusion is an isomorphism: if $x$ in a
$k$-algebra is linearly independent of $1$, then
$x\otimes1-1\otimes x$ is nonzero over $k$, contradicting the epimorphism
condition. Thus $B/\mathfrak n=k$ and $B=A+\mathfrak n=A$.

We have proved $C_{\mathfrak m}=R_{\mathfrak m}$ at every maximal ideal.
The module $C/R$ therefore vanishes, with no finite-generation hypothesis:
a nonzero element has a proper annihilator contained in a maximal ideal,
and survives localization there. Hence $C=R$. B.064 proves all remaining
equivalences in (B69.2) and the positive tensor isomorphisms.

Finally suppose $R$ is an $\mathbb F_p$-algebra with surjective Frobenius.
This property survives localization. For $a$ in a local maximal ideal
and any $N$, choose $e$ with $p^e\ge N$ and a $p^e$-th root $b$ of $a$.
Then $b$ also lies in the maximal ideal, since it is prime, and
$a=b^{p^e}=b^N b^{p^e-N}$. Every element of that maximal ideal lies in the
ideal generated by $N$-th powers of its elements, proving (B69.1).
The reverse containment is automatic. Nilpotents cause no change in
this argument.
\end{proof}

\begin{saresult}{B.070}{Proved}{Every binomial base ring and every coefficient extension}
Let $R$ be a nonzero binomial ring: its additive group is
$\mathbb Z$-torsion-free, and $\binom r n\in R$ for every $r\in R$ and
$n\ge0$, with the fraction interpreted in $\mathbb Q\otimes R$.
For every unital commutative inclusion $R\subset S$, the polynomial
coaddition and comultiplication tensors are uniquely determined if and
only if the additive group of $S$ is $\mathbb Z$-torsion-free.
In every positive case all multivariable polynomial tensor maps are
isomorphisms. No Noetherian, domain or reducedness hypothesis is required
on $R$, and $S$ may have torsion as an $R$-module.

This includes $R=\operatorname{Int}(\mathbb Z)$, and also the nonreduced
non-Noetherian binomial ring $R=\mathbb Z\oplus\mathbb Q\epsilon$ with
$\epsilon^2=0$.
\end{saresult}
\begin{proof}
Write $(X)_n=X(X-1)\cdots(X-n+1)$, with $(X)_0=1$.
For every polynomial $F\in S[X]$ of degree at most $d$, and every $r\in R$,
the finite Newton identity holds in $S$:
\[
               F(r)=\sum_{n=0}^d \Delta^nF(0)\binom r n.  \tag{B70.1}
\]
This does not cancel integers in $S$. For each monomial $X^j$, the
identity $r^j=\sum_n n!\,{j\brace n}\binom r n$ holds in the
torsion-free ring $R$, where ${j\brace n}$ are the integer Stirling
coefficients. Multiplying these identities by the coefficients of $F$
and adding proves (B70.1) even if $S$ has integer torsion.
Consequently $F$ preserves $R$ if and only if its finitely many values
at $0,1,\ldots,d$ lie in $R$.

The polynomials $(X)_n$ form a monic $S$-basis, and
$\Delta^i(X)_n(0)=n!$ for $i=n$, and zero otherwise. Hence there is an
exact, split coefficient description
\[
 \Omega(R,S)=\bigoplus_{n\ge0}(R:_S n!)(X)_n.              \tag{B70.2}
\]
For several variables, the identical argument on a rectangular integer
grid gives coefficient module
$(R:_S\prod_i n_i!)$ on the basis $\prod_i(X_i)_{n_i}$.

Every localization of $R$ is again binomial. Here is the needed direct
localization argument. If a polynomial $G$ in a coefficient extension
preserves $R$, then it preserves $R[1/s]$ with values in the localized
base: induct on its degree $d$, using that
$s^dG(X)-G(sX)$ has smaller degree and still preserves $R$; evaluation
at $r/s$ gives $s^dG(r/s)\in R[1/s]$. Apply this to each
$G(X)=\binom X n$ inside $\mathbb Q\otimes R$ and allow arbitrary
localization denominators. Integer torsion-freeness is also preserved
by localization. Finally the degree-bounded preservation condition in
(B70.1) is the kernel of the finite evaluation map
\[
          S^{d+1}\longrightarrow (S/R)^{d+1}.
\]
Flatness of localization therefore proves exact localization of
$\Omega$ for every $S$, including integer torsion, and the same statement
in every number of variables. No Noetherian argument is needed.

We now analyze a local binomial ring $(A,\mathfrak n)$.
If its residue field has characteristic zero, every nonzero integer is
a unit in $A$, and ordinary integer Vandermonde determinants show that
$\Omega(A,S)=A[X]$ for every extension $S$.
If the residue characteristic is $p$, one has
$a^p-a\in pA$ for every $a\in A$: the polynomial $(X)_p$ is congruent to
$X^p-X$ modulo $p$, and $(a)_p=p!\binom a p\in pA$.
Thus every element of the local ring $A/pA$ satisfies $x^p=x$.
An element $x$ of its maximal ideal satisfies
$x(1-x^{p-1})=0$, where the second factor is a unit, so $x=0$.
The residue field is consequently $\mathbb F_p$, and
\[
                           \mathfrak n=pA.                \tag{B70.3}
\]
The element $p$ is a nonzerodivisor of $A$.
By (B70.2), all occurring coefficients lie in the ring
\[
                   C=\{s\in S:p^Ns\in A
                                  \text{ for some }N\}.  \tag{B70.4}
\]
Conversely every such element occurs, by taking $n$ with $p^N\mid n!$;
the other integer prime factors are units in $A$.
The equality $\Omega(A,S)=\Omega(A,C)$ follows.

Suppose first that $p$ is a nonzerodivisor in $S$. Then $C$ embeds into
$A[1/p]$. This overring is either $A$ or all of $A[1/p]$.
Indeed if $c=a/p^N\notin A$, then $a\notin p^NA$; choose the largest
$j<N$ for which $a\in p^jA$ and write $a=p^ju$ with $u$ a unit, using
(B70.3). Thus $c=p^{j-N}u$, so $p^{-1}\in C$.
If $C=A$, the polynomial tensor maps are the ordinary ones.
If $C=A[1/p]$, equation (B70.2) gives the free $A$-basis
$\binom X n$, and its multivariable version gives the product of these
bases. The polynomial tensor maps are isomorphisms in every arity.

For necessity of the absence of $p$-torsion, put
\[
                          P=\bigcap_{j\ge0}p^jA.
\]
The quotient $\bar A=A/P$ is a discrete valuation ring with uniformizer
$p$. To verify this without a Noetherian assumption, $p$ stays regular
modulo $P$: from $pa\in P$ cancel the regular element $p$ in each
$pa=p^{j+1}b_j$ to obtain $a\in P$. Every nonzero element of $A/P$ is a
unit times a finite power of $p$, by choosing the largest $j$ with
$a\in p^jA$. Products are therefore nonzero, and every nonzero ideal is
generated by an element having the smallest such exponent. This proves
the domain, discrete valuation and Noetherian assertions directly.

The ideal $P$ is an ideal of $C$, with $PC=P$. If $p^Nc=a\in A$ and
$x\in P$, write $x=p^{N+j}b_j$ for every $j$; then
$xc=p^jb_ja\in p^jA$ for all $j$, proving $xc\in P$.
Reduction of coefficients gives a surjection
\[
 \Omega(A,C)\longrightarrow\Omega(\bar A,C/P),\qquad
 \ker=P[X]=P\Omega(A,C).                                  \tag{B70.5}
\]
Indeed lift the finitely many coefficients of any reduced polynomial.
Its values on $A$ lie in the preimage of $A/P$, which is $A$, so the
lift is integer-valued. The kernel assertion follows from $PC=P$.

Both existence of co-operations and tensor injectivity descend through
(B70.5). For the latter, lift any reduced kernel tensor. Its polynomial
image belongs to $P[X,Y]$; subtract the tensor of ordinary monomials
with those coefficients, which reduces to zero. The remaining tensor
has zero polynomial image upstairs, hence is zero. This proves descent
without flatness of $A/P$.
If $0\ne z\in S$ satisfies $pz=0$, then $z\in C$ and its image in $C/P$
is nonzero, since $P\subset A$ and $p$ is regular in $A$. This contradicts
the necessary criterion B.063 for the discrete valuation base $\bar A$.
Thus the two co-operations cannot both exist uniquely when $S$ has
$p$-torsion.

Return to the global inclusion. If $S$ is integer torsion-free, so is
every localization $S_{\mathfrak m}$. The local positive cases above
and the proved localization of $\Omega_n$ show that every tensor map
is an isomorphism, by testing its kernel and cokernel at maximal ideals.
Conversely, any nonzero integer-torsion element has a nonzero multiple
$z$ killed by a prime $p$. Choose a maximal ideal containing its proper
$R$-annihilator. Then $p$ lies in that maximal ideal and $z$ survives
localization. Localizing the purported co-operations contradicts the
local necessity just proved.

For the examples, $\operatorname{Int}(\mathbb Z)$ is a binomial ring
because $\binom{F(X)}n$ is integer-valued whenever $F$ is.
In $\mathbb Z\oplus\mathbb Q\epsilon$, evaluating any binomial polynomial
at $a+b\epsilon$ gives its integer value at $a$ plus a rational multiple
of $\epsilon$, so this ring is binomial. It is not reduced and its
square-zero ideal $\mathbb Q\epsilon$ is not finitely generated over
the ring, proving non-Noetherianity.
\end{proof}

\begin{saresult}{B.071}{Proved}{All finite-residue local bases with principal maximal ideal}
Let $(A,\mathfrak m)$ be any nonzero local commutative ring with finite
residue field and $\mathfrak m=\pi A$, and let $A\subset S$ be a unital
commutative inclusion. Put
\[
                  C=\{s\in S:\pi^Ns\in A
                                  \text{ for some }N\ge0\}.
\]
Then the polynomial coaddition and comultiplication tensors are uniquely
determined precisely in the following cases:
\[
 C=A,\quad\text{or}\quad
 \pi\text{ is regular in }A\text{ and }C=A[1/\pi]
 \text{ as a canonical subring of }S.                     \tag{B71.1}
\]
Every positive case has all polynomial tensor isomorphisms.
If $\pi$ is regular in $A$, the criterion is equivalently
$\operatorname{Ann}_S(\pi)=0$.
If $\pi$ is nilpotent, the criterion reduces to $S=A$.
Neither Noetherianity nor reducedness is assumed.
\end{saresult}
\begin{proof}
If $\pi$ is nilpotent, the finite residue field and the successive cyclic
quotients $\pi^jA/\pi^{j+1}A$ make $A$ finite. Also $C=S$, and B.013
gives exactly the stated conclusion. This includes the field case
$\pi=0$. We henceforth assume $\pi$ is not nilpotent.

Set $P=\bigcap_{j\ge0}\pi^jA$. First
\[
          \operatorname{Ann}_A(\pi)\subset P,
          \qquad P=\pi P,\qquad \bar A=A/P
          \text{ is a discrete valuation ring}.           \tag{B71.2}
\]
If $x\notin P$, choose the greatest $j$ with $x\in\pi^jA$ and write
$x=\pi^ju$; then $u\notin\mathfrak m$, so $u$ is a unit.
Such an $x$ cannot be annihilated by $\pi$, as that would make
$\pi^{j+1}=0$. This proves the first inclusion without cancelling a
zero divisor. For $x\in P$ write $x=\pi y$. For every $j$ one can also
write $x=\pi^{j+1}z_j$, so
$y-\pi^jz_j\in\operatorname{Ann}_A(\pi)\subset P$.
Hence $y\in P$, proving $P=\pi P$.
If $\pi x\in P$, the same equations show $x\in P$, so $\pi$ is regular
in $\bar A$. Every nonzero element of $\bar A$ is a unit times a finite
power of $\pi$. Products are nonzero and every nonzero ideal is generated
by an element of smallest exponent. Thus $\bar A$ is a DVR, as claimed.

All coefficients of integer-valued polynomials belong to $C$.
Indeed the quotients $A/\pi^nA=\bar A/\pi^n\bar A$ are finite and have
the usual discrete-valuation digit representatives. Choose enough
distinct such representatives for a Vandermonde grid. Every difference
is a unit times a finite power of $\pi$, so Cramer's rule gives the
asserted saturation for each coefficient. Conversely, if $\pi^Nc\in A$,
take the monic polynomial that is the product of $X-a$ over a complete
set of representatives of $A/\pi^NA$. Every value is in $\pi^NA$, so
$c$ occurs as its multiplied leading coefficient. Thus $C$ is the actual
coefficient ring and $\Omega(A,S)=\Omega(A,C)$.

One has $PC=P$. For $x\in P$ and $\pi^Nc=a\in A$, write
$x=\pi^{N+j}b_j$ for every $j$, giving $xc=\pi^jb_ja\in P$.
It follows, in every number of variables, that reduction gives
\[
  \Omega_n(A,C)/P\Omega_n(A,C)
       \simeq\Omega_n(\bar A,\bar C),\qquad
  P\Omega_n(A,C)=P[X_1,\ldots,X_n],\quad \bar C=C/P.       \tag{B71.3}
\]
Every reduced polynomial lifts because its lifted values lie in the
preimage of $\bar A$, namely $A$; the asserted kernel follows from
$PC=P$.

Crucially, tensor maps lift as well as descend. Let $M=\Omega(A,C)$.
The restriction of the polynomial tensor map is an isomorphism
\[
              P M^{\otimes_A n}\ \simeq P[X_1,\ldots,X_n]. \tag{B71.4}
\]
Here is an elementwise proof, valid with zero divisors. For $h\in P$
and $F\in M$, choose $N$ with $\pi^NF\in A[X]$ and write
$h=\pi^Nh'$ with $h'\in P$, using (B71.2). Moving $\pi^N$ across a tensor
sign expands a factor $h\otimes F$ as the ordinary monomial tensor with
coefficients $h$ times the coefficients of $F$, all of which belong to
$P$ because $PC=P$. For two factors the explicit equality is
$hF\otimes G=h'F\otimes\pi^NG$ when $h=\pi^Nh'$ and
$\pi^NG\in A[X]$; expand the second factor and use $PC=P$
for every coefficient of the first.
Repeating this step in each factor writes every
element of $PM^{\otimes n}$ as that ordinary monomial tensor of its
polynomial image. Conversely these monomial tensors exist, and their
polynomial images recover their coefficients. This proves both
injectivity and surjectivity in (B71.4).

Reduction of $M^{\otimes n}$ modulo $P$ is the tensor power over
$\bar A$ of the reduced module. Together with (B71.3)--(B71.4), this
shows that every $\theta_n$ is injective, surjective, or an isomorphism
exactly when its reduced counterpart is. For injectivity, reduce a
kernel element and then use (B71.4); for surjectivity, lift a reduced
preimage and correct the remaining polynomial in $P[X_1,\ldots,X_n]$
using (B71.4). The same correction proves that coaddition or
comultiplication exists upstairs exactly when it exists downstairs.
These arguments require no flatness of $A/P$.

Apply the DVR case B.063 to $\bar A\subset\bar C$. Every coefficient
of $\bar C$ is $\pi$-power saturated over $\bar A$, so the positive cases
are exactly $\bar C=\bar A$ and $\bar C=\operatorname{Frac}(\bar A)$.
The first says $C=A$. In the second, lift an inverse of $\pi$ to obtain
$\pi c=1+x$ with $x\in P\subset\mathfrak m$. The element $1+x$ is a unit
of $A$, so $\pi$ is a unit of $C$. Injectivity of $A\subset C$ therefore
forces $\pi$ to be regular in $A$, and the definition of $C$ gives
$C=A[1/\pi]$. Both cases lift with all tensor isomorphisms by
(B71.3)--(B71.4).

Finally suppose $\pi$ is regular in $A$. If $\pi$ is regular in $S$,
then $C\subset A[1/\pi]$; any $a/\pi^N\notin A$ has
$a=\pi^ju$ for $j<N$ and $u$ a unit, so it forces $\pi^{-1}\in C$.
Thus (B71.1) holds. Conversely an element annihilated by $\pi$ lies in
$C$, and either positive option in (B71.1) has no such nonzero element.
This proves the equivalent annihilator criterion.
\end{proof}

\begin{saresult}{B.072}{Proved}{Local bases whose maximal ideal is killed by Frobenius}
Let $(R,\mathfrak m)$ be a nonzero local $\mathbb F_p$-algebra with
finite residue field $\mathbb F_q$. Suppose that $x^p=0$ for every
$x\in\mathfrak m$; the ideal $\mathfrak m$ need not be nilpotent or
finitely generated. For every unital inclusion $R\subset S$, the following
are equivalent:
\[
\begin{gathered}
 \text{the polynomial co-operations exist uniquely};\\
 \theta_2\text{ is injective};\\
 R\longrightarrow S\text{ is a ring epimorphism and }
                     (R:_S\mathfrak m)=R.                \tag{B72.1}
\end{gathered}
\]
Every positive case has all multivariable polynomial tensor isomorphisms.
Thus the nonflat example B.067 lies in a fully classified class of
possibly infinite, non-Noetherian bases. If $\mathfrak m$ is nilpotent as
an ideal, these conditions force $S=R$.
\end{saresult}
\begin{proof}
The image $R^p$ of Frobenius is a coefficient field $k\simeq\mathbb F_q$.
Indeed the kernel of Frobenius is exactly $\mathfrak m$: it contains
$\mathfrak m$ by hypothesis and contains no unit. The image is a field
and maps isomorphically onto the residue field, since Frobenius on that
finite field is an automorphism. Every element of $R$ has a unique
expression $a+n$ with $a\in k$, $n\in\mathfrak m$.

Put
\[
                    Q(X)=(X^q-X)^p=X^{qp}-X^p.
\]
It is monic and vanishes on $R$. The Chinese remainder isomorphism
\[
 k[X]/Q\ \simeq\prod_{a\in k}k[t]/(t^p),
 \qquad X\longmapsto a+t                              \tag{B72.2}
\]
gives a basis $h_{a,j}$, $a\in k$, $0\le j<p$, of polynomials of degree
less than $qp$ whose image in the $a$-component is $t^j$ and whose other
components vanish. Monic division by $Q$ and this basis work over every
$k$-algebra $S$. For $n\in\mathfrak m$ one has
$h_{a,j}(b+n)=0$ if $a\ne b$, and $h_{a,j}(a+n)=n^j$.

Write $C_j=(R:_S\mathfrak m^j)$, so $C_0=R$.
For a low-degree remainder $H=\sum c_{a,j}h_{a,j}$, preservation of $R$
is exactly the condition $c_{a,j}\in C_j$. To verify necessity, evaluate
$H(a+\lambda n)$ for fixed $a,n$ and for $p$ distinct $\lambda\in k$.
The Vandermonde matrix is invertible over $k$, so each $c_{a,j}n^j$ lies
in $R$. For $1\le j<p$, the ideal generated by all $n^j$, $n\in\mathfrak m$,
is $\mathfrak m^j$: the polarization identity
\[
 j!\,n_1\cdots n_j=
 \sum_{J\subset\{1,\ldots,j\}}(-1)^{j-|J|}
                         \left(\sum_{i\in J}n_i\right)^j
\]
and invertibility of $j!$ prove the reverse inclusion. The case $j=0$
is ordinary evaluation. Sufficiency follows by direct substitution.
Hence, as an explicitly split $R$-module,
\[
 \Omega(R,S)=\bigoplus_{a\in k,\ 0\le j<p} C_jh_{a,j}
                                \ \oplus\ Q(X)S[X].       \tag{B72.3}
\]

Assume $\theta_2$ is injective. The last summand of (B72.3) makes the
multiplication map $S\otimes_R S\to S$ injective, and it is already
surjective; hence $R\to S$ is an epimorphism. The summand indexed by
$(a,j)=(0,1)$ similarly makes $C_1\otimes_R C_1\to S$ injective.
Here the leading nonzero coefficient of $h_{0,1}$ belongs to $k^\times$,
so a coefficient tensor has zero polynomial image precisely when its
product in $S$ is zero.

The module $C_1$ is in fact a subring of $S$, with
$\mathfrak mC_1=\mathfrak m$. If $x\in\mathfrak m$ and $c\in C_1$,
then $xc\in R$ and $(xc)^p=0$; thus $xc\in\mathfrak m$.
This proves the ideal equality and shows closure of $C_1$ under products:
$\mathfrak m C_1C_1\subset\mathfrak mC_1\subset R$.
Consequently multiplication $C_1\otimes_R C_1\to C_1$ is both injective
and surjective, making $R\to C_1$ an epimorphism.
Its reduction is an injective epimorphism from the field
$R/\mathfrak m$ to $C_1/\mathfrak m$. A field admits no proper injective
epimorphic algebra: the tensors $c\otimes1$ and $1\otimes c$ are distinct
whenever $1,c$ are linearly independent. Therefore
$C_1/\mathfrak m=R/\mathfrak m$, and $C_1=R$.
This proves necessity in (B72.1).

Suppose conversely that $R\to S$ is an epimorphism and $C_1=R$.
Iterating colons gives
\[
 (R:_S\mathfrak m^{j+1})
       =((R:_S\mathfrak m^j):_S\mathfrak m)=R
                         \qquad(j\ge0).                  \tag{B72.4}
\]
Thus (B72.3) simplifies to
$\Omega(R,S)=R[X]_{<qp}\oplus Q(X)S[X]$.
In several variables, divide by $Q$ in each variable. Every basis term
with at least one $Q$ factor is a null polynomial and has an arbitrary
coefficient in $S$. For the remaining terms use the product Hermite basis.
Independent scalar Vandermonde interpolation in each nilpotent increment
forces a coefficient with jet indices $j_1,\ldots,j_n$ to multiply
$\mathfrak m^{j_1+\cdots+j_n}$ into $R$; polarization in each coordinate
justifies this whole ideal, not only individual powers. Equation (B72.4)
puts that coefficient in $R$.

The tensor power of the one-variable decomposition has the same
coefficient modules: $R$ on every all-low-degree term, and $S$ on every
term with a high-degree factor, since $S\otimes_R S=S$.
The monic bases $X^i$ for $i<qp$ and $Q(X)X^{i-qp}$ for $i\ge qp$
therefore identify every polynomial tensor map with these coefficient
isomorphisms. All $\theta_n$ are isomorphisms.
For explicit co-operations on the new generators, one can also use
\[
\begin{aligned}
 Q(X+Y)&=Q(X)+Q(Y),\\
 Q(XY)&=Q(X)Q(Y)+Q(X)Y^p+X^pQ(Y).
\end{aligned}
\]
These identities put the coaddition and comultiplication of every $cQ$
in the tensor image. Together with ordinary polynomials they generate
the whole analyser, and uniqueness and compatibility follow from the
tensor isomorphisms.

Finally, if $\mathfrak m^N=0$, equation (B72.4) gives
$S=(R:_S0)=R$. This proves the last assertion and also distinguishes
elementwise Frobenius nilpotence from nilpotence of the ideal.
\end{proof}

\begin{saresult}{B.073}{Proved under stated hypotheses}{Zero-dimensional bases with T-nilpotent nilradical}
Let $R$ be a nonzero commutative ring with nilradical $N$. Suppose
$R/N$ is absolutely flat and $N$ is T-nilpotent: for every sequence
$n_1,n_2,\ldots$ in $N$, some finite product $n_1\cdots n_j$ is zero.
Then every injective ring epimorphism out of $R$ is an isomorphism.
Consequently, for every unital inclusion $R\subset S$, the following
are equivalent:
\[
\begin{gathered}
 \text{the polynomial co-operations exist uniquely};\\
 \theta_2\text{ is injective};\qquad C=R;\\
 (R:_S D_d(R))=R\quad\text{for every }d\ge1;\\
 \operatorname{Ann}_{S/R}(D_d(R))=0\quad\text{for every }d\ge1;\\
 \Omega(R,S)=R[X].                                       \tag{B73.1}
\end{gathered}
\]
All positive cases have every multivariable polynomial tensor
isomorphism. In particular the theorem applies to every zero-dimensional
ring with T-nilpotent nilradical; the nilradical need not be nilpotent.
\end{saresult}
\begin{proof}
Let $R\subset B$ be a ring epimorphism. The ideal $NB$ is nil: each of
its elements is a finite sum of multiples of nilpotents in a commutative
ring, and is therefore nilpotent. Thus $NB\cap R=N$, and base change
gives an injective epimorphism
\[
                         R/N\longrightarrow B/NB.
\]
The absolutely flat argument in B.065 makes this an isomorphism.
Consequently $B=R+NB$.

For completeness we prove that T-nilpotence suffices to conclude
$B=R$, without imposing finite generation on $B$ or $N$. Fix $b\in B$.
Write $b=r+\sum_{i=1}^s n_i b_i$ with $r\in R$, $n_i\in N$, and repeat
this finite expansion for every $b_i$, then for each new coefficient,
and so forth. This gives a rooted, finitely branching tree. At a node
the coefficient of its remaining $B$-element is the product of the
successive $N$-labels along the path from the root. Remove nodes where
that product is zero. There is no infinite path in the remaining tree,
by T-nilpotence. Such a finitely branching tree is finite: if it had
infinitely many nodes, the root would have a child with infinitely many
descendants, and repeating that finite-choice argument would give an
infinite path. Hence after a finite expansion all remaining coefficients
are zero and the displayed expression for $b$ lies in $R$. Thus $B=R$.

If $\theta_2$ is injective, B.064 applies this fact to the actual
coefficient ring $C$ and gives $C=R$. The converse and every equivalent
ideal condition in (B73.1) follow from B.064. The positive tensor maps
are ordinary polynomial tensor maps.

Finally a reduced zero-dimensional commutative ring is absolutely flat;
this is the equivalence in the opened Stacks Lemma 15.106.5, tag 092F,
already used in B.065. Thus the stated zero-dimensional class satisfies
the hypotheses. The T-nilpotence requirement is substantive: in B.067,
a sequence of degree-zero monomials from distinct tree components has
every finite product nonzero, so that example lies outside this class.
\end{proof}

\begin{saresult}{B.074}{Proved under stated hypotheses}{Regular sequences at finite-residue points without Noetherianity}
Let $R$ be a nonzero commutative ring such that, for every maximal ideal
$\mathfrak m$ with finite residue field, the maximal ideal of
$R_{\mathfrak m}$ contains a regular sequence of length two.
For every unital extension $R\subset S$, the polynomial co-operations
exist uniquely exactly when
\[
 (R:_S D_d(R))=R\quad\text{for every }d\ge1,
 \quad\text{equivalently}\quad
 \operatorname{Ann}_{S/R}(D_d(R))=0\quad\text{for every }d. \tag{B74.1}
\]
This is also equivalent to tensor injectivity, to $C=R$, and to
$\Omega(R,S)=R[X]$. All positive cases have all tensor isomorphisms.
There is no Noetherian, reducedness or coefficient-torsion hypothesis.
The proof does not presume localization of $\Omega$.
\end{saresult}
\begin{proof}
Use the actual coefficient ring $C$ of B.064. If the second tensor map
is injective, $R\to C$ is an injective epimorphism. Localize that map
at a maximal ideal to obtain an injective epimorphism $A\subset B$.
The exact determinant-ideal localization (B69.3), which holds for
every ring, implies that each $b\in B$ is multiplied into $A$ by some
$D_d(A)$. If the residue field is infinite, that determinant ideal is
$A$, so $B=A$.

Suppose the residue field is finite, and let $a,t$ be a regular sequence
in its maximal ideal. The case $d=1$ is immediate. For $d\ge2$ the
explicit polynomial $P_d$ of B.064 gives
$a^{e_d},t^{e_d}\in D_d(A)$, since its values are those powers times
units, as in (B69.4). Write $e=e_d$. The pair $a^e,t^e$ is still a
regular sequence. For clarity, regularity of $t$ modulo $a^j$ follows
by induction: if $tx\in a^jA$, reduction modulo $a$ gives $x=ay$;
cancelling the regular $a$ reduces the assertion to exponent $j-1$.
Iterating regularity of $t$ gives regularity of $t^e$ modulo $a^e$.

Let $W$ denote the regular elements of $A$. The inclusion $A\subset B$
induces an inclusion of the total quotient ring $Q(A)$ in $W^{-1}B$.
For $b\in B$ with $D_d(A)b\subset A$, put $x=a^eb$ and $y=t^eb$,
both in $A$. The equality $t^ex=a^ey$ and regularity modulo $a^e$ show
that $x=a^ez$ for some $z\in A$. Thus the image of $b$ in $W^{-1}B$
is exactly $z$. In particular every rationalized element of $B$ lies
in the embedded copy of $A$.

Restriction of rationalization therefore gives a ring retraction
$\rho:B\to A$ of the inclusion. The two endomorphisms of $B$ consisting
of its identity and of inclusion composed with $\rho$ agree on $A$.
The epimorphism property forces them to agree on $B$, hence $B=A$.
This holds at every maximal ideal, so $C/R=0$ by the arbitrary-module
localization test. B.064 gives precisely the equivalences in (B74.1)
and the positive ordinary-polynomial tensor isomorphisms.
\end{proof}

\begin{saresult}{B.075}{Proved}{Infinite-residue conductor cusps with arbitrary coefficient extensions}
Let $k=\mathbb F_q\subset l$ with $l$ an infinite field, let $a\ge1$, and put
\[
 V=l[[t]],\qquad K=l((t)),\qquad R=k+t^aV,\qquad
                         \mathfrak m=t^aV.
\]
For every unital commutative inclusion $R\subset S$, with arbitrary
$R$-torsion in $S$, the following are equivalent:
\[
\begin{gathered}
 \text{the polynomial co-operations exist uniquely};\\
 \theta_2\text{ is injective};\qquad
 (R:_S\mathfrak m)=R;\\
 \operatorname{Ann}_{S/R}(\mathfrak m)=0;\qquad
 \Omega(R,S)=R[X].                                       \tag{B75.1}
\end{gathered}
\]
Every positive case has all multivariable tensor isomorphisms.
This removes the torsion-free restriction from the entire family B.043,
including the explicit base $\mathbb F_2+t\mathbb F_2(u)[[t]]$.
\end{saresult}
\begin{proof}
We first identify the coefficient ring without any localization theorem
for integer-valued modules. The determinant ideals $D_d(R)$ are cofinal
with the powers of $\mathfrak m$. For fixed $d$, choose $d+1$ elements
of $t^a l$ with distinct coefficients in $l$. Their determinant has
the form $t^{ad(d+1)/2}u$ with $u\in V^\times$. Multiplying it by
$\mathfrak m$ shows that $D_d(R)$ contains a power of $\mathfrak m$.
In the opposite direction, every tuple of $d+1$ nodes has residue
multiplicities $n_1,\ldots,n_q$; its determinant belongs to
$\mathfrak m^{\sum_i\binom{n_i}{2}}$. The minimum of this exponent
tends to infinity with $d$. Thus B.064 gives the actual coefficient ring
\[
 C=\bigcup_{n\ge0}(R:_S\mathfrak m^n),\qquad
                         \Omega(R,S)=\Omega(R,C).         \tag{B75.2}
\]
It is a ring. If $(R:_S\mathfrak m)=R$, iterating colons gives $C=R$,
so every positive conclusion follows. We prove that tensor injectivity
forces $C=R$.

Assume $\theta_2$ is injective, so $R\to C$ is an epimorphism by B.064.
Localizing at the nonzero elements of $R$ maps every element of $C$ into
the embedded fraction field $K$: some nonzero element of
$\mathfrak m^n$ multiplies it into $R$. Let $L\subset K$ be this image,
and $T$ its kernel. Every element of $T$ is annihilated by a power of
$\mathfrak m$, because its products with that power lie in $R$ and
vanish after localization. The composite $R\to L$ is an epimorphism.

If $L\subset V$, the ideal $\mathfrak m$ is an ideal of $L$ with
$\mathfrak mL=\mathfrak m$. Reduction gives an injective epimorphism
from $k$ to $L/\mathfrak m$, forcing equality and $L=R$.
If $L$ contains an element of negative valuation, a positive power of
its inverse belongs to $\mathfrak m$. Thus $L$ contains the inverse of
some nonzero $b\in\mathfrak m$, and $bV\subset\mathfrak m\subset R$
implies $V\subset L$. An overring of the discrete valuation ring $V$
with an element of negative valuation is $K$. Hence $L=R$ or $K$.
In the former case rationalization is a ring retraction $C\to R$;
epicness forces $C=R$.

In the latter case $C$ has a precise elementary decomposition. Put
$s=t^a\in R$, choose $c\in C$ with rational image $s^{-1}$, and set
$x=sc$. Since $x-1\in T$, choose $N\ge1$ with $s^N(x-1)=0$.
Multiplication by $c^N$ gives $x^N(x-1)=0$. Therefore
$e=x^N$ is an idempotent and $ex=e$. On $eC$, the element $s$ is
invertible with inverse $ec$. All $s$-power torsion vanishes there,
so rationalization identifies $eC$ with $K$.
The unit $1-e$ of $D=(1-e)C$ belongs to $T$, so
$\mathfrak m^M D=0$ for some $M$.
The epimorphism $R\to D$ is surjective. Indeed, modulo its nilpotent
maximal-ideal image it is an epimorphism from the field $k$, hence an
isomorphism if $D\ne0$. Thus $D$ is the image of $R$ plus
$\mathfrak mD$, and finite iteration of $\mathfrak m^MD=0$ proves
surjectivity. Consequently
\[
             C\simeq K\times R/I,\qquad
             \mathfrak m^M\subset I,                     \tag{B75.3}
\]
with the diagonal inclusion of $R$; $I=R$ is allowed and means $C=K$.

We exclude all possibilities in (B75.3) by nonzero tensor kernels.
Write $P(X)=X^q-X$ and $U=P(X)/t^a$. The following explicit ring is
central to the argument:
\[
 B=\{F\in K[X]:F(R)\subset V\}
   =V[X,U]/(X^q-X-t^aU)
   =\bigoplus_{0\le i<q,\ j\ge0}VX^iU^j.                 \tag{B75.4}
\]
Here the presentation means its embedded image in $K[X]$.
The displayed monomials are a $K$-basis of $K[X]$, by successive
division by the monic $P$. They all preserve $R$ into $V$.
If an integer-valued polynomial had a coefficient of negative minimum
$t$-valuation in this basis, multiplication by the opposite power of
$t$ would give a polynomial with nonzero reduction modulo $t$ whose
values are all zero modulo $t$. But
\[
 B/tB\simeq l[X,U]/(X^q-X)\simeq\prod_{b\in k}l[U],       \tag{B75.5}
\]
and evaluating $X=b+t^av$ with $v\in l$ sends $U$ modulo $t$ to $-v$.
The infinitude of $l$ makes such a nonzero polynomial impossible.
This proves (B75.4).

For $I=R$, the same evaluation modulo $t^a$ gives the split module
formula
\[
 \Omega(R,K)=
 \bigoplus_{0\le i<q}RX^i
 \ \oplus\!!\bigoplus_{0\le i<q,\ j\ge1}t^aVX^iU^j.     \tag{B75.6}
\]
To justify all $a$ layers, on constant $v\in l$ one has
$X\equiv b$ and $U\equiv-v$ modulo $t^a$.
The coefficient of $t^0$ must be a $k$-valued polynomial on the
infinite field $l$, hence constant in $k$; every coefficient of
$t^r$, $0<r<a$, must vanish. Interpolation on the $q$ possible $b$
then gives exactly (B75.6).
Choose $\alpha\in l\setminus k$. The tensor
\[
                  \alpha P\otimes P-P\otimes\alpha P     \tag{B75.7}
\]
has zero polynomial image and is nonzero. Projection to the
$(i,j)=(0,1)$ summand of (B75.6), followed by division by $t^a$ and
reduction modulo $t$, sends it to
$\alpha\otimes1-1\otimes\alpha\ne0$ in $l\otimes_k l$.
The projection is $R$-linear via $R\to k$. This excludes $C=K$.

Now suppose $I\ne R$. Choose the least valuation $b\ge a$ of a
nonzero element of $I$. Multiplying such an element by
$\mathfrak m=t^aV$ gives
\[
                      t^{b+a}V\subset I\subset t^bV.
\]
Put $H=\{F\in K[X]:F(R)\subset I\}$ and let $E$ be the first projection
of $\Omega(R,K\times R/I)$ into $K[X]$. Lifting the coefficients of
its second polynomial component to $R$ shows exactly that
\[
                  E=R[X]+H,\qquad
                  t^{b+a}B\subset H\subset t^bB.           \tag{B75.8}
\]
We claim $H\not\subset\mathfrak mE$.

For a polynomial in $t^jB$, call its image after division by $t^j$
in (B75.5) its leading residue at level $j$.
Let $d$ be the least level for which some element of $H\cap t^dB$
has a leading residue with a nonconstant $U$-component. Such a level
exists with $b\le d\le b+a$, since $t^{b+a}U\in H$.
Choose an element witnessing it and a component indexed by
$b_0\in k$ whose residue polynomial $h(U)$ is nonconstant.
Precompose it with
\[
                    \psi_N(X)=b_0+t^aU(X)^N.
\]
This polynomial maps $R$ into $R$, so the composition remains in $H$.
Its leading residue at level $d$ is $h(-U^N)$ in every component,
since $\psi_N$ has residue $b_0$ and
$U(\psi_N)\equiv-U^N\pmod t$.
Taking $N$ large gives an element $F\in H\cap t^dB$ whose leading
$U$-degree is strictly larger than $\lfloor d/a\rfloor-1$.

This $F$ cannot lie in $\mathfrak mE$.
In the monic basis $X^iP^j$, one has
\[
 \mathfrak mR[X]=
       \bigoplus_{0\le i<q,\ j\ge0}t^{a(j+1)}VX^iU^j.
\]
Therefore the leading residue at level $d$ of an element of
$\mathfrak mR[X]\cap t^dB$ has $U$-degree at most
$\lfloor d/a\rfloor-1$.
Also $\mathfrak mH\subset t^{b+a}B$. If $d<b+a$ it contributes
nothing at level $d$. If $d=b+a$, its leading residues are linear
combinations over $l$ of leading residues of $H$ at level $b$;
these are $U$-constant by the minimality of $d$.
Thus adding $\mathfrak mH$ cannot supply the missing higher degree.
This proves $F\notin\mathfrak mE$.

The polynomial $f=(F,0)$ belongs to $\Omega(R,C)$, and is nonzero
modulo $\mathfrak m\Omega$ because its first projection is outside
$\mathfrak mE$. Choose $n$ with $an\ge b+a$ and put
$g=(0,P^n\bmod I)$. It is a null polynomial on $R$, so also belongs
to $\Omega(R,C)$. Its second component is monic and remains nonzero
over $k=R/\mathfrak m$, so $g\notin\mathfrak m\Omega$.
The pure tensor
\[
                              f\otimes g                  \tag{B75.9}
\]
has zero polynomial image, because its two coefficient components
have disjoint idempotent support. It is nonzero: reduction modulo
$\mathfrak m$ gives the tensor of two nonzero vectors over the field
$k$. This excludes every proper $I$ in (B75.3).

All cases with $C\ne R$ have contradicted tensor injectivity.
Finally $C=R$ is equivalent to $(R:_S\mathfrak m)=R$ by the iterated
colon identity, and the latter is precisely the annihilator condition
in (B75.1). B.064 supplies the remaining conclusions.
\end{proof}

\begin{saresult}{B.076}{Proved}{Principal-conductor pullbacks in arbitrary valuation rings}
Let $V$ be a valuation domain of characteristic $p$ with infinite residue
field $l$, containing a finite subfield $k=\mathbb F_q$. Let
$0\ne x\in\mathfrak m_V$, and put
\[
                 R=k+xV,\qquad\mathfrak m=xV.
\]
For every unital commutative inclusion $R\subset S$, the polynomial
coaddition and comultiplication tensors exist uniquely if and only if
\[
                           (R:_S\mathfrak m)=R.           \tag{B76.1}
\]
This is equivalent to tensor injectivity, to the absence of a
$\mathfrak m$-annihilator in $S/R$, and to $\Omega(R,S)=R[X]$.
All positive cases have every multivariable tensor isomorphism.
No discreteness, rank, completeness, or coefficient-torsion condition
is imposed. The conductor $xV$ is principal as a $V$-ideal; it need
not be principal as an $R$-ideal.
\end{saresult}
\begin{proof}
Let $v$ be the valuation, $\eta=v(x)>0$, and $K=\operatorname{Frac}(V)
=\operatorname{Frac}(R)$. The ring $R$ is local with maximal ideal
$\mathfrak m$ and residue field $k$. The determinant-ideal proof in
B.075 works verbatim with $x$ in place of $t^a$: a grid $xv_i$ with
distinct residues of $v_i\in V$ gives a determinant $x^{d(d+1)/2}$ times
a unit of $V$, and multiplication by $\mathfrak m$ puts a power of
$\mathfrak m$ in its principal $R$-ideal. The opposite cofinality follows
by residue multiplicities. Thus the actual coefficient ring is
\[
                 C=\bigcup_{n\ge0}(R:_S\mathfrak m^n).    \tag{B76.2}
\]
Under tensor injectivity $R\to C$ is epic. Rationalization has image
$L$ satisfying $R\subset L\subset R[1/x]=V[1/x]$, and kernel whose
elements are annihilated by powers of $\mathfrak m$.

There are only two possibilities for this epimorphic overring $L$.
If every $c\in L$ has $v(c)\ge-\eta$, then
$\mathfrak m^2L\subset\mathfrak m\subset R$.
The conductor $J=(R:L)$, an ideal of both rings, therefore contains
$\mathfrak m^2$. The inclusion $R/J\subset L/J$ is epic and its base
has nilpotent maximal ideal and residue field $k$. Reducing modulo
that ideal gives a field epimorphism, hence an isomorphism; finite
iteration of the nilpotent ideal then gives $L/J=R/J$ and $L=R$.
Consequently if $L\ne R$, some $c\in L$ has $v(c)<-\eta$.
Then $c^{-1}\in xV\subset R$ and $c^{-1}V\subset xV$, so $V\subset L$.
Also $x^{-1}=c(xc)^{-1}$ and $(xc)^{-1}\in V$, yielding $L=V[1/x]$.
If $L=R$, rationalization retracts $C$ to $R$ and epicness gives $C=R$.

If $L=V[1/x]$, the polynomial idempotent argument of B.075 applies
unchanged: lift $x^{-1}$ to $c\in C$, use
$x^N(xc-1)=0$ to make $(xc)^N$ idempotent, and split its two factors.
The $x$-invertible factor is $L$, and the other has a power of
$\mathfrak m$ acting by zero. A ring epimorphism from a local ring
with nilpotent maximal ideal is surjective, by the field quotient
and finite iteration just used. Hence
\[
                    C\simeq L\times R/I,
                    \mathfrak m^M\subset I               \tag{B76.3}
\]
for some $M$, allowing $I=R$.

Put $P(X)=X^q-X$ and $U=P(X)/x$. One has
\[
 \begin{aligned}
 B&=\{F\in K[X]:F(R)\subset V\}=V[X,U]/(X^q-X-xU)\\
  &=\bigoplus_{0\le i<q,\ j\ge0}VX^iU^j,
 & B/\mathfrak m_VB&\simeq\prod_{b\in k}l[U].
 \end{aligned}                                      \tag{B76.4}
\]
Every polynomial has a finite expansion in this basis over $K$.
If its minimum coefficient valuation were negative, divide by a
coefficient attaining that minimum. The resulting polynomial has
nonzero residue, but its values on $R$ have positive valuation.
Evaluation at $b+xz$, with $z\in V$ varying, has residues
$X=b$, $U=-\bar z$; the infinite field $l$ makes this impossible.
This proves (B76.4), including non-Archimedean value groups.
In particular all these integer-valued polynomials have their
coefficients in $L$.

In this basis the $R$-valued module is
\[
 \Omega(R,K)=
 \bigoplus_{0\le i<q}RX^i
 \ \oplus\!!\bigoplus_{0\le i<q,\ j\ge1}xVX^iU^j.       \tag{B76.5}
\]
To see this directly, subtract the degree-less-than-$q$ interpolation
polynomial with the same values on $k$; it has coefficients in $R$.
The remainder has no $j=0$ coefficients. If its minimum valuation
$\gamma$ is positive but less than $\eta$, all its values lie in
$R\cap\mathfrak m_V=xV$; division by a minimum coefficient then gives
a nonzero residue polynomial vanishing on all $k\times l$, impossible.
If $\gamma=0$, the residue polynomial in each component has values
in the finite field $k$, hence is constant; its zero constant term
makes it zero, again impossible. Thus $\gamma\ge\eta$.
The converse is immediate. If $I=R$ in (B76.3), formula (B76.5) gives
exactly the nonzero tensor kernel $\alpha P\otimes P-P\otimes\alpha P$
for $\alpha\in V$ whose residue lies outside $k$, as in B.075.

For a proper $I$ in (B76.3), let
$H=\{F\in L[X]:F(R)\subset I\}$ and $E=R[X]+H$, the first projection of
$\Omega(R,L\times R/I)$. We prove $H\not\subset\mathfrak mE$ by an
integer-bound argument which requires no least valuation in an ideal.
For $F\in H$, subtract its degree-less-than-$q$ interpolation on $k$;
that polynomial has coefficients in $I$. Denote the resulting module
of polynomials with all $j=0$ coefficients zero by $H_0$.
Every nonzero member of $H_0$ has minimum coefficient valuation at
least $\eta$, by (B76.5). Moreover $x^MU\in H_0$ by (B76.3).
Choose the least positive integer $n$ for which some nonzero
$G\in H_0$ has minimum coefficient valuation $\gamma\le n\eta$.
This integer exists, and it is a minimum in the ordinary integers,
not an assumed minimum in the valuation group.

After division by a minimum coefficient, $G$ has a nonzero residue
component $h(U)$ at some $b_0\in k$. Since its $j=0$ coefficients
vanish, $h(0)=0$ and $h$ is nonconstant.
Precompose with $\psi_N(X)=b_0+xU(X)^N$, for an integer $N>n$.
This maps $R$ into $R$ and preserves vanishing on $k$, so
$F=G\circ\psi_N\in H_0$.
The minimum coefficient valuation of $F$ is still $\gamma$, and its
normalized residue is $h(-U^N)$ in every component. Hence some
coefficient of $U$-degree at least $N$ has valuation $\gamma$.

If $H\subset\mathfrak mE$, write
$F=A+\sum_i r_iF_i$, where $A\in\mathfrak mR[X]$,
$r_i\in\mathfrak m$, and $F_i\in H$.
Subtract the interpolation polynomials of the $F_i$ and absorb their
products into $A$, so that every $F_i\in H_0$.
The coefficient of $X^iU^j$ in $\mathfrak mR[X]$ lies in
$x^{j+1}V$, since the corresponding monic $R$-basis is $X^iP^j$.
At the chosen degree $j\ge N>n$ its valuation is therefore greater
than $\gamma$. Some summand $r_iF_i$ must have a coefficient there
of valuation at most $\gamma$. The minimum coefficient valuation
of that nonzero $F_i$ is at most
\[
                  \gamma-v(r_i)\le\gamma-\eta
                                     \le(n-1)\eta,
\]
If $n=1$, this contradicts the lower bound $\eta>0$ already proved
for every nonzero element of $H_0$. If $n>1$, it contradicts the
minimal choice of $n$. Thus $H\not\subset\mathfrak mE$.

Choose $F\in H\setminus\mathfrak mE$. As in B.075,
$f=(F,0)$ survives modulo $\mathfrak m\Omega$ by its first projection.
The polynomial $g=(0,P^M\bmod I)$ is null on $R$ and survives there
by its monic second projection over $k$.
Their pure tensor is nonzero modulo $\mathfrak m$ but its polynomial
image is zero. This excludes every proper $I$ in (B76.3).
We have proved that tensor injectivity forces $C=R$.
The iterated colon identity makes $C=R$ equivalent to (B76.1), and
B.064 proves all remaining positive assertions.
\end{proof}

\begin{saresult}{B.077}{Proved}{Every conductor cut in a real-valued valuation domain}
Let $V$ be a characteristic-$p$ valuation domain with infinite residue
field $l$ and a nontrivial valuation
$v:K^\times\to\Gamma\subset\mathbb R$, where
$K=\operatorname{Frac}(V)$. Suppose $V$ contains $k=\mathbb F_q$.
Let $J$ be any nonzero proper ideal of $V$ and put $R=k+J$.
For every unital commutative inclusion $R\subset S$, with arbitrary
base torsion, the following are equivalent:
\[
\begin{gathered}
 \text{the polynomial co-operations exist uniquely};\\
 \theta_2\text{ is injective};\qquad (R:_S J)=R;\\
 \operatorname{Ann}_{S/R}(J)=0;\qquad \Omega(R,S)=R[X].    \tag{B77.1}
\end{gathered}
\]
All positive cases have every multivariable polynomial tensor
isomorphism. The ideal $J$ may be an open or irrational valuation cut;
it need not have a generator or a least element of its value set.
In particular this settles the previously listed Hahn example with
residue $\mathbb F_2(u)$, value group $\mathbb Q$, and $J=\{v>\sqrt2\}$.
\end{saresult}
\begin{proof}
The ring $R$ is local with maximal ideal $J$ and residue field $k$.
Let $\alpha=\inf v(J\setminus\{0\})\ge0$.
If $\alpha=0$, then $J$ is the maximal ideal of $V$. For $r\in J$ and
any $n$, choose $y\in J$ with $nv(y)<v(r)$. The quotient $r/y^n$ has
positive valuation and belongs to $J\subset R$. Thus $J$ is generated
as an $R$-ideal by its $n$-th powers for every $n$. B.069, including
its determinant calculation, gives exactly (B77.1).
We henceforth assume $\alpha>0$.

The determinant ideals are cofinal with powers of $J$.
Residue multiplicities give $D_d(R)\subset J^n$ for sufficiently
large $d$, as in B.075. Conversely choose $c\in J$ and a grid
$cv_i$ whose $v_i\in V$ have distinct residues in $l$.
Its determinant has valuation $h v(c)$, where $h=d(d+1)/2$.
Choose $n$ with $n\alpha>hv(c)+\alpha$. Every element of $J^n$ divided
by that determinant then has valuation greater than $\alpha$, hence
belongs to $J\subset R$. Thus $J^n\subset D_d(R)$.
Consequently the actual coefficient ring is
\[
                    C=\bigcup_{n\ge0}(R:_S J^n).          \tag{B77.2}
\]
The colon identity already proves the positive implication from
$(R:_S J)=R$. Assume tensor injectivity and hence epicness of $R\to C$.

Rationalization maps $C$ to an epimorphic overring $L\subset K$, with
kernel killed elementwise by powers of $J$. If $L\subset V$, the common
ideal $J$ reduces this to a field epimorphism $k\to L/J$, so $L=R$.
If $L$ contains an element of negative valuation, a power of its inverse
has valuation greater than $\alpha$ and belongs to $J$; multiplying
$V$ by that element puts it in $J$, so $V\subset L$.
The real-valued rank-one property then gives $L=K$.
If $L=R$, the ring retraction forces $C=R$.
If $L=K$, choose any nonzero $x\in J$ and apply the polynomial
idempotent construction of B.075 to a lift of $x^{-1}$.
The same nilpotent-quotient argument gives
\[
                  C\simeq K\times R/I,\qquad J^M\subset I
                    \quad\text{for some }M\ge1,           \tag{B77.3}
\]
allowing $I=R$.

Put $P(X)=X^q-X$. The monomials $X^iP(X)^j$, $0\le i<q$, $j\ge0$,
form a monic $K$-basis of $K[X]$.
Their precise $R$-valued coefficient modules are
\[
 \Omega(R,K)=\bigoplus_{0\le i<q}RX^i
 \ \oplus\!\bigoplus_{0\le i<q,\ j\ge1}
                          (J:_KJ^j)X^iP(X)^j.            \tag{B77.4}
\]
Here is a direct proof of necessity, including the valuation boundary.
Subtract the degree-less-than-$q$ interpolation polynomial of the
values on $k$; its coefficients lie in $R$, leaving all $j=0$
coefficients zero. Fix $\gamma\in v(J)$ and choose $c\in J$ of that
value. In evaluating at $b+cz$, with $b\in k$ and $z\in V$, the minimum
term valuation is
\[
                   \delta=\min_{0\le i<q,\ j\ge1}
                                 \bigl(v(a_{ij})+j\gamma\bigr).
\]
Divide by a coefficient attaining this minimum after multiplication
by $c^j$. In residue, interpolation over $b\in k$ shows that for some
$b$ the result is a nonzero polynomial in $\bar z\in l$, with no
constant term. If $\delta<0$, a value avoiding its finitely many roots
contradicts $V$-valuedness. If $\delta=0$, the unnormalized residue polynomial takes the
infinite field $l$ into $k$. Equivalently, the normalized polynomial
takes values in a finite scalar multiple of $k$. Either way a polynomial
with finite image on the infinite field $l$ is constant, again impossible. If $\delta>0$ but is outside the value set
of $J$, a value with valuation exactly $\delta$ contradicts
$R\cap\mathfrak m_V=J$. Thus $\delta$ belongs to the value set of $J$.
Every coefficient consequently satisfies $a_{ij}c^j\in J$ for every
$c\in J$. In a valuation ring, products of $j$ elements of $J$ have
valuation at least $j$ times the least valuation among those elements.
This proves $a_{ij}J^j\subset J$. Sufficiency of these conditions is
termwise evaluation, proving (B77.4).

In particular $(J:_KJ)=V$. A coefficient of negative valuation cannot
stabilize $J$, since an element with valuation sufficiently close to
$\alpha$ is sent below $\alpha$; every coefficient in $V$ does stabilize
the ideal. Therefore the $j=1,i=0$ summand of (B77.4) is $VP$.
If $I=R$ in (B77.3), choose $\beta\in V$ with residue outside $k$.
The tensor $\beta P\otimes P-P\otimes\beta P$ has zero polynomial
image and maps to a nonzero tensor in $l\otimes_k l$ under this split
coefficient projection and residue. Thus $C=K$ is excluded.

Suppose $I$ is proper. Define explicitly
\[
 H=\{F\in K[X]:F(R)\subset I\},\qquad E=R[X]+H.
\]
The latter is exactly the first projection of $\Omega(R,K\times R/I)$.
Let $H_0$ consist of the polynomials in $H$ vanishing on $k$; subtraction
of interpolation on $k$, with coefficients in $I$, maps $H$ into $H_0$.
For a nonzero $F=\sum_{0\le i<q,\ j\ge1}a_{ij}X^iP^j\in H_0$, define its real
weighted value by
\[
             w(F)=\min_{0\le i<q,\ j\ge1}\bigl(v(a_{ij})+j\alpha\bigr).
\]
In this formula $i$ ranges over $0\le i<q$; only $j$ is required to
be positive. Taking $\gamma$ down to $\alpha$ in the coefficient
criterion (B77.4) shows $w(F)\ge\alpha$.
The module $H_0$ is nonzero: if $0\ne c\in J^{M-1}$, then $cP\in H_0$.
Hence
\[
             \lambda=\inf\{w(F):0\ne F\in H_0\}
                   \quad\text{is finite and }\lambda\ge\alpha>0.
                                                               \tag{B77.5}
\]

We show $H\not\subset JE$. Choose $G\in H_0$ with
$w(G)<\lambda+\alpha/4$ and let $D=\deg G\ge1$.
For some $b_0\in k$, the ordinary polynomial $G(b_0+Y)$ has Gauss
value $w(G)$ when $Y$ is assigned weight $\alpha$.
Indeed every term of its expansion has at least that weight. Choose
a $P$-index $j$ attaining the minimum; interpolation on $k$ makes the
coefficient polynomial of that $P^j$ have minimum valuation at some
$b_0$. Its contribution to the coefficient of $Y^j$ has weight
$w(G)$, while contributions from smaller $P$-indices have strictly
larger weight because $\alpha>0$. This proves the assertion.

Choose an integer $N$ with $(N+1)\alpha>\lambda+\alpha$, and choose
$c_N\in K$ such that
\[
        0\le\varepsilon=v(c_N)+(N-1)\alpha<\frac\alpha{4D}.
                                                               \tag{B77.6}
\]
If $\alpha\in\Gamma$, choose equality zero. Otherwise $\Gamma$ is
dense in $\mathbb R$ and an element of the indicated interval exists.
For completeness, an additive subgroup of $\mathbb R$ which is not
dense has a least positive element and is cyclic; every nonzero ideal
cut in that cyclic value group is attained, so its infimum belongs
to the group. Thus these two choices cover every case.
The polynomial
\[
                       \psi_N(X)=b_0+c_NP(X)^N
\]
maps $R$ into $R$. For $P(r)\ne0$, its valuation lies in the value set
of $J$, so the nonconstant term has valuation
$\alpha+\varepsilon+N(v(P(r))-\alpha)$, in that same value set.
If the cut is open this is strictly greater than $\alpha$; if it is
closed equality is allowed.

Put $F=G\circ\psi_N$. It belongs to $H_0$, and its expansion is a
polynomial in $P^N$ with zero constant term. Its weighted value is
exactly the Gauss value of $G(b_0+Y)$ with weight
$\alpha+\varepsilon$, since distinct powers of $Y$ become distinct
powers of $P^N$. Consequently
\[
                  w(F)\le w(G)+D\varepsilon
                                  <\lambda+\alpha/2,      \tag{B77.7}
\]
and every nonzero $P$-index of $F$ is at least $N$.

If $H\subset JE$, write $F=A+\sum_i r_iF_i$, with
$A\in JR[X]$, $r_i\in J$, $F_i\in H$.
Remove the interpolation polynomials of the $F_i$ and absorb them
into $A$, so every $F_i\in H_0$.
Every coefficient of $A$ in the monic $X^iP^j$ basis belongs to $J$,
so its weight at a degree $j\ge N$ is at least
$(j+1)\alpha>\lambda+\alpha$.
At a monomial attaining $w(F)$, a summand $r_iF_i$ must therefore
have weight at most $w(F)$. This implies
\[
                   w(F_i)\le w(F)-v(r_i)
                         \le w(F)-\alpha<\lambda,
\]
contrary to (B77.5). Thus $H\not\subset JE$.

Choose $F\in H\setminus JE$. The polynomial $f=(F,0)$ survives modulo
$J\Omega$ by its first projection. The null polynomial
$g=(0,P^M\bmod I)$ survives there by its monic second projection over
$k$. Their tensor is nonzero modulo $J$ and has zero polynomial image,
excluding every proper $I$ in (B77.3).
We have proved $C=R$ under tensor injectivity. The iterated-colon
identity and B.064 now give every assertion in (B77.1).
\end{proof}

\begin{saresult}{B.078}{Proved under stated hypotheses}{Every valuation conductor ideal, without a rank restriction}
Let $V$ be a valuation domain with fraction field $K$ and infinite
residue field $l$. Let $R\subset V$ be a unital subring containing a
nonzero proper $V$-ideal $J$, and suppose $R/J=k$ is a finite field
of cardinality $q$. No coefficient-field section or characteristic
assumption is imposed. For every unital
commutative extension $R\subset S$, the following are equivalent:
\[
\begin{gathered}
 \text{the standard polynomial co-operations exist uniquely};\\
 \theta_2\text{ is injective};\qquad (R:_S J)=R;\\
 \operatorname{Ann}_{S/R}(J)=0;\qquad \Omega(R,S)=R[X].
                                                        \tag{B78.1}
\end{gathered}
\]
Every positive case has all multivariable polynomial tensor isomorphisms.
The valuation may have arbitrary rank and arbitrary cardinality; the ideal
may be nonprincipal and need not have a value-group infimum.
\end{saresult}
\begin{proof}
$R$ is local with maximal ideal $J$, and $R\cap\mathfrak m_V=J$.
Indeed the unital map from the field $R/J$ to $V/\mathfrak m_V$ is
injective. For $r\in R\setminus J$, choose $a\in R$ with $ra=1+j$,
$j\in J$; the inverse of $1+j$ lies in $1+J\subset R$, since $J$
is an ideal of $V$. Thus $r$ is a unit of $R$.
We first record the ideal identity that replaces a real-valued infimum:
\[
                 J^N(J:_K J^N)=J\qquad(N\ge1).           \tag{B78.2}
\]
Here all ideal products and ideal quotients in the proof are taken over
$V$, unless an $R$-module is explicitly specified. This is the valuation
ideal form of Clifford regularity: Bazzoni--Salce, Theorem 3(3), pp.~139--140,
and the complete elementary proof in Halter--Koch, Lemmas 5.1--5.3 and
Theorem 5.4(1), author-PDF pp.~8--9, were opened
\cite{BazzoniSalce1996,HalterKoch2009}. For completeness we give the proof
needed here. Set $E=(J:J)$; it is a valuation overring of $V$, with maximal
ideal $\mathfrak p$. Every overring of a valuation domain is a localization:
an element outside it has its inverse inside it, its maximal ideal is
contained in $V$, and inverting the elements of $V$ outside that ideal
gives the overring. In particular $\mathfrak p\subset V$.
If $J=aE$ is principal over $E$, then $J^2a^{-1}=J$.
Otherwise $J\mathfrak p=J$: if $a\in J\setminus J\mathfrak p$, comparison
of any $b\in J$ with $a$ in the valuation ring $E$ shows $b\in aE$.
Indeed if $b\notin aE$, then $a/b\in\mathfrak p$ and
$a\in b\mathfrak p\subset J\mathfrak p$, a contradiction.
In this nonprincipal case $J(E:J)=\mathfrak p$. The left side is a proper
ideal of $E$ (otherwise $J$ would be invertible and hence principal in
the local valuation ring). Conversely, for $z\in\mathfrak p$ one has
$zJ\subsetneq J$, since equality would put $z^{-1}$ in $(J:J)=E$.
Choose $a\in J\setminus zJ$. Comparison of ideals gives $zJ\subset aE$,
so $z/a\in(E:J)$ and $z\in J(E:J)$. Consequently
$J^2(E:J)=J\mathfrak p=J$. In both cases $J=J^2A$ for a fractional
$V$-ideal $A$. Iteration gives $J=J^N A^{N-1}$ and
$A^{N-1}\subset(J:J^N)$, proving (B78.2); the reverse inclusion is
the definition of the ideal quotient.

If $J^2=J$, then $J$ is generated as an $R$-ideal by its $N$th powers
for every $N$. For $0\ne r\in J=J^{N+1}$ choose, in a finite expression
of $r$ as a sum of products of $N+1$ elements of $J$, a product of
valuation at most $v(r)$. If $y$ is a factor of smallest valuation, then
$(N+1)v(y)\le v(r)$, so $r/y^N\in J\subset R$.
Thus B.069 applies and proves (B78.1) in this case.
We may therefore suppose $J^2\ne J$ and fix
\[
             0\ne c\in J\setminus J^2,\qquad W=V[1/c]=R[1/c].
                                                        \tag{B78.3}
\]
The latter equality follows from $cV\subset J\subset R$.
For every $y\in J\setminus\{0\}$ one has $2v(y)>v(c)$, since the
opposite inequality would give $c\in y^2V\subset J^2$.
Taking a factor of smallest valuation in a product gives
\[
                      J^{2h+1}\subset c^hJ\quad(h\ge0). \tag{B78.4}
\]
All statements about valuations here concern inequalities in the actual
ordered value group, with no Archimedean assumption.

The Vandermonde ideals $D_d(R)$ and the powers of $J$ are cofinal.
For one direction, sufficiently many nodes in the $q$ residue classes
force arbitrarily many differences into $J$, so $D_d\subset J^n$ for
large $d$ and any prescribed $n$. For the other, choose $d+1$ elements
$u_0,\ldots,u_d\in V$ with distinct residues in $l$. The determinant
on the nodes $cu_i\in J$ is $\delta=c^h u$, where
$h=d(d+1)/2$ and $u\in V^\times$. By (B78.4),
$J^{2h+1}\subset c^hJ=\delta J\subset\delta R\subset D_d$.
B.064 therefore identifies the effective coefficient ring as
\[
                         C=\bigcup_{n\ge1}(R:_S J^n).   \tag{B78.5}
\]
If $(R:_S J)=R$, induction gives every term in this union equal to $R$,
and B.064 proves all the positive assertions. It remains to show that
injectivity of $\theta_2$ forces $C=R$.

Assume that injectivity. By B.064 the inclusion $R\to C$ is epic.
Localizing at the nonzero elements of $R$ maps $C$ into $K$: an element
of the $n$th term of (B78.5) has its product with $c^n$ in $R$.
Let $L\subset W$ be its rational image and $T$ its kernel. The map
$R\to L$ is epic, and every element of $T$ is annihilated by some power
of $J$. If all $b\in L\setminus\{0\}$ satisfy $v(b)\ge-v(c)$, then
(B78.4) gives $J^3L\subset J$. The common ideal $J^3L$ is an ideal of
both $R$ and $L$, contains $J^3$, and is proper in $R$. The induced
injective epimorphism $R/J^3L\to L/J^3L$ has a local source with
nilpotent maximal ideal. B.073, or its finite-depth expansion proof,
makes this map an isomorphism; hence $L=R$.
Otherwise choose $b\in L$ with $v(b)<-v(c)$. Then $b^{-1}V\subset J$,
so $V\subset L$, and $c^{-1}=b(cb)^{-1}\in L$ because $(cb)^{-1}\in V$.
It follows that $L=W$. If $L=R$, rationalization is a retraction of
the epimorphism $R\to C$, forcing $C=R$. We may thus assume $L=W$.

Lift $c^{-1}$ to $a\in C$ and set $z=ca$. Since $z-1\in T$, some
$c^n$ annihilates $z-1$, whence $z^n(z-1)=0$. Increasing $n$ to be
positive, $e=z^n$ is an idempotent. It has rational image one and
annihilates $T$: the equality $e=z^N$ holds for every $N\ge n$, and
each element of $T$ is killed by some $c^N$. Thus $eC=W$, canonically;
indeed $c(ea)=e$. The element $1-e\in T$ is killed by some $J^M$, so
$D=(1-e)C$ satisfies $J^MD=0$. The epic map $R\to D$ factors through
$R/J^M$, whose maximal ideal is nilpotent. Applying the same
nilpotent-ideal epimorphism argument to its image gives
\[
                          C=W\times R/I,
             \qquad J^M\subset I\subset R,\quad M\ge1,  \tag{B78.6}
\]
with the diagonal inclusion of $R$.

We need an exact coefficient description and an ideal norm.
Choose a set $B\subset R$ of representatives of $k$, and put
$P(X)=\prod_{b\in B}(X-b)$. Differences between distinct elements of
$B$ are units of $R$, $P(R)\subset J$, and $P'(b)\in R^\times$
for every $b\in B$. The monic basis $X^iP^j$, $0\le i<q$, $j\ge0$, gives
\[
 \Omega(R,K)=
 \bigoplus_{0\le i<q}RX^i\ \oplus\!
 \bigoplus_{\substack{0\le i<q\\j\ge1}}
                (J:_KJ^j)X^iP^j.                         \tag{B78.7}
\]
To prove necessity, first subtract the polynomial of degree below $q$
interpolating the values on $B$; its coefficients are in $R$.
For the resulting $F=\sum_{i<q,\,j\ge1}a_{ij}X^iP^j$ and any
$0\ne y\in J$, let $\delta=\min_{i,j}v(a_{ij}y^j)$.
For $b\in B$, substitute $X=b+yZ$ and divide by a coefficient of
valuation $\delta$. Modulo the maximal ideal of $V$, the resulting
polynomial is
\[
             \sum_{j\ge1}\left(\sum_{i<q}
                  \overline{a_{ij}y^j/\mu}\,\bar b^i\right)
                       (\overline{P'(b)}Z)^j,
                      \qquad v(\mu)=\delta.              \tag{B78.8}
\]
For some $b$ this is nonzero by the invertible Vandermonde matrix on
$k$ and the nonzero residues of $P'(b)$, and it has zero constant term. If $\delta<0$, a residue outside
its finite zero set gives a value of $F$ outside $V$. If $\delta=0$,
its values on the infinite field $l$ would have to lie in the finite
set $\bar\mu^{-1}k$, which is impossible for a nonconstant polynomial.
If $\delta>0$ but $\delta\notin v(J)$, a nonzero residue gives a
positive-valued element of $R$ outside $J$, also impossible.
Thus every $a_{ij}y^j$ lies in $J$. Products of $j$ elements of $J$
are divisible by the $j$th power of a factor of smallest valuation;
hence $a_{ij}J^j\subset J$. This proves necessity; sufficiency is
immediate from $P(R)\subset J$. The coefficient ideals in (B78.7)
all lie in $W$, since $c^j a_{ij}\in J\subset R$.

For a polynomial $F$ vanishing on $B$, define its $J$-norm to be the
$V$-ideal
\[
                   \nu_J(F)=\sum_{i<q,\,j\ge1}a_{ij}J^j.
                                                        \tag{B78.9}
\]
Write $F(b+Y)=\sum_{h\ge1}g_{bh}Y^h$ for $b\in B$.
The same reduction (B78.8), now considered for each fixed $y\in J$,
proves equality of the principal ideals of least coefficient value
after multiplication by the corresponding powers of $y$. Taking the
union over $y\in J$, which computes powers of a valuation ideal,
gives
\[
                         \nu_J(F)=\sum_{b\in B,h\ge1}g_{bh}J^h.
                                                        \tag{B78.10}
\]
This equality uses actual ideals rather than a numerical infimum.

If $I=R$ in (B78.6), then $C=W$ and (B78.7) is an $R$-module direct
sum for $\Omega(R,C)$. Its coefficient module of $P$ is
$U=(J:J)$, with $JU=J$. Choose $\beta\in V$ whose residue is not in
$k$. The tensor $\beta P\otimes P-P\otimes\beta P$ has zero polynomial
image, but its projection to $U\otimes_R U$ is nonzero. Indeed modulo
$J$ it becomes $\bar\beta\otimes1-1\otimes\bar\beta$ in
$(U/J)\otimes_k(U/J)$, which is nonzero since $\bar\beta\notin k$.
This contradicts injectivity.

Finally suppose $I\subsetneq R$ in (B78.6). Define, with the indicated
value conditions,
\[
 H=\{F\in W[X]:F(R)\subset I\},\qquad E_0=R[X]+H,
 \qquad H_0=\{F\in H:F|_B=0\}.
                                                        \tag{B78.11}
\]
The first projection of $\Omega(R,C)$ is exactly $E_0$: lift the
coefficients of its second component to $R[X]$ and subtract.
Set $K_0=\sum_{F\in H_0}\nu_J(F)$, a $V$-ideal. Formula (B78.7)
and the polynomials $aP$, $a\in J^{M-1}$, give
\[
                            J^M\subset K_0\subset J.    \tag{B78.12}
\]
We claim $H\not\subset JE_0$. Suppose the contrary and fix $N\ge1$.
For $F\in H_0$, $b\in B$, and $\lambda\in(J:J^N)$, the substitution
$\psi=b+\lambda P^N$ preserves $R$ and sends $F$ into $H_0$.
All its nonzero $P$-degrees are at least $N$. Moreover, (B78.2), the
fact that products of a valuation ideal are generated by pure powers,
and (B78.10) give
\[
 \sum_{b\in B,\,\lambda\in(J:J^N)}\nu_J(F\circ\psi)
       =\sum_{b,h}g_{bh}(J:J^N)^hJ^{Nh}
       =\sum_{b,h}g_{bh}J^h=\nu_J(F).                    \tag{B78.13}
\]
The coefficients $\lambda$ belong to $W$, since $c^N\lambda\in J$.
By the assumed containment write
$F\circ\psi=A+\sum r_tF_t$ with $A\in JR[X]$, $r_t\in J$,
and $F_t\in H$. Subtracting the interpolants on $B$ from all $F_t$
absorbs only elements of $JR[X]$ into $A$, so we may take $F_t\in H_0$.
Project to $P$-degrees at least $N$. The norm of this part of $A$ is
contained in $J^{N+1}$, while the remaining terms have norm contained
in $JK_0$. Consequently (B78.13) implies
$K_0\subset J^{N+1}+JK_0$. Taking $N=M$ and using (B78.12) gives
$K_0=JK_0$. Iterating and using $K_0\subset J$ gives
$J^M\subset K_0\subset J^{M+1}$, so $J^M=J^{M+1}$.
This forces $J=J^2$: for $r\in J$, the equality
$r^M\in J^M=J^{2M}$ provides a product of $2M$ elements of $J$
of valuation at most $Mv(r)$; its least-valued factor $y$ satisfies
$2v(y)\le v(r)$, hence $r\in y^2V\subset J^2$.
The contradiction proves the claim.

Choose $F\in H\setminus JE_0$. In $\Omega(R,C)$ the two elements
\[
                       f=(F,0),\qquad g=(0,P^M\bmod I)
\]
both survive modulo $J\Omega(R,C)$. For $f$ this follows by first
projection; for $g$ it follows from its monic coefficient and
$I\subset J$. Their pure tensor is therefore nonzero over the residue
field $R/J=k$. But its polynomial image is zero, since the two
factors lie in the distinct product components in (B78.6).
This final contradiction excludes every $C\ne R$.
Since $(R:_S J)\subset C$ and (B78.5) is obtained by iterating this
colon operation, $C=R$ is equivalent to the ideal condition in
(B78.1). The proof and all positive tensor assertions are complete.
\end{proof}

\begin{saresult}{B.079}{Proved}{Coefficient base change and an explicit monic saturation obstruction}
For an arbitrary unital inclusion $R\subset S$, let
$C=\bigcup_d(R:_S D_d(R))$ be the effective coefficient ring of B.064,
and put $Q=C/R$. Define the $R$-module directly by finite coefficient
relations
\[
 N(R,Q)=\left\{(q_0,q_1,\ldots)\in\bigoplus_{i\ge0}Q:
             \sum_i r^i q_i=0\text{ for every }r\in R\right\}.
                                                        \tag{B79.1}
\]
If the polynomial tensor map $\theta_2$ is injective, then $R\to C$
is epic and
\[
 C\otimes_R\Omega(R,S)\ \xrightarrow{\ \sim\ } C[X],
 \qquad C\otimes_RN(R,Q)=0.                              \tag{B79.2}
\]
The two assertions in (B79.2) are equivalent, without assuming
injectivity of $\theta_2$. They are necessary conditions, not a
universal sufficient criterion. Even when $C$ is a flat localization
of a domain $R$, they do not imply injectivity of $\theta_2$ or
injectivity of multiplication by monic polynomials on its source.
\end{saresult}
\begin{proof}
Write $A=\Omega(R,S)$. By B.064 every coefficient of $A$ lies in $C$,
and injectivity of $\theta_2$ makes $R\to C$ epic.
The multiplication map $\mu:C\otimes_R A\to C[X]$ is always surjective,
because $R[X]\subset A$. Suppose that
$z=\sum_{i=1}^m c_i\otimes F_i$ lies in its kernel.
For each $c_i$, B.064 supplies a monic $P_{d_i}\in R[X]$ such that
$c_iP_{d_i}\in A$. Their product $H=\prod_iP_{d_i}$ is monic and
$c_iH\in A$ for every $i$. The tensor
\[
                         w=\sum_i c_iH\otimes F_i
\]
has polynomial image $H(X)\sum_i c_iF_i(Y)=0$. Injectivity gives $w=0$.
Apply extraction of the coefficient of $X^{\deg H}$ in the first
factor, as an $R$-linear map $A\to C$. The result is $z=0$, proving
the first isomorphism in (B79.2).

Reduction of coefficient vectors modulo $R$ gives an exact sequence
\[
                0\longrightarrow R[X]\longrightarrow A
                    \longrightarrow N(R,Q)\longrightarrow0.  \tag{B79.3}
\]
Surjectivity on the right follows by lifting the finitely many
coefficients of a vector in (B79.1): its defining equations say
precisely that the lifted polynomial has values in $R$.
After tensoring, the composite
$C[X]=C\otimes_RR[X]\to C\otimes_RA\to C[X]$ is the identity.
Thus the first arrow is split injective and its cokernel is
$C\otimes_RN(R,Q)$. The second arrow is an isomorphism exactly when
this cokernel vanishes, proving the equivalence claimed.

Here is the explicit obstruction to sufficiency and to monic
saturation. Let $k=\mathbb F_q\subset l$ with $l$ infinite, let
$V=l[[t]]$, $K=l((t))$, $R=k+tV$, and $S=K$. Put $P=X^q-X$ and choose
$\beta\in l\setminus k$. Formula (B77.4), equivalently (B78.7), gives
the direct $R$-module decomposition
\[
 A=\bigoplus_{i<q}RX^i\ \oplus\!
      \bigoplus_{\substack{i<q\\j\ge1}}t^{1-j}V X^iP^j.
                                                        \tag{B79.4}
\]
The effective coefficient ring is $K$: for every $n\ge0$ and
$v\in V$, the polynomial $t^{-n}vP^{n+1}$ has values in $tV\subset R$.
Consequently $C=K$ is a flat epic localization, and
$K\otimes_RA=K[X]$.
Nevertheless
\[
                       z=\beta P\otimes P-P\otimes\beta P
                                                        \tag{B79.5}
\]
is a nonzero tensor with zero polynomial image. Its projection using
the coefficient of $P$ in (B79.4), followed by reduction modulo $tV$,
is $\beta\otimes1-1\otimes\beta$ in $l\otimes_k l$, which is nonzero.
For example choose a $k$-linear functional on $l$ that vanishes on
$1$ and takes value one on $\beta$.
But $P^2/t\in A$, since $P(R)\subset tV$, and $t,t\beta\in R$.
The balancing relations in the tensor product give
\[
 \beta P^2\otimes P
       =(P^2/t)\otimes(t\beta P)
       =P^2\otimes\beta P.
                                                        \tag{B79.6}
\]
Therefore $P(X)z=0$. This is a concrete nonzero element killed by a
monic polynomial in the first variable, despite flat epic coefficient
base change and both conditions (B79.2). It prevents that necessary
condition from closing the unrestricted structural question.
\end{proof}

\begin{saresult}{B.080}{Proved under stated hypotheses}{Canonical arithmetic Tor obstructions for a flat effective coefficient ring}
Let $R\subset S$ be any unital commutative inclusion. Form the explicit
coefficient ring $C=\bigcup_d(R:_S D_d(R))$ of B.064, and assume that
$R\to C$ is flat and epic. Put $Q=C/R$ and define the $R$-module
\[
 E=\left\{\left(r\longmapsto\sum_{i=0}^d c_i r^i+R\right):
                       d\ge0,\ c_i\in C\right\}
       \subset\operatorname{App}(R,Q).                   \tag{B80.1}
\]
Thus $E$ is specified by coefficient and value relations in $R\subset C$,
without choosing tensor presentations of integer-valued polynomials.
For $A=\Omega(R,S)$ there is a canonical exact sequence
\[
\begin{gathered}
 0\longrightarrow\operatorname{Tor}_2^R(E,E)
 \longrightarrow A\otimes_R A\\
 \xrightarrow{\theta_2} C[X,Y]
 \xrightarrow{\partial}\operatorname{Tor}_1^R(E,E)
 \longrightarrow0.
\end{gathered}                                          \tag{B80.2}
\]
The map $\partial$ is given explicitly below, with all connecting maps
fixed by the evaluation sequence. In particular the two tensors on
153/004 exist uniquely exactly when
\[
\begin{gathered}
 \operatorname{Tor}_2^R(E,E)=0,\\
 \partial\!\left(\sum_{i,j\ge0}\binom{i+j}{i}
                         a_{i+j}X^iY^j\right)=0,\\
 \partial\!\left(\sum_{i\ge0}a_iX^iY^i\right)=0
                                                        \tag{B80.3}
\end{gathered}
\]
for every finite vector $(a_i)$ in $C$ satisfying
$\sum_i a_i r^i\in R$ for all $r\in R$; missing coefficients are zero.
If only flatness of $C$ is initially assumed, epicness is an additional
necessary condition, and epicness together with (B80.3) is sufficient.
This is an exact obstruction calculation for the flat-coefficient
scope, not a classification of arbitrary nonflat coefficient extensions.
\end{saresult}
\begin{proof}
Let $\mathcal F=C[X]$, and let
$\varepsilon:\mathcal F\to E$ send a coefficient vector to its
function modulo $R$. By the definition (B80.1),
\[
                0\longrightarrow A\longrightarrow\mathcal F
                    \xrightarrow{\varepsilon}E\longrightarrow0
                                                        \tag{B80.4}
\]
is exact. Also $E$ is a quotient of $Q[X]$, since every ordinary
$R$-coefficient polynomial maps to zero. Epicness gives
$C\otimes_R C\simeq C$, and the sequence $R\to C\to Q\to0$ then
gives $C\otimes_R Q=0$. Therefore
\[
 C\otimes_R E=0,\qquad \mathcal F\otimes_R E=0.           \tag{B80.5}
\]
Since $C$ is flat, so is $\mathcal F$. Tensoring (B80.4) with
$\mathcal F$ gives an isomorphism
\[
 \mathcal F\otimes_R A\ \xrightarrow{\sim}
 \mathcal F\otimes_R\mathcal F\simeq C[X,Y].             \tag{B80.6}
\]
Here the last map multiplies coefficients, using the ring epimorphism;
the variables in the two factors remain distinct. Explicitly the inverse
sends $H=\sum_j G_j(X)Y^j$ to $\sum_jG_j(X)\otimes Y^j$, since
$R[Y]\subset A$.

We use the long exact Tor sequence with its usual connecting maps.
For clarity, it is obtained by tensoring a free resolution with the
short exact sequence and applying the connecting homomorphisms of the
resulting degreewise exact sequence of complexes; this fixes all signs.
Flatness makes every positive Tor group with $\mathcal F$ vanish.
Applying this construction to (B80.4), first with the module $A$,
gives
\[
 0\longrightarrow\operatorname{Tor}_1^R(E,A)
 \longrightarrow A\otimes_R A
 \longrightarrow\mathcal F\otimes_R A
 \longrightarrow E\otimes_R A\longrightarrow0.          \tag{B80.7}
\]
Applying it with $E$ and using (B80.5) gives the canonical isomorphisms
\[
\begin{aligned}
 \operatorname{Tor}_2^R(E,E)&\xrightarrow{\sim}
                         \operatorname{Tor}_1^R(E,A),\\
 \tau:\operatorname{Tor}_1^R(E,E)&\xrightarrow{\sim}E\otimes_R A.
\end{aligned}                                           \tag{B80.8}
\]
Combining (B80.6)--(B80.8) gives (B80.2). These uses of Tor agree
with the opened Stacks Section 10.75, Lemmas 10.75.2, 10.75.5 and
10.75.8 \cite{StacksTor}; the construction above specifies the maps actually used.

The promised formula, containing only the displayed coefficient array
and the fixed connecting isomorphism, is
\[
 \partial(H)=\tau^{-1}
       \left(\sum_j\varepsilon(G_j(X))\otimes Y^j\right),
             \qquad H=\sum_jG_j(X)Y^j.                  \tag{B80.9}
\]
Every term $Y^j$ here is an ordinary monomial in the specified copy
$R[Y]\subset A$. Thus no unknown co-operation or decomposition of
$H$ into integer-valued factors enters the definition of $\partial$.
One can compute its class with any free resolution of the explicitly
defined function module $E$; the comparison maps between resolutions
identify the resulting classes. Every individual relation or boundary
in such a computation uses finitely many coefficients. This statement
does not claim finite generation of $E$, a uniform degree bound, or a
terminating decision algorithm for arbitrary presented rings.

The kernel in (B80.2) proves that uniqueness is equivalent to the
first condition in (B80.3). The other two conditions say, by the exact
sequence, exactly that the binomial coefficient array of $F(X+Y)$
and the diagonal coefficient array of $F(XY)$ lift. When the kernel
vanishes their lifts are unique, so (B80.9) recovers precisely the two
specified polynomial tensors. If epicness was not initially assumed,
B.064 proves it from uniqueness, giving the variant in the statement.
No higher tensor injectivity or plethory identities are inferred from
this binary calculation alone.

There is also an intrinsic distinction between the two specified
existence tests and full bivariate polynomial compositeness.
Let $\varepsilon_2:C[X,Y]\to\operatorname{App}(R^2,Q)$ be evaluation
modulo $R$. It vanishes on the image of $\theta_2$, so (B80.2)
defines a unique map
\[
 \lambda:\operatorname{Tor}_1^R(E,E)\longrightarrow
                  \operatorname{App}(R^2,Q),\qquad
                         \lambda\partial=\varepsilon_2. \tag{B80.10}
\]
Writing $\Omega_2=\ker\varepsilon_2$, exactness gives
\[
\begin{gathered}
 0\longrightarrow\operatorname{Tor}_2^R(E,E)
 \longrightarrow A\otimes_R A\\
 \longrightarrow\Omega_2
 \longrightarrow\ker\lambda\longrightarrow0.
\end{gathered}                                          \tag{B80.11}
\]
Indeed every class in $\ker\lambda$ has a preimage under the
surjective map $\partial$, and that preimage lies in $\Omega_2$.
Thus the full bivariate tensor map is an isomorphism exactly when
$\operatorname{Tor}_2^R(E,E)=0$ and $\lambda$ is injective. This
identification makes the remaining existence obstruction explicit
without assuming that the stronger condition follows from (B80.3).
The nonflat example B.067 is outside the flat-coefficient hypothesis,
so the present calculation does not close the original arbitrary-ring
structural task.
\end{proof}

\begin{saresult}{B.081}{Proved under stated hypotheses}{Self-injective finite approximations and the unbounded nilpotent test ring}
Every injective ring epimorphism $R\to C$ is an isomorphism in either
of the following situations:
\begin{enumerate}
\item $R$ is a directed union of unital self-injective subrings, where
self-injective means injective as a module over itself.
\item $R$ is local with nil maximal ideal $\mathfrak m$, and it has
an ideal $I$ such that $\mathfrak m^N I=0$ for some finite $N$, while
every injective ring epimorphism out of $R/I$ is an isomorphism.
\end{enumerate}
For any such base and any unital inclusion $R\subset S$, the polynomial
co-operations exist uniquely exactly when
\[
 (R:_S D_d(R))=R\quad\text{for all }d\ge1,
 \quad\text{equivalently}\quad \Omega(R,S)=R[X].          \tag{B81.1}
\]
Every positive case has all multivariable tensor isomorphisms.
In particular this completely decides all coefficient extensions of
the previously specified test base
\[
 R=\mathbb F_2[z,x_1,x_2,\ldots]/
       (z^2,\ zx_i\ (i\ge1),\ x_i^{2^{i+1}}\ (i\ge1)).   \tag{B81.2}
\]
This base is local of dimension zero and non-Noetherian; its maximal
ideal is not T-nilpotent and its elementwise nilpotence has no uniform
bound. Thus the conclusion does not follow from those earlier
bounded-nilpotence cases.
\end{saresult}
\begin{proof}
For the first assertion write $R=\bigcup_\alpha R_\alpha$ as a directed
union, and suppose $R\subset C$ is epic. For $c\in C$ the equality
$c\otimes1=1\otimes c$ holds in $C\otimes_R C$. The balancing
relations defining this tensor product are the union of those for
the $R_\alpha$: every individual equality is a finite sum of
balancing relations and therefore uses coefficients from one
$R_\alpha$. Consequently
\[
          c\otimes1=1\otimes c
                   \quad\text{in }C\otimes_{R_\alpha}C
                                                        \tag{B81.3}
\]
for some $\alpha$. Since $R_\alpha$ is self-injective, its inclusion
as an $R_\alpha$-submodule of $C$ has an $R_\alpha$-linear retraction
$\rho:C\to R_\alpha$. In particular $\rho(1)=1$.
Applying $\rho\otimes1$ to (B81.3) gives $\rho(c)=c$ in $C$.
Hence $c\in R_\alpha\subset R$, proving $C=R$. This argument uses
neither flatness of the transition maps nor a module-finiteness
assumption on $C$.

For the second assertion, let again $R\subset C$ be epic. The ideal
$\mathfrak mC$ is nil: each element is a finite sum of multiples of
nilpotent elements, and the ideal generated by any finite collection
of such elements is nilpotent in a commutative ring. Hence
$C/\mathfrak mC$ is nonzero. Base change gives an epimorphism from
the field $R/\mathfrak m$ to this nonzero algebra, so it is an
isomorphism. For completeness a field extension algebra $B$ with
$B\otimes_kB\simeq B$ must equal $k$: for any $b\notin k$ a
$k$-linear functional vanishing on $1$ and nonzero on $b$ contradicts
$b\otimes1=1\otimes b$. We conclude $C=R+\mathfrak mC$.
Multiplying by $I$ and iterating this equality gives
\[
                 IC=I+\mathfrak mIC
                    =I+\mathfrak m^NIC=I.               \tag{B81.4}
\]
Thus $I$ is an ideal of $C$ with precisely its original underlying
subset of $R$. The induced map $R/I\to C/I$ is therefore an
injective ring epimorphism. By hypothesis it is an isomorphism, and
then $C=R+I=R$.

For any base with this epimorphism property, B.064 applies to the
explicit effective coefficient ring $C=\bigcup_d(R:_S D_d)$.
Tensor injectivity forces $R\to C$ epic, hence $C=R$. Conversely
$C=R$ gives ordinary polynomials in all arities. This proves (B81.1)
and all positive tensor assertions.

Finally put
\[
 A_n=\mathbb F_2[x_1,\ldots,x_n]/
                (x_i^{2^{i+1}}:1\le i\le n),
 \qquad A=\bigcup_{n\ge1}A_n.                            \tag{B81.5}
\]
The transition maps are the evident injections on the monomial
bases. Every $A_n$ is self-injective. Indeed the functional taking
the coefficient of the top monomial
$\prod_{i=1}^n x_i^{2^{i+1}-1}$ gives a nondegenerate bilinear form
$(a,b)\mapsto\ell(ab)$: a monomial pairs with its complementary
monomial, and the resulting matrix is a permutation matrix.
It identifies $A_n$ with $\operatorname{Hom}_{\mathbb F_2}
(A_n,\mathbb F_2)$ as an $A_n$-module. This dual module is injective,
since for every $A_n$-module $M$ there is a natural identification
\[
 \operatorname{Hom}_{A_n}
       (M,\operatorname{Hom}_{\mathbb F_2}(A_n,\mathbb F_2))
       \simeq\operatorname{Hom}_{\mathbb F_2}(M,\mathbb F_2),
\]
and vector-space duality is exact. The first part of the theorem
therefore proves that all injective epimorphisms out of $A$ are
isomorphisms.

The ring (B81.2) is $A\oplus\mathbb F_2z$, where $z^2=0$ and the
augmentation ideal of $A$ annihilates $z$. Its maximal ideal
$\mathfrak m=(z,x_1,x_2,\ldots)$ is nil, since any element uses only
finitely many variables, and the corresponding finite algebra has
nilpotent maximal ideal. Thus $R$ is local of dimension zero.
The ideal $I=\mathbb F_2z$ satisfies $\mathfrak mI=0$ and $R/I=A$,
so the second part applies. The classes of all $x_i$ are independent
in $\mathfrak m/\mathfrak m^2$, proving non-Noetherianity. Each product
$x_1\cdots x_n$ is a nonzero basis monomial, proving failure of
T-nilpotence. The exact nilpotence orders $2^{i+1}$ of the $x_i$ are
unbounded, which checks the remaining assertions about the test ring.
\end{proof}

\begin{saresult}{B.082}{Proved}{A universal coefficient-syzygy criterion for the two polynomial tensors}
Let $R\subset S$ be any unital inclusion of commutative rings. The
following construction uses only ideals, coefficient functions and
finite linear relations associated with this inclusion. It imposes
no finiteness, reducedness, flatness or domain hypothesis.

First form
\[
\begin{gathered}
 C=\bigcup_{d\ge1}(R:_S D_d(R)),\qquad \mathcal F=C[X],\\
 E=\operatorname{im}\left(\mathcal F\xrightarrow{\varepsilon}
                       \operatorname{App}(R,C/R)\right),\\
 \varepsilon\left(\sum c_iX^i\right)(r)=\sum c_ir^i+R.
\end{gathered}
                                                        \tag{B82.1}
\]
The first arithmetic test is that every $c\in C$ have a finite
matrix presentation
\[
\begin{gathered}
 c=Y M Z,\qquad M\in\operatorname{Mat}_{n}(R),\\
 Y\in C^{1\times n},\quad Z\in C^{n\times1},\quad
 YM\in R^{1\times n},\quad MZ\in R^{n\times1}.
\end{gathered}                                           \tag{B82.2}
\]
This test is equivalent to epicness of $R\to C$.
If it fails, the requested tensors are not uniquely determined.
When it holds, define the following finite-depth syzygy invariants.
Set
\[
\begin{aligned}
 U_q&=\operatorname{Tor}_q^R(\mathcal F,\mathcal F),\\
 V_q&=\operatorname{Tor}_q^R(\mathcal F,E)
                    \oplus\operatorname{Tor}_q^R(E,\mathcal F),\\
 W_q&=\operatorname{Tor}_q^R(E,E),\qquad
 U_q\xrightarrow{a_q}V_q\xrightarrow{b_q}W_q,
\end{aligned}                                           \tag{B82.3}
\]
where
\[
\begin{aligned}
 a_q&=(\operatorname{Tor}_q(1,\varepsilon),
                         \operatorname{Tor}_q(\varepsilon,1)),\\
 b_q&=\operatorname{Tor}_q(\varepsilon,1)
                        -\operatorname{Tor}_q(1,\varepsilon).
\end{aligned}
                                                        \tag{B82.4}
\]
The two terms of $b_q$ act on the respective summands of $V_q$;
thus $b_qa_q=0$. Put
\[
 Z_q=\ker a_q,\quad T_q=\operatorname{coker}b_q,\quad
 K_1=\ker b_1/\operatorname{im}a_1.                       \tag{B82.5}
\]
There are canonical maps, explicitly constructed in the proof,
\[
\begin{gathered}
 \delta_0:C[X,Y]=Z_0\longrightarrow T_1,
 \qquad \delta_1:Z_1\longrightarrow T_2,\\
 K_2=T_2/\operatorname{im}\delta_1.
\end{gathered}                                           \tag{B82.6}
\]
The two polynomial tensors on 153/004 exist uniquely if and only if
(B82.2) holds and
\[
\begin{gathered}
                         K_1=K_2=0,\\
 \delta_0\left(\sum_{i,j\ge0}\binom{i+j}{i}
                        a_{i+j}X^iY^j\right)=0,\\
 \delta_0\left(\sum_{i\ge0}a_iX^iY^i\right)=0            \tag{B82.7}
\end{gathered}
\]
for every finite coefficient vector $(a_i)$ in $C$ satisfying
$\sum_i a_ir^i\in R$ for every $r\in R$. Missing coefficients are zero.
All quantities in this test are defined before forming any tensor
power of $\Omega(R,S)$; in particular it is not a test that presupposes
the image or kernel of that tensor map. Its free-module implementation
below uses only finite coefficient relations in homological degrees
at most three. The criterion does not assert a uniformly terminating
algorithm for arbitrary presentations of infinite rings, and it
does not replace the separate higher-coherence conditions for a
plethory by a binary test.
\end{saresult}
\begin{proof}
We first check the matrix test and then construct the invariants and
prove their meaning by explicit chain calculations.
If (B82.2) holds, balancing its entries shows
$c\otimes1=1\otimes c$ in $C\otimes_R C$: use $YM$ to move the first
sum to the second factor, and then $MZ$ to move it back.
Conversely suppose this tensor equality holds. A finite presentation
of that equality gives finitely many $s_i\in C$, with labeled
$s_0=1,s_1=c$, and a finite matrix $(m_{ij})$ over $R$ such that
\[
 \sum_i m_{i0}s_i=c,\qquad \sum_i m_{i1}s_i=-1,
 \qquad \sum_i m_{ij}s_i=0\ (j\ne0,1),\qquad
                       \sum_jm_{ij}s_j=0\ (\text{all }i).
                                                        \tag{B82.8}
\]
Indeed present the second copy of $C$ by the free module on its
elements. The vector having entries $c$ at $s_0$, $-1$ at $s_1$,
and zero elsewhere lies in the image of its relation module tensored
with the first copy of $C$. Only finitely many relations and elements
are needed, giving (B82.8), after enlarging the common index set.
Remove the row and column indexed by zero, and call the remaining
matrix $M'$, with row $Y'=(s_i)_{i\ne0}$ and column
$Z'=(s_j)_{j\ne0}$. Then $Y'M'$ and $M'Z'$ have entries in $R$,
and $Y'M'Z'=m_{00}-c$. Appending a one-by-one block $m_{00}$ to
$-M'$, and a coordinate $1$ to both $Y'$ and $Z'$, gives (B82.2).
This proves the equivalence, including the finite certificate.
It is the opened Stacks Lemma 10.107.11 and Remark 10.107.12
\cite{StacksEpicMatrices};
the argument also supplies the final elementary step omitted there.
B.064 now proves that failure of this test precludes uniqueness.

Assume henceforth (B82.2). Since $C\otimes_R C\simeq C$, one has
$C\otimes_R(C/R)=0$. The module $E$ is a quotient of $(C/R)[X]$,
and $\mathcal F\to E$ is surjective. Right exactness therefore gives
\[
 \mathcal F\otimes_R E=E\otimes_R\mathcal F=0,
                         \qquad E\otimes_R E=0.          \tag{B82.9}
\]
In particular $V_0=W_0=0$, while
$U_0=\mathcal F\otimes_R\mathcal F=C[X,Y]$ by coefficient
multiplication. Hence $Z_0=U_0$ as claimed.

Here is a fixed, functorial free construction, so that the
transgressions in (B82.6) are unambiguous. For any $R$-module $M$,
let $P_0(M)$ be the free $R$-module on its underlying set, augmented
by $[m]\mapsto m$. Take $P_1(M)$ to be free on the elements of the
kernel, with its evident map to $P_0(M)$, and repeat on successive
kernels. A module map sends the indicated generators to their
images at each stage, making $P_\bullet$ functorial. Every generator
of a kernel is a finite linear relation in the preceding free
module. Thus this construction is explicit even without finite
generation or any choice of a tensor presentation of $\Omega$.

Write $P=P_\bullet(\mathcal F)$, $Q=P_\bullet(E)$, and
$u=P_\bullet(\varepsilon):P\to Q$. For nonnegative vertical degree
$q$, form three complexes of free $R$-modules
\[
\begin{aligned}
 \mathcal T_{0,\bullet}&=\operatorname{Tot}(P\otimes_RP),\\
 \mathcal T_{1,\bullet}&=\operatorname{Tot}(P\otimes_RQ)
                         \oplus\operatorname{Tot}(Q\otimes_RP),\\
 \mathcal T_{2,\bullet}&=\operatorname{Tot}(Q\otimes_RQ).
\end{aligned}                                           \tag{B82.10}
\]
Here $\operatorname{Tot}(P\otimes Q)_q=
\bigoplus_{a+b=q}P_a\otimes_RQ_b$, with vertical differential
$d=d_P\otimes1+(-1)^a1\otimes d_Q$ on that summand.
Define the horizontal chain maps
\[
 h_0=(1\otimes u,u\otimes1),\qquad
 h_1=(u\otimes1,-1\otimes u).                            \tag{B82.11}
\]
They commute with $d$ and satisfy $h_1h_0=0$.
Their induced maps on vertical homology are exactly $a_q,b_q$.
The vertical homology identifications with Tor follow by the free
double-resolution argument, as in opened Stacks Lemma 10.75.5
\cite{StacksTor}.

For $q=0,1$ and $[z]\in Z_q$, choose a vertical cycle
$z\in\mathcal T_{0,q}$ representing it. Since $a_q[z]=0$,
there exists $w\in\mathcal T_{1,q+1}$ such that
\[
                dw=(-1)^{q+1}h_0z.
\]
The element $(-1)^{q+1}h_1w$ is a vertical cycle in
$\mathcal T_{2,q+1}$. Define
\[
        \delta_q[z]=[(-1)^{q+1}h_1w]
                  \pmod{\operatorname{im}b_{q+1}}\in T_{q+1}.
                                                        \tag{B82.12}
\]
Changing $w$ by a vertical cycle changes this class by an image of
$b_{q+1}$. Changing $z$ to $z+dv$, with
$v\in\mathcal T_{0,q+1}$, permits replacing $w$ by
$w+(-1)^{q+1}h_0v$, which leaves $h_1w$ unchanged.
This proves independence of choices and $R$-linearity. Standard
comparison maps, or the same calculation with homotopic chain maps,
identify these definitions for any other free resolutions; the fixed
functorial resolution already defines them without that ambiguity.
For $q=0$ an arbitrary coefficient polynomial $H\in C[X,Y]$ is
represented by any preimage in $\mathcal T_{0,0}$, and (B82.12)
specifies its obstruction using one finite chain lift $w$.

We now identify what these invariants measure. Give
$\mathcal T_{p,q}$ total homological degree $q-p$, and set
\[
                      D=d+(-1)^q h.                     \tag{B82.13}
\]
The commutation $dh=hd$ and $h^2=0$ give $D^2=0$.
This total complex has homology in degree $n$ equal to
$\operatorname{Tor}_n^R(A,A)$, where $A=\ker\varepsilon=\Omega(R,S)$,
and zero homology in negative degrees. We verify this assertion
without a convergence assumption. The two-column total complex
$L$ formed from $P\xrightarrow{u}Q$, in degrees $q-p$ with
$p=0,1$, has $L_n=P_n\oplus Q_{n+1}$ and differential (B82.13).
The exact sequence $0\to A\to\mathcal F\to E\to0$ shows
$H_0(L)=A$ and all other homology zero. Also $L_n=0$ for $n<-1$,
and $L_{-1}=Q_0$ is free. The surjection $L_0\to L_{-1}$ therefore
splits, separating a contractible two-term complex. The remainder
is a nonnegative projective resolution of $A$: its degree-zero
module is a direct summand of the free module $L_0$, and its
higher modules are free. Tensoring two copies of this decomposition
shows that $L\otimes_RL$ computes $\operatorname{Tor}_\bullet^R(A,A)$.

There is an explicit sign isomorphism from $L\otimes_RL$ to the
three-column complex (B82.10)--(B82.13). On the component whose
first factor has column $p_1\in\{0,1\}$ and whose second factor
has vertical resolution degree $b$, multiply by $(-1)^{p_1b}$.
For the second vertical differential this changes its sign
$(-1)^{a-p_1}$ to $(-1)^a$; for the first horizontal differential
it changes $(-1)^a$ to $(-1)^{a+b}$; for the second horizontal
differential it leaves $(-1)^{a+b-p_1}$, which is the sign in
(B82.11)--(B82.13). The first vertical sign is unchanged.
This verifies the chain isomorphism term by term and proves the
homology assertion. In particular
\[
                 H_0(\operatorname{Tot}\mathcal T)=A\otimes_RA,
                 \qquad H_{-1}(\operatorname{Tot}\mathcal T)=0.
                                                        \tag{B82.14}
\]
The edge map obtained by taking the column-zero component and
then its vertical homology is the natural map
$A\otimes_RA\to\mathcal F\otimes_R\mathcal F=C[X,Y]$.
Indeed $H_0(L)=A$ is identified with the kernel of $\varepsilon$
by the augmentation to $\mathcal F$. Thus this edge map is
precisely $\theta_2$.

For completeness the low-degree exactness can now be checked
directly. A total degree-zero cycle consists of
\[
 x_0\in\mathcal T_{0,0},\quad
 x_1\in\mathcal T_{1,1},\quad
 x_2\in\mathcal T_{2,2},\qquad
             dx_1=-h_0x_0,\quad dx_2=h_1x_1.             \tag{B82.15}
\]
It follows from (B82.12) that its edge value lies in
$\ker\delta_0$. Conversely, if $\delta_0[x_0]=0$, choose $x_1$
solving the first equation. The zero class of $-h_1x_1$ modulo
$\operatorname{im}b_1$ allows adding a vertical cycle to $x_1$
so that $h_1x_1$ becomes a vertical boundary; then choose $x_2$
solving the second equation. Therefore
$\operatorname{im}\theta_2=\ker\delta_0$.
If $y\in\mathcal T_{2,1}$ is a vertical cycle, it is a total
degree-minus-one cycle. By (B82.14), it is the boundary of a
triple $(x_0,x_1,x_2)$ of total degree zero. The column-one
equation is $dx_1=-h_0x_0$, and the column-two equation is
$-h_1x_1+dx_2=y$. Thus $\delta_0[x_0]=[y]$ in $T_1$.
Every class in $T_1$ has such a representative, proving surjectivity
and the canonical cokernel identification
\[
                    \operatorname{coker}\theta_2\simeq T_1.
                                                        \tag{B82.16}
\]

If a total degree-zero cycle has zero edge value, its $x_0$ is a
vertical boundary. Subtracting a total boundary makes $x_0=0$.
Then $x_1$ is a vertical cycle, and (B82.15) says that its class
lies in $\ker b_1$. Different ways to remove $x_0$ change that
class by an image of $a_1$. This defines a canonical surjection
\[
                       \ker\theta_2\longrightarrow K_1. \tag{B82.17}
\]
Surjectivity follows by starting with a representative of a class
in $\ker b_1$ and lifting $h_1x_1$ to $dx_2$.
If the class in $K_1$ is zero, subtracting further total boundaries
makes $x_1=0$ as well. What remains is a vertical cycle
$x_2\in\mathcal T_{2,2}$. Such a pure column-two cycle is a total
boundary exactly when, modulo a vertical boundary, it is of the form
$h_1w$, where
\[
 dz=0,\quad z\in\mathcal T_{0,1},\qquad
 w\in\mathcal T_{1,2},\quad dw=h_0z.                     \tag{B82.18}
\]
Indeed the column-zero and column-one components of a boundary
of total degree one must vanish, giving precisely these equations;
the column-two component is $h_1w$ plus a vertical boundary.
Changing $w$ by a cycle contributes $\operatorname{im}b_2$,
and (B82.18) is exactly the defining lift for $\delta_1$.
Consequently (B82.17) has kernel $K_2$, with a canonical exact
sequence
\[
             0\longrightarrow K_2\longrightarrow\ker\theta_2
                     \longrightarrow K_1\longrightarrow0.   \tag{B82.19}
\]
This proves that uniqueness is equivalent to $K_1=K_2=0$.

All coefficient vectors of integer-valued polynomials belong to
$C$ by B.064. The two explicit arrays in (B82.7) are exactly the
coefficient polynomials of $F(X+Y)$ and $F(XY)$.
Equations (B82.15)--(B82.16) say that they lift exactly when their
$\delta_0$ classes vanish. If (B82.19) is zero, these lifts are
unique and recover the requested tensors. This proves both
directions of (B82.7) for every original inclusion $R\subset S$.

Every complex used to define the criterion is a free complex built
from $\mathcal F$, $E$ and $\varepsilon$, which themselves come
directly from (B82.1). Relations in the free resolutions and the
lifts in (B82.12), (B82.15), and (B82.18) have finite support.
Only $U_1,V_1,W_1$, $V_2,W_2$ and boundaries in vertical degree
three are needed. Thus the test is an explicit arithmetic syzygy
criterion, though infinite rings may give infinitely many such
tests. When $C$ is flat, all positive Tor groups involving
$\mathcal F$ vanish, so $K_1=0$, $K_2=W_2$, $T_1=W_1$ and the
criterion reduces to B.080, with its connecting map fixed by the
following sign comparison. Under the usual double-resolution
augmentation and the connecting convention of (B80.9),
$\delta_0=-\partial$. One may check the sign using a flat resolution
of $E$ whose degree-zero module is $\mathcal F$ and whose positive
degrees resolve $A$: the component $w_2$ of $w$ in
$Q_1\otimes\mathcal F$ satisfies $dw_2=-H$. The usual connecting
map $\tau$ therefore takes its image in $Q_1\otimes E$ to the
negative of the class used in (B80.9). Thus the vanishing tests
are identical, and negating the last cokernel identification gives
exactly (B80.2). The correction terms in (B82.19) are essential
in the nonflat scope; they are not discarded.

Here are two checks which also show how the invariants operate
without assuming flatness. For the forest inclusion of B.067,
write $D=S/R$ and use its monic basis
$b_i=X^i$ for $i<4$, $b_i=(X^4+X^2)X^{i-4}$ for $i\ge4$.
Then the coefficient and value modules in (B82.1) are
\[
 \mathcal F=\bigoplus_{i\ge0}S b_i,\qquad
 E=\bigoplus_{i<4}D b_i,
                                                        \tag{B82.20}
\]
and $\varepsilon$ is the quotient $S\to D$ on the four low
summands and zero on the other summands. This follows directly
from the Hermite value calculation in B.067. Put
$M_q=\operatorname{Tor}_q^R(S,S)$ for $q\ge1$.
The exact sequence $0\to R\to S\to D\to0$ and
$S\otimes_RS=S$ give natural isomorphisms
\[
 M_q\simeq\operatorname{Tor}_q^R(S,D)
     \simeq\operatorname{Tor}_q^R(D,S)\quad(q\ge1).
                                                        \tag{B82.21}
\]
Indeed the surrounding positive Tor groups with the free module
$R$ vanish, and in degree one the remaining map
$S\otimes_RR\to S\otimes_RS$ is an isomorphism. The same
exact sequence gives an injection
\[
 0\longrightarrow M_1\longrightarrow
 \operatorname{Tor}_1^R(D,D)\longrightarrow D\longrightarrow0
                                                        \tag{B82.22}
\]
and isomorphisms $M_q\simeq\operatorname{Tor}_q^R(D,D)$
for $q\ge2$. Tor commutes with these direct sums, so the rows
in (B82.3) decompose by pairs of basis indices. For two low
indices, $a_1$ is the diagonal of $M_1$, while $b_1$ is the
difference followed by the injection in (B82.22); hence the
row homology $K_1$ is zero. For exactly one low index, $a_1$
is an isomorphism and $W_1=0$, again giving zero row homology.
For two high indices both $V_1,W_1$ are zero. Also $b_2$ is
surjective on every pair with nonzero $W_2$, so $T_2=K_2=0$.
Consequently, in this explicitly nonflat case,
\[
                  K_1=K_2=0,\qquad T_1\simeq D^{16}.
                                                        \tag{B82.23}
\]
The connecting map records, up to its fixed minus sign, the
sixteen low-low coefficients modulo $R$. This follows either
by the lift (B82.12) and (B82.22), or by taking the coefficient
cokernel of the displayed split summands. The identities
\[
\begin{aligned}
 Q(X+Y)&=Q(X)+Q(Y),\\
 Q(XY)&=Q(X)Q(Y)+Q(X)Y^2+X^2Q(Y),\qquad Q=X^4+X^2,
\end{aligned}
\]
show that the two arrays for every $QG\in Q S[X]$ have zero
low-low remainder; those for $R[X]_{<4}$ have that remainder
over $R$. Thus both obstruction classes in (B82.7) vanish.
This verifies the positive case directly in the new invariants,
without falsely imposing $M_q=0$.

For the flat cusp of B.079, let $R=k+tV$, $V=l[[t]]$, $K=l((t))$,
$k=\mathbb F_q$, $l$ infinite, and $P=X^q-X$. Its value module
is specified explicitly by the monic basis:
\[
 E=\bigoplus_{i<q}(K/R)X^i\ \oplus\!
   \bigoplus_{\substack{i<q\\j\ge1}}
        (K/t^{1-j}V)X^iP^j.                              \tag{B82.24}
\]
In particular $K/V$ is a direct summand. Choose
$\beta\in l\setminus k$. The relation
$\beta\otimes1-1\otimes\beta$ is nonzero in
$V\otimes_RV$, because reduction modulo $tV$ sends it to
the nonzero element of $l\otimes_k l$ detected by a
$k$-linear functional taking $1$ to zero and $\beta$ to one.
It maps to zero in $K\otimes_RV=K$. The two long exact
sequences for $0\to V\to K\to K/V\to0$, with $K$ flat,
therefore give a unique nonzero class
\[
       \eta\in\operatorname{Tor}_2^R(K/V,K/V)
\]
whose double connecting image is that relation. The split
inclusion of $K/V$ in (B82.24) makes its image nonzero in
$W_2=\operatorname{Tor}_2^R(E,E)$. Since $C=K$ is flat,
this is a nonzero class in $K_2$. The criterion rejects the
cusp at the uniqueness test. Its nonzero kernel class corresponds,
with the connecting sign fixed as above, to
$\beta P\otimes P-P\otimes\beta P$ from B.079.
These calculations check both a nonflat positive inclusion and
a flat negative inclusion within the universal criterion.
\end{proof}
